% =====================================================================
%  Capítulo: Resultados de flexibilidad
%  Plantilla TFG/TFM EPS - Universidad de Alicante
%  Registrar en main.tex con: % =====================================================================
%  Capítulo: Resultados de flexibilidad
%  Plantilla TFG/TFM EPS - Universidad de Alicante
%  Registrar en main.tex con: % =====================================================================
%  Capítulo: Resultados de flexibilidad
%  Plantilla TFG/TFM EPS - Universidad de Alicante
%  Registrar en main.tex con: % =====================================================================
%  Capítulo: Resultados de flexibilidad
%  Plantilla TFG/TFM EPS - Universidad de Alicante
%  Registrar en main.tex con: \input{contenido/capitulos/flexibilidad}
%  (dentro de \mainmatter, tras el capítulo de resultados de rendimiento)
% =====================================================================

\chapter{Resultados de flexibilidad}\label{chap:flexibilidad}

Este capítulo presenta el tercer eje de la comparación: la flexibilidad de
cada framework. A diferencia de la comparativa de rendimiento, que es
cuantitativa, y de la valoración de usabilidad, que recoge la percepción de uso,
la flexibilidad describe una propiedad arquitectónica: hasta qué punto la
herramienta permite adaptarse a escenarios distintos de los previstos por defecto.
Conviene recordar que este eje no debe confundirse con la participación parcial
de clientes medida en el experimento e3, que es una variación cuantitativa del
escenario; aquí el término flexibilidad se reserva para la capacidad de
modificar, sustituir o ampliar el comportamiento del framework.

La valoración se apoya en los diez criterios F1--F10 definidos en la
Sección~\ref{sec:flex-criterios} y se puntúa con la escala de tres niveles
(0, 1, 2) de la Tabla~\ref{tab:escala-flexibilidad}. Para cada framework se
distinguen dos planos de evidencia. El primero es la evidencia experimental:
cambios efectivamente realizados en el código, la configuración y los scripts, y
ejecuciones reales sobre el escenario común. El segundo es la evidencia
documental: funcionalidades avanzadas descritas en la documentación oficial pero
no implementadas en los experimentos del trabajo. Esta separación es deliberada y
evita presentar como probado aquello que solo se ha revisado en la documentación,
manteniendo así la comparabilidad del diseño experimental.

\section{Flexibilidad de Flower}\label{sec:flex-flower}

Flower es un framework agnóstico respecto a la biblioteca de aprendizaje
automático y con una separación clara entre la lógica del cliente y la del
servidor~\parencite{beutel2020flower}, rasgos que condicionan directamente su
flexibilidad. La evaluación experimental se centró en cinco capacidades: modificar
la configuración federada, sustituir la estrategia de agregación, integrar un
modelo y un conjunto de datos propios en PyTorch, registrar métricas personalizadas
y alternar el modo de ejecución. Todas las ejecuciones se realizaron con
\texttt{flwr} 1.29.0 sobre el backend de simulación Ray~\parencite{moritz2018ray},
con Python 3.12 (entorno \texttt{tfg-flower}), una GPU NVIDIA GeForce RTX~3050
Laptop de 6\,GB y el modelo ResNet-18 adaptado sobre CIFAR-10 descrito en el
capítulo de metodología.

\subsection{Adaptaciones para habilitar la evaluación}\label{subsec:flex-flower-adapt}

La primera evidencia de flexibilidad es la propia parametrización del escenario.
Para que Flower recibiera la configuración de federación en cada ejecución, se
ajustó el script de lanzamiento de modo que toda la configuración viajara en una
única cadena pasada a \texttt{flwr run}, como muestra el Código~\ref{cod:flower-fedconfig}.
Con ello fue posible variar el número de supernodos o clientes y la fracción de
GPU asignada a cada uno en función del modo de ejecución, sin tocar el código de
la aplicación. El parámetro \texttt{init-args-num-cpus=2} no es cosmético: se
incorporó para resolver un agotamiento de memoria RAM del anfitrión, una
limitación que se detalla en la Sección~\ref{subsec:flex-flower-lim}.

\begin{bashcode}[title={run\_flower\_experiments.sh}, label={cod:flower-fedconfig}]
--federation-config="num-supernodes=$clients \
  init-args-num-cpus=2 \
  client-resources-num-cpus=1 \
  client-resources-num-gpus=$gpu_share"
\end{bashcode}

Una segunda adaptación afectó a la declaración de la estrategia. Flower exige que
las claves que se sobrescriben mediante \texttt{-{}-run-config} estén previamente
declaradas en la sección \texttt{[tool.flwr.app.config]} del fichero
\texttt{pyproject.toml}; de lo contrario, la ejecución aborta con el mensaje
\mintinline{text}{Key 'strategy' is not present in the main dictionary}. Por ese
motivo se añadió la clave \texttt{strategy} a la configuración de la aplicación,
junto con el resto de hiperparámetros del escenario (Código~\ref{cod:flower-pyproject}).
Asimismo, el valor de \texttt{learning-rate} se fijó en 0,02 para que coincidiera
con el que el script aplica de forma efectiva, eliminando una posible fuente de
confusión entre el valor declarado y el realmente utilizado.

\begin{codigo}[title={pyproject.toml}, label={cod:flower-pyproject}]{toml}
[tool.flwr.app.config]
num-server-rounds = 3
fraction-train = 0.5
fraction-evaluate = 0.0
local-epochs = 1
learning-rate = 0.02
batch-size = 32
strategy = "fedavg"
\end{codigo}

\subsection{Selección dinámica de la estrategia de agregación}\label{subsec:flex-flower-strategy}

El núcleo de la evaluación experimental de flexibilidad fue la sustitución de la
estrategia de agregación sin alterar el resto del sistema. Para ello se modificó
el servidor de Flower de manera que la estrategia se seleccionara dinámicamente a
partir de la configuración de ejecución, leyendo la clave \texttt{strategy} de
\mintinline{python}{context.run_config}. El servidor instancia entonces la clase
correspondiente y, en el caso de las estrategias adaptativas, fija los
hiperparámetros del optimizador del servidor antes de iniciar el entrenamiento,
como recoge el Código~\ref{cod:flower-strategy}. Gracias a este diseño, el mismo
código de cliente, modelo y conjunto de datos puede ejecutarse con distintas
estrategias modificando únicamente la lógica del servidor.

\begin{pythoncode}[title={server\_app.py}, label={cod:flower-strategy}]
from flwr.serverapp.strategy import FedAvg, FedAdagrad, FedAdam

strategy_name = str(context.run_config.get("strategy", "fedavg")).lower()

if strategy_name == "fedadam":
    strategy_class = FedAdam
    adaptive_kwargs = dict(
        eta=float(context.run_config.get("server-lr", 0.01)),
        tau=float(context.run_config.get("tau", 0.01)),
    )
elif strategy_name == "fedadagrad":
    strategy_class = FedAdagrad
    adaptive_kwargs = dict(
        eta=float(context.run_config.get("server-lr", 0.01)),
        tau=float(context.run_config.get("tau", 0.01)),
    )
else:
    strategy_class = FedAvg
    adaptive_kwargs = {}

strategy = strategy_class(**common_kwargs, **adaptive_kwargs)
\end{pythoncode}

Las estrategias FedAdam y FedAdagrad pertenecen a la familia de optimización
federada adaptativa propuesta por \textcite{reddi2021adaptive}, que generaliza
FedAvg sustituyendo el promedio del servidor por un optimizador con momento y tasa
de aprendizaje propios. Con los valores por defecto de Flower
(\texttt{eta}~=~0,1 y \texttt{tau}~=~0,001), ambas estrategias divergieron
numéricamente en este escenario: la pérdida del modelo global creció hasta el
orden de $10^{9}$ mientras las pérdidas locales de entrenamiento descendían con
normalidad. El problema se localiza en el paso de agregación del servidor, donde
una tasa de aprendizaje alta combinada con un factor de adaptabilidad pequeño
produce primeros pasos desproporcionados antes de que el segundo momento se
estabilice. Reducir la tasa del servidor a \texttt{eta}~=~0,01 y elevar el factor
a \texttt{tau}~=~0,01 estabilizó el entrenamiento. Lejos de restar valor a la
evaluación, este ajuste evidencia precisamente la flexibilidad buscada: Flower
expone esos hiperparámetros de agregación y permite configurarlos desde el
servidor.

\subsection{Mini-experimento de flexibilidad}\label{subsec:flex-flower-mini}

Para comprobar la sustitución de estrategias de forma controlada se definió una
variante reducida del experimento e3, denominada \texttt{e3\_flex\_mini}, cuyos
parámetros recoge la Tabla~\ref{tab:flex-flower-config}. La finalidad de esta
prueba no era comparar el rendimiento de aprendizaje, sino demostrar que Flower
permite cambiar la estrategia federada y la configuración de participación sin
rehacer la implementación, reduciendo a la vez el coste temporal de las pruebas.
Se ejecutó con una única semilla (42) y una repetición, por lo que las cifras de
aprendizaje son auxiliares y no llevan estimación de variabilidad.

\begin{table}[htbp]
  \centering
  \caption{Configuración del mini-experimento de flexibilidad de Flower.}
  \label{tab:flex-flower-config}
  \begin{tabular}{lc}
    \toprule
    \textbf{Parámetro} & \textbf{Valor} \\
    \midrule
    Clientes totales & 8 \\
    Clientes por ronda & 4 \\
    Fracción de entrenamiento & 0,5 \\
    Rondas federadas & 3 \\
    Épocas locales & 1 \\
    Tamaño de lote & 32 \\
    Tasa de aprendizaje (cliente) & 0,02 \\
    \texttt{eta} servidor (FedAdam/FedAdagrad) & 0,01 \\
    \texttt{tau} (FedAdam/FedAdagrad) & 0,01 \\
    Modo de ejecución & Secuencial \\
    Repeticiones / semilla & 1 / 42 \\
    \bottomrule
  \end{tabular}
\end{table}

Los tres mini-experimentos terminaron correctamente, lo que confirma que el mismo
flujo se ejecutó con FedAvg, FedAdam y FedAdagrad manteniendo constantes el
conjunto de datos, el modelo, el número de clientes, los clientes por ronda, las
rondas, las épocas locales, el tamaño de lote y la tasa de aprendizaje del
cliente. La Tabla~\ref{tab:flex-flower-resultados} resume las métricas obtenidas.

\begin{table}[htbp]
  \centering
  \caption{Resultados del mini-experimento de flexibilidad de Flower con tres
  estrategias de agregación. Las cifras son auxiliares y no deben interpretarse
  como una comparación de rendimiento. La memoria de GPU corresponde al pico
  reservado por PyTorch al proceso del cliente, métrica comparable entre
  estrategias.}
  \label{tab:flex-flower-resultados}
  \begin{tabular}{lcccrcr}
    \toprule
    \textbf{Estrategia} & \textbf{Estado} & \textbf{Acc.} & \textbf{Loss} &
    \textbf{Dur. (s)} & \textbf{GPU proc. (MiB)} & \textbf{Energía (Wh)} \\
    \midrule
    FedAvg     & success & 0,4257 & 1,5265 & 132 & $\approx$289 & 1,272 \\
    FedAdam    & success & 0,1410 & 2,3401 & 124 & $\approx$289 & 1,252 \\
    FedAdagrad & success & 0,1000 & 18,7832 & 144 & $\approx$289 & 1,123 \\
    \bottomrule
  \end{tabular}
\end{table}

Estos resultados no permiten concluir que una estrategia sea superior a otra, por
dos motivos. En primer lugar, se trata de un mini-experimento de solo tres rondas
y una época local, insuficiente para que las estrategias adaptativas estabilicen
su segundo momento y converjan; con los hiperparámetros ajustados ya no divergen,
pero todavía no alcanzan el comportamiento de FedAvg. En segundo lugar, se ejecutó
con una única semilla, sin medida de variabilidad. Conviene además una
advertencia de medición: la memoria de GPU reportada por \texttt{nvidia-smi} a
nivel de sistema mostró un valor anómalo de 1623~MiB en la ejecución de
FedAdagrad frente a 757~MiB en las otras dos, atribuible a memoria ocupada por
otras aplicaciones del equipo durante esa corrida; el consumo real por proceso fue
prácticamente idéntico ($\approx$289~MiB) en las tres estrategias, por lo que para
cualquier comparación debe emplearse la métrica por proceso. La conclusión
experimental válida es que Flower permitió sustituir la estrategia federada entre
FedAvg, FedAdam y FedAdagrad sin modificar el cliente, el modelo ni el conjunto de
datos, demostrando su flexibilidad para alterar la lógica de agregación desde el
servidor y la configuración de ejecución.

\subsection{Integración de modelo, datos y métricas}\label{subsec:flex-flower-integracion}

Más allá de las estrategias, la evaluación confirmó que Flower integra sin
fricción un modelo y un conjunto de datos propios. El modelo es la ResNet-18
adaptada a CIFAR-10, construida en PyTorch sustituyendo la primera convolución por
una de núcleo 3$\times$3 y eliminando el \textit{maxpool} inicial, mientras que el
conjunto de datos se carga mediante \mintinline{python}{FederatedDataset} con un
particionador IID. El bucle de entrenamiento local es código propio del cliente:
emplea descenso de gradiente estocástico con momento y \textit{weight decay},
entropía cruzada como función de pérdida y un número de épocas, tasa y tamaño de
lote configurables. Además, el cliente registra métricas personalizadas, entre
ellas la pérdida de entrenamiento, el número de ejemplos, la duración y el pico de
memoria de GPU asignada por PyTorch, que el script de ejecución vuelca junto con
los registros completos en ficheros de resumen reutilizables. Esta combinación de
modelo propio, bucle de entrenamiento editable y métricas a medida sustenta una
valoración alta en los criterios de personalización, integración y
reproducibilidad.

\subsection{Capacidades avanzadas según la documentación}\label{subsec:flex-flower-doc}

Más allá de lo ejecutado, la documentación oficial de Flower (versión 1.31) describe
un conjunto amplio de técnicas de flexibilidad que se valoran como evidencia
documental. En el plano de la configuración y el control del servidor, se detallan
cuatro vías de personalización de las estrategias: ajustar los parámetros de una
estrategia existente (\texttt{fraction\_train}, \texttt{fraction\_evaluate},
\texttt{min\_available\_nodes}); inyectar funciones de devolución como
\texttt{evaluate\_fn} que el servidor ejecuta en momentos concretos; sobrescribir
métodos internos como \texttt{configure\_train}, \texttt{aggregate\_fit} o
\texttt{aggregate\_evaluate}; e implementar una estrategia nueva heredando de
\texttt{BaseStrategy}~\parencite{beutel2020flower}. La misma documentación muestra cómo
enviar configuraciones por ronda mediante \texttt{ConfigRecord} —por ejemplo, reducir
progresivamente la tasa de aprendizaje leyendo \texttt{server\_round}— y cómo agregar
métricas con criterios propios a través de \texttt{evaluate\_metrics\_aggr\_fn}, en
lugar del promedio ponderado por número de ejemplos. Los clientes, sin estado por
defecto, pueden conservar información entre rondas mediante \texttt{Context.state}, lo
que habilita técnicas como el \textit{checkpointing} en el cliente.

Una pieza distintiva son los \textit{Mods}, funciones que interceptan los mensajes
entre servidor y cliente para transformarlos sin alterar la lógica principal; se
encadenan en un orden definido y permiten, por ejemplo, recortar gradientes
(\texttt{fixedclipping\_mod}, \texttt{adaptiveclipping\_mod}) o limitar el tamaño de
los mensajes (\texttt{message\_size\_mod})~\parencite{beutel2020flower}. Sobre esta
base se construye la privacidad diferencial, documentada en modo preliminar en tres
variantes —central del lado del servidor, central del lado del cliente y local
mediante \texttt{LocalDpMod}—, todas con recorte y ruido gaussiano configurables, a las
que se suma la agregación segura. La documentación también ejemplifica esquemas no
estándar, como FedBN —que excluye los parámetros de normalización por lotes del
intercambio para preservar las estadísticas locales de cada cliente—, el aprendizaje
federado vertical y la integración con XGBoost.

La documentación recoge asimismo soporte para el entrenamiento asíncrono, si bien de
carácter experimental y fuera de las guías principales. Las notas de versión describen
una \textit{Driver API} experimental (versión 1.2) que ejecuta un servidor programable,
asíncrono y multi-inquilino, y la API \textit{Flower Next} (versión 1.8), un conjunto de
primitivas de bajo nivel que permite personalizar el bucle de entrenamiento e incluye,
de forma explícita, soporte para aprendizaje federado asíncrono y esquemas cíclicos,
con un \texttt{ServerApp} capaz de permanecer activo de forma
indefinida~\parencite{beutel2020flower}. Por su naturaleza en desarrollo y de bajo
nivel, esta capacidad se valora como soporte parcial y no como una funcionalidad
directa y estable. Todas estas capacidades se reconocen como existentes, pero no se
contabilizan como demostradas empíricamente en este trabajo.

\subsection{Ejecución, despliegue y operación según la documentación}\label{subsec:flex-flower-deploy}

La flexibilidad de Flower se extiende a su modelo de ejecución y despliegue. En local,
los comandos \texttt{flwr run} y \texttt{flwr list} levantan automáticamente un
\textit{SuperLink} que gestiona las ejecuciones, mientras que las simulaciones se
lanzan sobre el \textit{Simulation Runtime}, cuyo número de \textit{SuperNodes} y
recursos por cliente se fijan de forma persistente o se sobrescriben por ejecución;
la documentación contempla incluso simulaciones sobre un clúster Ray
multi-anfitrión~\parencite{beutel2020flower}. Para despliegues distribuidos, la
federación se compone de un \textit{SuperLink} que coordina y varios
\textit{SuperNodes} que ejecutan los clientes, modelo que la documentación lleva hasta
entornos de producción en la nube —Google Cloud sobre Kubernetes, Microsoft Azure y
Red Hat OpenShift— e incluso a la interconexión de varios clústeres mediante Skupper.

Este recorrido se acompaña de mecanismos de seguridad operativa que también cuentan
para la flexibilidad: cifrado de las comunicaciones mediante TLS con certificados
propios, autenticación de \textit{SuperNodes} basada en firmas y lista blanca de
claves públicas, autenticación de cuentas de usuario mediante OpenID Connect con
autorización fina y un registro de auditoría en formato JSON que traza qué actor
ejecuta cada acción~\parencite{beutel2020flower}. Por último, la interfaz de línea de
comandos puede emitir su salida en formato JSON mediante la opción
\texttt{-{}-format json} para integrarse en flujos automatizados, y las herramientas de
gestión de federaciones permiten inspeccionar nodos, ejecuciones y estado. Ninguna de
estas capacidades se ejercitó en los experimentos, pero refuerzan la valoración de
Flower en los modos de ejecución y despliegue y en la integración con herramientas
externas.

\subsection{Limitaciones observadas}\label{subsec:flex-flower-lim}

El hallazgo más relevante de la evaluación afecta a la memoria RAM del anfitrión y
no a la GPU. El motor de simulación de Flower delega en Ray, que al iniciarse
detecta los procesadores del equipo (veinte en este caso) y prelanza un proceso
\textit{worker} ocioso por cada uno; cada worker carga el intérprete y las
bibliotecas (en torno a 0,28~GB), de modo que el conjunto de workers, sumado al
proceso de simulación y al resto del sistema, agotaba los aproximadamente 14,6~GB
de RAM disponibles. Ray respondía terminando procesos con un error de memoria
agotada, lo que ocurría incluso en modo secuencial, donde a priori solo hay un
cliente entrenando a la vez. La solución fue limitar el número de procesadores que
ve la \textit{Simulation Runtime} mediante \texttt{init-args-num-cpus=2}, lo que
reduce drásticamente la huella de RAM. Se trata de un resultado comparativo de
interés: el backend basado en Ray presenta una huella de memoria de anfitrión
notablemente mayor que el simulador de NVIDIA FLARE o el backend de proceso único
de FedML, que se ejecutan en un único proceso, por lo que en equipos con poca RAM
la simulación de Flower obliga a acotar los recursos manualmente.

La ejecución concurrente, probada también mediante Ray con una fracción de GPU por
cliente de 0,5, quedó condicionada por esta misma limitación de memoria, en línea
con lo descrito en el capítulo de resultados de rendimiento. Por su parte, la
divergencia de las estrategias adaptativas con los valores por defecto, ya
comentada, es una limitación de uso y no del framework: las estrategias de la
familia FedOpt son sensibles a la tasa de aprendizaje del servidor y requieren un
ajuste deliberado, especialmente con pocas rondas~\parencite{reddi2021adaptive}.
Ninguna de estas limitaciones impidió completar la evaluación, pero todas matizan
la valoración final.

\subsection{Valoración de la flexibilidad de Flower}\label{subsec:flex-flower-score}

La Tabla~\ref{tab:flex-flower-score} resume la valoración de Flower sobre los diez
criterios F1--F10, distinguiendo el tipo de evidencia que sustenta cada
puntuación. Los criterios de personalización del modelo y los datos (F1), del
entrenamiento local (F2), de la estrategia de agregación (F3) y de los algoritmos
federados disponibles (F4) se valoran con la máxima puntuación a partir de
evidencia experimental directa, al igual que los modos de ejecución y despliegue
(F6), la integración con herramientas externas (F7) y la reproducibilidad (F10).
El soporte de aprendizaje asíncrono (F5), la seguridad y privacidad (F8) y la
modificación de la arquitectura federada (F9) reciben puntuación parcial por estar
respaldados solo por documentación; en el caso de F5, mediante interfaces de bajo
nivel todavía en desarrollo (la \textit{Driver API} y \textit{Flower Next}), por lo
que se reconoce el soporte sin elevarlo a la máxima puntuación.

\begin{table}[htbp]
  \centering
  \caption{Valoración de la flexibilidad de Flower según los criterios F1--F10.
  El tipo de evidencia distingue las capacidades ejecutadas en los experimentos
  (experimental) de las descritas en la documentación oficial (documental).}
  \label{tab:flex-flower-score}
  \begin{tabular}{clcl}
    \toprule
    \textbf{Criterio} & \textbf{Descripción} & \textbf{Punt.} & \textbf{Evidencia} \\
    \midrule
    F1  & Personalización del modelo y el dataset      & 2/2 & Experimental \\
    F2  & Personalización del entrenamiento local      & 2/2 & Experimental \\
    F3  & Estrategias de agregación personalizadas     & 2/2 & Experimental \\
    F4  & Algoritmos federados disponibles             & 2/2 & Experimental \\
    F5  & Aprendizaje federado asíncrono               & 1/2 & Documental \\
    F6  & Modos de ejecución y despliegue              & 2/2 & Exp.\ + documental \\
    F7  & Integración con herramientas externas        & 2/2 & Exp.\ + documental \\
    F8  & Seguridad, privacidad y extensiones          & 1/2 & Documental \\
    F9  & Modificación de la arquitectura federada     & 1/2 & Documental \\
    F10 & Reproducibilidad y configuración             & 2/2 & Experimental \\
    \midrule
    \multicolumn{2}{r}{\textbf{Total}} & \multicolumn{2}{l}{\textbf{17/20}} \\
    \bottomrule
  \end{tabular}
\end{table}

En conjunto, Flower obtiene una valoración alta en flexibilidad (17/20). La
evidencia experimental muestra una herramienta cómoda de adaptar:
permitió integrar un modelo y un conjunto de datos propios en PyTorch, editar el
bucle de entrenamiento local, sustituir la estrategia de agregación entre FedAvg,
FedAdam y FedAdagrad desde el servidor —incluyendo el ajuste de sus
hiperparámetros— y automatizar la recogida de métricas. La valoración no alcanza
el máximo porque varias capacidades avanzadas —el aprendizaje asíncrono, la
privacidad diferencial, la agregación segura o los esquemas no horizontales— solo se
verificaron documentalmente, y porque el soporte asíncrono se apoya en interfaces
todavía experimentales; a ello se añade que la viabilidad práctica del modo
concurrente quedó condicionada por la elevada huella de memoria del backend Ray en
el equipo empleado. Estas conclusiones se contrastan con las de NVIDIA FLARE y FedML
en la comparativa final del capítulo.

% =====================================================================
%  SECCIÓN AÑADIDA: Flexibilidad de NVIDIA FLARE
% =====================================================================

\section{Flexibilidad de NVIDIA FLARE}\label{sec:flex-nvflare}

NVIDIA FLARE se presenta como un SDK de aprendizaje federado agnóstico al dominio
y compatible con flujos de trabajo de aprendizaje automático ya existentes, con
una orientación que abarca desde la simulación local hasta el despliegue en
producción y en el borde~\parencite{roth2022nvflare,nvflareWelcome}. La evaluación
experimental de su flexibilidad se apoya en un paquete de cinco ejecuciones de la
variante reducida del experimento e3, la misma \texttt{e3\_flex\_mini} empleada con
Flower (Subsección~\ref{subsec:flex-flower-mini}), con ocho clientes totales,
cuatro por ronda, tres rondas, una época local, tamaño de lote 32, tasa de
aprendizaje 0,02 y semilla 42, en modo secuencial con un único hilo del simulador.
Conviene precisar el alcance: el paquete inspeccionado contiene una selección de
suites concretas y no la totalidad del histórico de ejecuciones, por lo que las
conclusiones experimentales se refieren únicamente a los archivos incluidos en él.

\subsection{Diseño modular: jobs, componentes y workflows}\label{subsec:flex-nvflare-modular}

La flexibilidad de NVIDIA FLARE descansa en su modelo de configuración modular. La
unidad de ejecución es el \textit{job}, que describe por completo un experimento
mediante ficheros de configuración de servidor y de cliente
(\texttt{config\_fed\_server.conf} y \texttt{config\_fed\_client.conf}), en los que
cada elemento del sistema se declara como un componente con su ruta de clase y sus
argumentos de construcción~\parencite{nvflareComponents}. Al desplegar el job, el
servidor y los clientes analizan esos ficheros e instancian dinámicamente los
componentes declarados, registrándolos en el bucle de eventos del framework; cada
componente se identifica con un \texttt{component\_id} que permite referenciarlo
desde otros componentes o workflows~\parencite{nvflareComponents}. Este diseño es
el que permite sustituir un agregador, un persistor, un \textit{executor}, un
generador de objetos compartibles (\textit{Shareable}) o un filtro modificando la
configuración, sin reescribir la aplicación completa.

Sobre esa base de bajo nivel, la documentación ofrece además una capa de
\textit{recipes} que empaqueta algoritmos y workflows habituales —FedAvg, FedProx,
FedOpt, SCAFFOLD, aprendizaje cíclico, XGBoost, Sklearn, estadísticas federadas,
evaluación cruzada, PSI, integración con Flower y swarm learning, entre
otros~\parencite{nvflareRecipes}. Esta dualidad, configuración granular por
componentes y configuración de alto nivel por recetas, sustenta una valoración
elevada en personalización y reproducibilidad, y constituye el punto de partida de
las adaptaciones realizadas para probar estrategias alternativas.

\subsection{Adaptaciones para soportar las estrategias}\label{subsec:flex-nvflare-adapt}

El objetivo de la prueba fue comprobar si un mismo job podía adaptarse para
ejecutar FedAvg, FedProx y variantes FedOpt del servidor (FedAdagrad y FedAdam),
manteniendo constante la estructura general del experimento. Para ello se
introdujeron tres cambios de naturaleza distinta. En primer lugar, la selección de
estrategia se trasladó al cliente, que admite los valores \texttt{fedavg},
\texttt{fedprox}, \texttt{fedadam} y \texttt{fedadagrad} mediante un argumento de
línea de órdenes que, por defecto, toma su valor de la variable de entorno
\texttt{TFG\_FL\_STRATEGY}. En segundo lugar, FedProx se implementó en el
entrenamiento local: cuando la estrategia es \texttt{fedprox} y el coeficiente
proximal es positivo, el entrenamiento añade un término que penaliza la desviación
de los pesos locales respecto al modelo global recibido al inicio de la tarea,
manteniendo el descenso de gradiente estocástico como optimizador
local~\parencite{nvflareAlgorithms}.

El tercer cambio afecta al servidor y es el más relevante para la flexibilidad de
agregación. Para las variantes FedOpt, el script de ejecución parchea el fichero de
configuración del servidor sustituyendo el generador de compartibles por defecto,
\mintinline{text}{FullModelShareableGenerator}, por
\mintinline{text}{PTFedOptModelShareableGenerator}, y añade un componente de modelo
fuente requerido por este generador, como ilustra el
Código~\ref{cod:nvflare-fedopt}. Este generador actualiza el modelo global mediante
un optimizador de PyTorch —y, opcionalmente, un planificador de tasa de
aprendizaje— cuyos argumentos incluyen \texttt{optimizer\_args},
\texttt{lr\_scheduler\_args}, \texttt{source\_model} y
\texttt{device}~\parencite{nvflareFedopt}. En consecuencia, FedAdam y FedAdagrad no
son recetas independientes, sino configuraciones de FedOpt que se obtienen
seleccionando \mintinline{text}{torch.optim.Adam} o
\mintinline{text}{torch.optim.Adagrad} como optimizador del servidor.

\begin{codigo}[title={Fragmento representativo de config\_fed\_server.conf (FedOpt)}, label={cod:nvflare-fedopt}]{json}
"shareable_generator": {
  "path": "nvflare.app_opt.pt.fedopt.PTFedOptModelShareableGenerator",
  "args": {
    "source_model": "model",
    "optimizer_args": { "path": "torch.optim.Adagrad" }
  }
},
"components": [
  { "id": "model", "path": "<ruta de la ResNet-18 adaptada>" }
]
\end{codigo}

Una precaución metodológica relevante surgió con el dispositivo de ejecución. En
todos los ficheros de configuración del cliente aparece \texttt{-{}-device cuda},
pero en la ejecución de FedAdam en CPU el propio cliente seleccionó CPU porque
\mintinline{python}{torch.cuda.is_available()} devolvió falso al lanzarse con
\texttt{CUDA\_VISIBLE\_DEVICES} vacío. Por ello, el dispositivo realmente empleado
debe leerse del registro de ejecución y no del fichero de configuración, criterio
que se ha seguido al elaborar la Tabla~\ref{tab:flex-nvflare-resultados}.

\subsection{Mini-experimento de flexibilidad}\label{subsec:flex-nvflare-mini}

La Tabla~\ref{tab:flex-nvflare-resultados} resume las cinco ejecuciones. FedAvg
sirvió de referencia y se ejecutó correctamente con el workflow \textit{Scatter and
Gather}, en el que el servidor distribuye los pesos iniciales, los clientes
entrenan localmente y devuelven sus pesos, y el sistema los agrega por promedio
ponderado~\parencite{nvflareAlgorithms}. A partir de esa base, FedProx finalizó las
tres rondas y los registros confirmaron \texttt{strategy=fedprox} con un coeficiente
proximal \texttt{fedprox\_mu=0,001} activo durante el entrenamiento. FedAdagrad
también completó las tres rondas como FedOpt real —y no como mera etiqueta— porque
la configuración del servidor incorporaba el generador
\mintinline{text}{PTFedOptModelShareableGenerator} y el registro mostraba las
actualizaciones del servidor con \mintinline{text}{torch.optim.Adagrad}.

\begin{table}[htbp]
  \centering
  \caption{Resultados del mini-experimento de flexibilidad de NVIDIA FLARE. Las
  cifras de aprendizaje son auxiliares (tres rondas, una época, una semilla) y no
  constituyen una comparación de rendimiento. El dispositivo real procede del
  registro de ejecución, no del fichero de configuración del cliente.}
  \label{tab:flex-nvflare-resultados}
  \begin{tabular}{llccccrr}
    \toprule
    \textbf{Estrategia} & \textbf{Estado} & \textbf{Disp.} & \textbf{FedOpt} &
    \textbf{Rondas} & \textbf{Dur. (s)} & \textbf{Acc.} & \textbf{Loss} \\
    \midrule
    FedAvg      & success & \texttt{cuda:0} & No  & 3 & 205  & 0,4178 & 1,5601 \\
    FedProx     & success & \texttt{cuda:0} & No  & 3 & 225  & 0,4045 & 1,5876 \\
    FedAdagrad  & success & \texttt{cuda:0} & Sí  & 3 & 199  & 0,1402 & 2,3785 \\
    FedAdam     & failed  & \texttt{cuda:0} & Sí  & 0 & 77   & ---    & ---    \\
    FedAdam     & success & \texttt{cpu}    & Sí  & 3 & 1280 & 0,2586 & 3,4439 \\
    \bottomrule
  \end{tabular}
\end{table}

El caso de FedAdam exige una lectura cuidadosa. La estrategia llegó a configurarse
realmente mediante FedOpt con \mintinline{text}{torch.optim.Adam}, pero su ejecución
en GPU abortó durante el paso de actualización del servidor
(\mintinline{python}{torch.optim.Adam.step()}) con un error de inicialización de
CUDA, registrado como \texttt{FATAL\_SYSTEM\_ERROR}. Para aislar el problema, la
prueba se repitió en CPU, donde FedAdam completó las tres rondas; esta ejecución
valida FedAdam como estrategia configurable, pero su duración (1280~s) es muy
superior precisamente por entrenarse en CPU, de modo que no debe mezclarse con la
comparativa de rendimiento en GPU del capítulo correspondiente. En coherencia con
ello, la fila de la ejecución fallida no reporta precisión ni pérdida, ya que no
completó ninguna ronda.

Estos resultados no permiten ordenar las estrategias por calidad, por los mismos
motivos señalados para Flower: tres rondas, una época local y una sola semilla son
insuficientes para que las variantes adaptativas converjan. Su valor es
demostrar que NVIDIA FLARE permite adaptar tanto el entrenamiento local como la
lógica de actualización del servidor mediante configuración y componentes. Importa
además una cautela de interpretación: la validez de FedAdagrad y FedAdam no se
sustenta en el nombre de la ejecución, sino en la presencia del generador FedOpt y
del optimizador del servidor en los registros. Del mismo modo, el valor
\texttt{fedprox\_mu=0,001} que aparece en la configuración no implica que FedProx se
aplique a todas las estrategias: en FedAdagrad y FedAdam el entrenamiento imprime
\texttt{fedprox\_mu=0,0}, porque el término proximal solo se activa cuando la
estrategia es FedProx.

De esta experiencia se desprende una observación comparativa de interés. Frente a
Flower, donde la estrategia se sustituía cambiando una única clase en el servidor,
NVIDIA FLARE exigió modificar componentes del servidor, añadir un modelo fuente para
FedOpt y comprobar en los registros que la estrategia configurada era realmente la
que se ejecutaba. La flexibilidad de NVIDIA FLARE es, por tanto, amplia, pero su
ejercicio práctico conlleva una complejidad de configuración mayor.

\subsection{Capacidades avanzadas según la documentación}\label{subsec:flex-nvflare-doc}

Como en el caso de Flower, varias capacidades relevantes no se ejecutaron en los
experimentos y se valoran únicamente como evidencia documental. La documentación
oficial las articula en torno a las \textit{Job Recipes}, una API de alto nivel que
encapsula un trabajo federado exponiendo solo los parámetros clave —\texttt{min\_clients},
\texttt{num\_rounds}, modelo y \textit{script} de entrenamiento— y ocultando la
complejidad de controladores y \textit{workflows}~\parencite{nvflareRecipes}. Estas
recetas se ejecutan sobre tres entornos intercambiables: \texttt{SimEnv} simula los
clientes en procesos locales, \texttt{PocEnv} los lanza como procesos separados en la
misma máquina y \texttt{ProdEnv} despliega el trabajo en una instalación de
producción; un mismo trabajo puede así trasladarse de la simulación al despliegue sin
cambiar el código~\parencite{nvflareRecipes}. El catálogo de recetas cubre FedAvg,
FedProx, FedOpt —donde el optimizador del servidor se fija con un argumento que indica
su ruta de clase e hiperparámetros, por ejemplo \texttt{torch.optim.SGD} con tasa y
momento—, SCAFFOLD y el aprendizaje cíclico, además de XGBoost federado en sus
variantes horizontal, \textit{bagging} y vertical~\parencite{nvflareAlgorithms,nvflareFedopt}.

En cuanto al aprendizaje asíncrono, la receta \texttt{EdgeFedBuffRecipe}, orientada a
dispositivos de borde, admite entrenamiento síncrono y asíncrono: en modo síncrono el
servidor espera todas las actualizaciones antes de generar un nuevo modelo, mientras
que en modo asíncrono lo genera en cuanto recibe la primera, lo que reduce el tiempo
total cuando los dispositivos presentan latencias heterogéneas~\parencite{nvflareEdge,nvflareWelcome}.
La receta expone parámetros como el número máximo de versiones de modelo activas o el
tamaño de selección de dispositivos por ronda. No obstante, la propia documentación
advierte de que la asincronía general aún no está disponible para todas las recetas y
se concentra en el escenario de borde, por lo que su valoración se mantiene como
soporte documental parcial.

En materia de seguridad y privacidad, la documentación es notablemente extensa.
NVIDIA FLARE implementa la privacidad mediante un mecanismo general de filtros que
procesan los objetos compartibles antes o después de la comunicación, configurables en
la receta como \texttt{task\_data\_filters} o \texttt{task\_result\_filters}. Soporta
privacidad diferencial a nivel de muestra mediante la integración de DP-SGD con Opacus
—controlada por los parámetros \texttt{epsilon} y \texttt{delta}— y filtros de DP local
como \texttt{PercentilePrivacy}, que comparte solo las diferencias de pesos por encima
de un percentil, \texttt{SVTPrivacy}, basado en la técnica del vector disperso con ruido
de Laplace, y \texttt{StatisticsPrivacyFilter} para estadísticas
federadas~\parencite{nvflareDP}. A un nivel superior, la solución de computación
confidencial ejecuta la agregación dentro de un entorno de ejecución confiable
(\textit{Trusted Execution Environment}) sobre hardware como AMD SEV-SNP o Intel TDX,
protege el modelo y el código del cliente en máquinas virtuales confidenciales y
permite combinar privacidad diferencial con cifrado homomórfico o agregación
segura~\parencite{nvflareSecurity,nvflareConfidential}.

Por último, NVIDIA FLARE documenta esquemas descentralizados mediante los
\textit{workflows} controlados por clientes, construidos sobre las clases base
\texttt{ServerSideController} y \texttt{ClientSideController} para escenarios donde el
servidor no es plenamente confiable~\parencite{nvflareSwarm}. Entre ellos destacan el
aprendizaje cíclico, en el que cada cliente entrena con el modelo recibido del anterior
según un orden fijo o aleatorio~\parencite{nvflareCyclic}; el aprendizaje \textit{swarm},
que designa un cliente agregador distinto por ronda y mejora la resiliencia frente a
fallos del servidor~\parencite{nvflareSwarm}; y la evaluación cruzada entre sitios, que
permite a los clientes validar modelos de otros mediante comunicación entre pares sin
exponerlos al servidor. A esto se suman el aprendizaje vertical con XGBoost y las
estadísticas federadas, así como las utilidades para integrar el seguimiento de
experimentos con TensorBoard, MLflow o Weights \& Biases sobre cualquier
receta~\parencite{nvflareTracking,nvflareRecipes}. La documentación reconoce, además,
dos cautelas: la API de recetas se encuentra en estado de \textit{technical preview} y
no cubre todavía todos los algoritmos, y la configuración de la computación confidencial
exige infraestructura especializada.

\subsection{Limitaciones observadas}\label{subsec:flex-nvflare-lim}

La limitación más concreta de la evaluación fue el fallo de FedAdam en GPU por un
error de inicialización de CUDA durante la actualización FedOpt del servidor. Debe
interpretarse como una incidencia técnica del entorno, no como una carencia del
framework, puesto que la misma configuración completó las tres rondas en CPU; aun
así, impide presentar FedAdam como ejecución correcta en GPU. La segunda limitación
es la complejidad práctica ya señalada: a diferencia de Flower, alterar la lógica de
agregación obligó a intervenir en componentes del servidor y a validar
cuidadosamente los registros, lo que incrementa la probabilidad de errores de
configuración. Por último, el modo concurrente quedó condicionado por los recursos
del equipo, en línea con lo expuesto en el capítulo de resultados de rendimiento;
aunque la documentación describe un sistema de ejecución multi-job con gestión de
recursos, esa capacidad no pudo validarse de forma estable en el hardware empleado.

\subsection{Valoración de la flexibilidad de NVIDIA FLARE}\label{subsec:flex-nvflare-score}

La Tabla~\ref{tab:flex-nvflare-score} recoge la valoración sobre los criterios
F1--F10, con el mismo convenio aplicado a Flower: puntuación máxima cuando existe
evidencia experimental directa, puntuación parcial cuando la capacidad solo consta
en la documentación y puntuación nula cuando no se identifica soporte. NVIDIA FLARE
obtiene la máxima puntuación en personalización del modelo y los datos (F1), del
entrenamiento local (F2), de la estrategia de agregación (F3) y de los algoritmos
federados disponibles (F4), todos respaldados por las ejecuciones, así como en los
modos de ejecución y despliegue (F6), la integración con herramientas externas (F7)
y la reproducibilidad (F10). La diferencia respecto a Flower aparece en el
aprendizaje asíncrono (F5): mientras que en Flower no se identificó soporte, NVIDIA
FLARE lo documenta de forma explícita mediante FedBuff, lo que justifica una
puntuación parcial. La seguridad y privacidad (F8) y la modificación de la
arquitectura federada (F9) reciben también puntuación parcial por estar respaldadas
solo documentalmente.

\begin{table}[htbp]
  \centering
  \caption{Valoración de la flexibilidad de NVIDIA FLARE según los criterios
  F1--F10. El tipo de evidencia distingue las capacidades ejecutadas en los
  experimentos de las descritas en la documentación oficial.}
  \label{tab:flex-nvflare-score}
  \begin{tabular}{clcl}
    \toprule
    \textbf{Criterio} & \textbf{Descripción} & \textbf{Punt.} & \textbf{Evidencia} \\
    \midrule
    F1  & Personalización del modelo y el dataset      & 2/2 & Exp.\ + documental \\
    F2  & Personalización del entrenamiento local      & 2/2 & Experimental \\
    F3  & Estrategias de agregación personalizadas     & 2/2 & Experimental \\
    F4  & Algoritmos federados disponibles             & 2/2 & Exp.\ + documental \\
    F5  & Aprendizaje federado asíncrono               & 1/2 & Documental \\
    F6  & Modos de ejecución y despliegue              & 2/2 & Exp.\ + documental \\
    F7  & Integración con herramientas externas        & 2/2 & Exp.\ + documental \\
    F8  & Seguridad, privacidad y extensiones          & 1/2 & Documental \\
    F9  & Modificación de la arquitectura federada     & 1/2 & Documental \\
    F10 & Reproducibilidad y configuración             & 2/2 & Exp.\ + documental \\
    \midrule
    \multicolumn{2}{r}{\textbf{Total}} & \multicolumn{2}{l}{\textbf{17/20}} \\
    \bottomrule
  \end{tabular}
\end{table}

En conjunto, NVIDIA FLARE alcanza una valoración de flexibilidad de 17/20. La
evidencia experimental
demuestra que permite adaptar el entrenamiento local —incluido el término proximal
de FedProx— y, sobre todo, la lógica de agregación del servidor mediante FedOpt, con
FedAvg, FedProx y FedAdagrad ejecutados correctamente y FedAdam validado en CPU. A
ello se suma una de las superficies de capacidades documentadas más amplias del
estudio, especialmente sólida en seguridad y privacidad, que abarca aprendizaje
asíncrono, workflows cíclicos y \textit{swarm}, privacidad
diferencial, cifrado homomórfico y computación confidencial. El contrapunto es una
complejidad de configuración superior a la de Flower y el hecho de que esas
capacidades avanzadas solo se verificaron documentalmente; además, la ejecución de
FedAdam en GPU falló por un error de CUDA y el modo concurrente quedó limitado por el
hardware. Esta lectura se contrasta con la de los otros frameworks en la comparativa
final del capítulo.

% =====================================================================
%  SECCIÓN AÑADIDA: Flexibilidad de FedML
% =====================================================================

\section{Flexibilidad de FedML}\label{sec:flex-fedml}

FedML es una biblioteca de investigación y banco de pruebas para aprendizaje
federado~\parencite{he2020fedml}, documentada actualmente dentro de la plataforma
TensorOpera aunque la librería y los ejemplos siguen empleando el paquete
\texttt{fedml}. Su documentación organiza el framework en varios escenarios
—simulación, cross-silo y cross-device— y ofrece un catálogo amplio de algoritmos
de referencia. La evaluación de su flexibilidad combina la implementación propia
realizada para la comparativa, que extiende las abstracciones de FedML, con una
prueba reducida de sustitución de estrategias. Conviene una precisión de alcance:
por el coste computacional, la prueba reducida de FedML se ejecutó con cuatro
clientes totales, dos por ronda y dos rondas, una configuración algo menor que la
empleada con Flower y NVIDIA FLARE; las métricas de aprendizaje son, por tanto, una
comprobación funcional y no una comparación de rendimiento.

\subsection{Configuración y componentes personalizados}\label{subsec:flex-fedml-config}

La configuración de FedML se apoya en un fichero \texttt{fedml\_config.yaml}
estructurado en secciones (\texttt{common\_args}, \texttt{data\_args},
\texttt{model\_args}, \texttt{train\_args}, \texttt{device\_args},
\texttt{comm\_args} y \texttt{tracking\_args}), donde se declaran los parámetros
federados como \texttt{federated\_optimizer}, \texttt{client\_num\_in\_total},
\texttt{client\_num\_per\_round}, \texttt{comm\_round}, \texttt{epochs},
\texttt{batch\_size}, \texttt{client\_optimizer} y \texttt{learning\_rate}~\parencite{fedmlSimApi}.
Esta organización permitió variar clientes, rondas, épocas e hiperparámetros entre
ejecuciones sin reescribir la aplicación.

Más allá del fichero de configuración, la flexibilidad de FedML en este trabajo se
sustenta en evidencia experimental directa: la implementación de la comparativa
define un entrenador y un agregador propios que extienden las abstracciones
\mintinline{python}{ClientTrainer} y \mintinline{python}{ServerAggregator} del
framework. El entrenador \texttt{FedMLCifar10Trainer} encapsula el bucle de
entrenamiento local sobre la ResNet-18 adaptada —descenso de gradiente estocástico,
entropía cruzada y métricas propias—, mientras que el agregador
\texttt{FedMLCifar10Aggregator} gobierna la evaluación del modelo global; ambos se
conectan al flujo mediante \texttt{FedMLRunner}. La documentación respalda esta vía,
al indicar que el usuario puede definir modelo, datos, optimizadores, funciones de
pérdida y lógica de entrenamiento arbitrarios, así como dataloaders propios que
respeten la estructura esperada por el simulador~\parencite{fedmlTrainer,fedmlDataloader}.
Esta capacidad de sustituir componentes centrales mediante herencia sostiene una
valoración alta en personalización del entrenamiento local, de la agregación y de la
integración.

\subsection{Sustitución de la estrategia federada}\label{subsec:flex-fedml-strategy}

La prueba de flexibilidad consistió en ejecutar cuatro estrategias modificando
únicamente la configuración, manteniendo constantes el conjunto de datos, el
modelo y el resto de hiperparámetros. FedAvg y FedProx se seleccionan asignando
\texttt{federated\_optimizer} a \texttt{"FedAvg"} o \texttt{"FedProx"}, mientras que
FedAdam y FedAdagrad se obtienen mediante \texttt{federated\_optimizer:\ "FedOpt"}
combinado con un optimizador de servidor \texttt{Adam} o \texttt{Adagrad}, tal como
ilustra el Código~\ref{cod:fedml-config}. Durante las primeras pruebas, las
variantes FedOpt fallaban porque FedML esperaba el atributo \texttt{ci} en los
argumentos de configuración; tras añadir \texttt{ci:\ 0} a la configuración común,
ambas pudieron completar la ejecución. La selección del algoritmo desde la
configuración está respaldada por la documentación, que controla el optimizador
federado a través del mismo campo~\parencite{fedmlSimApi}.

\begin{yamlcode}[title={Fragmento de fedml\_config.yaml (FedAdam vía FedOpt)}, label={cod:fedml-config}]
train_args:
  federated_optimizer: "FedOpt"
  server_optimizer: "Adam"
  ci: 0
  client_num_in_total: 4
  client_num_per_round: 2
  comm_round: 2
  epochs: 1
  batch_size: 32
  learning_rate: 0.02
\end{yamlcode}

\subsection{Mini-experimento de flexibilidad}\label{subsec:flex-fedml-mini}

La Tabla~\ref{tab:flex-fedml-resultados} resume las cuatro ejecuciones, todas
finalizadas con estado \texttt{success}. FedAvg y FedProx concluyeron con métricas
coherentes para una prueba de dos rondas, con precisiones del 43,01\,\% y el
42,60\,\% respectivamente, lo que confirma el soporte experimental de una estrategia
alternativa a FedAvg con comportamiento correcto. FedAdam y FedAdagrad también
completaron la ejecución, lo que demuestra que fue posible activar FedOpt con
optimizadores de servidor Adam y Adagrad; sin embargo, sus métricas finales fueron
claramente deficientes, con una precisión próxima al 10\,\% —equivalente al azar en
un problema de diez clases— y pérdidas muy elevadas.

\begin{table}[htbp]
  \centering
  \caption{Resultados del mini-experimento de flexibilidad de FedML. Las cifras de
  aprendizaje son auxiliares (dos rondas, una época, una semilla) y constituyen una
  comprobación funcional, no una comparación de rendimiento entre estrategias.}
  \label{tab:flex-fedml-resultados}
  \begin{tabular}{llrrrrr}
    \toprule
    \textbf{Estrategia} & \textbf{Estado} & \textbf{Dur. (s)} & \textbf{Acc.} &
    \textbf{Loss} & \textbf{GPU (MiB)} & \textbf{Util. (\%)} \\
    \midrule
    FedAvg     & success & 118 & 0,4301 & 1,5248   & 416 & 75,29 \\
    FedProx    & success & 123 & 0,4260 & 1,5137   & 484 & 79,01 \\
    FedAdam    & success & 117 & 0,0998 & 809,7697 & 416 & 76,09 \\
    FedAdagrad & success & 116 & 0,0998 & 191,2948 & 416 & 76,85 \\
    \bottomrule
  \end{tabular}
\end{table}

La lectura correcta de estos datos es que FedML permite cambiar la estrategia
federada mediante configuración, pero que la mera disponibilidad de una estrategia
no garantiza su estabilidad ni un buen rendimiento. Las variantes FedOpt
ejecutaron de principio a fin, de modo que la flexibilidad queda demostrada, pero su
falta de aprendizaje en esta prueba indica que requieren un ajuste específico de
hiperparámetros —como la tasa de aprendizaje del servidor o un mayor número de
rondas— para resultar competitivas. Este patrón es coherente con lo observado en los
otros frameworks: en Flower las estrategias adaptativas divergían con los valores por
defecto y exigieron ajustar \texttt{eta} y \texttt{tau}, y en NVIDIA FLARE FedAdam no
llegó a completarse en GPU. La fragilidad de las variantes FedOpt con poca
configuración es, por tanto, un comportamiento transversal y no una carencia
particular de FedML.

\subsection{Capacidades avanzadas según la documentación}\label{subsec:flex-fedml-doc}

La documentación de FedML respalda una superficie de capacidades más amplia que la
probada experimentalmente, organizada en varias capas. En el plano de uso, ofrece una
interfaz de línea de comandos —integrada en la plataforma TensorOpera— con órdenes para
autenticación y gestión de dispositivos (\texttt{fedml login}, \texttt{fedml device
bind}), lanzamiento y gestión de trabajos y clústeres (\texttt{fedml launch},
\texttt{fedml cluster}, \texttt{fedml run}), y gestión y despliegue de modelos como
servicios de inferencia (\texttt{fedml model})~\parencite{fedmlExamples}. En el plano de
programación, expone una API de Python con una llamada de una línea
(\texttt{fedml.run\_simulation()}) y una API más explícita de varias líneas que
inicializa dispositivo, datos y modelo y ejecuta la federación mediante
\texttt{SimulatorMPI} o las clases \texttt{Client}/\texttt{Server} envueltas por
\texttt{FedMLRunner}~\parencite{fedmlSimApi}.

En cuanto al catálogo algorítmico, el simulador Parrot enumera implementaciones de
referencia como FedOpt, FedNova, FedGKT, aprendizaje federado descentralizado, vertical
y jerárquico, FedNAS y \textit{Split Learning}, seleccionables cambiando el campo
\texttt{federated\_optimizer}~\parencite{fedmlSimOverview,fedmlSimApi}, y el módulo de
seguridad declara compatibilidad con optimizadores como FedAvg, FedSGD, FedOpt, FedProx,
FedGKT y FedNova~\parencite{fedmlSecurity}. A este catálogo se añaden la analítica
federada, que calcula estadísticas distribuidas como promedios o percentiles mediante un
ejecutor específico sin compartir datos crudos, y un proyecto independiente orientado al
ajuste federado de grandes modelos de lenguaje, que reutiliza las mismas API y la
plataforma de operaciones~\parencite{fedmlExamples}. En cuanto al aprendizaje asíncrono,
FedML Octopus documenta un escenario jerárquico heterogéneo en el que cada silo puede
ejecutar entrenamiento distribuido con varias GPU y orquestarse mediante optimización
federada asíncrona o síncrona~\parencite{fedmlCrossSilo}; se trata, no obstante, de un
soporte menos específico que el de NVIDIA FLARE, ya que no se documenta una receta directa
equivalente a FedBuff, aprendizaje cíclico o \textit{swarm}.

En privacidad y seguridad, el perfil de FedML difiere del de NVIDIA FLARE. Su
fortaleza documental está en la simulación de ataques y defensas mediante
\texttt{FedMLAttacker} y \texttt{FedMLDefender}, con ataques de tipo bizantino,
\textit{label flipping} y reemplazo de modelo, y defensas como recorte de norma, Krum,
mediana geométrica o FoolsGold, activables desde el fichero de configuración con
\texttt{enable\_attack} y \texttt{enable\_defense}~\parencite{fedmlSecurity}, además de un
ejemplo de agregación segura basada en computación multiparte~\parencite{fedmlExamples}.
En cambio, en las páginas revisadas no figura una integración de privacidad diferencial o
cifrado homomórfico comparable a la de NVIDIA FLARE, por lo que su valoración en este
criterio se sostiene en el \textit{benchmarking} de robustez más que en la privacidad
criptográfica. Por último, para integración y trazabilidad, FedML documenta el registro
de métricas, artefactos y modelos mediante \texttt{fedml.log()}, la captura automática de
métricas de hardware y GPU, los backends de comunicación MPI, gRPC, PyTorch RPC y MQTT+S3,
y la integración con la plataforma de operaciones de TensorOpera para gestionar agentes,
experimentos y despliegues~\parencite{fedmlTracking,fedmlExamples}.

\subsection{Limitaciones observadas}\label{subsec:flex-fedml-lim}

La limitación más visible es que las variantes FedOpt finalizaron sin aprender bajo
la configuración empleada, lo que confirma que su soporte es real pero condicionado a
un ajuste de hiperparámetros que queda fuera del alcance de esta prueba de
flexibilidad. A ello se añade el detalle de configuración del atributo \texttt{ci},
que tuvo que fijarse manualmente para que FedOpt funcionara, y el hecho de que la
prueba reducida de FedML emplea una configuración menor que la de los otros dos
frameworks, por lo que sus métricas auxiliares no son directamente comparables. En
cuanto a la ejecución, el modo concurrente de FedML mediante el backend MPI fue el
único que llegó a completarse entre los tres frameworks, según se documenta en el
capítulo de resultados de rendimiento, si bien se mostró sensible al consumo de
memoria en las configuraciones más grandes.

\subsection{Valoración de la flexibilidad de FedML}\label{subsec:flex-fedml-score}

La Tabla~\ref{tab:flex-fedml-score} recoge la valoración de FedML sobre los criterios
F1--F10, con el convenio ya aplicado a los otros frameworks. FedML obtiene la máxima
puntuación en personalización del modelo y los datos (F1), del entrenamiento local
(F2) y de la agregación (F3), respaldadas por la implementación de entrenador y
agregador propios; en los algoritmos federados disponibles (F4), por haber ejecutado
cuatro estrategias mediante configuración; en los modos de ejecución (F6), por ser el
único framework cuyo modo concurrente se completó; y en integración (F7) y
reproducibilidad (F10). El aprendizaje asíncrono (F5), la seguridad y privacidad (F8)
y la modificación de la arquitectura federada (F9) reciben puntuación parcial por
estar respaldados únicamente por la documentación.

\begin{table}[htbp]
  \centering
  \caption{Valoración de la flexibilidad de FedML según los criterios F1--F10. El
  tipo de evidencia distingue las capacidades ejecutadas en los experimentos de las
  descritas en la documentación oficial.}
  \label{tab:flex-fedml-score}
  \begin{tabular}{clcl}
    \toprule
    \textbf{Criterio} & \textbf{Descripción} & \textbf{Punt.} & \textbf{Evidencia} \\
    \midrule
    F1  & Personalización del modelo y el dataset      & 2/2 & Exp.\ + documental \\
    F2  & Personalización del entrenamiento local      & 2/2 & Exp.\ + documental \\
    F3  & Estrategias de agregación personalizadas     & 2/2 & Experimental \\
    F4  & Algoritmos federados disponibles             & 2/2 & Exp.\ + documental \\
    F5  & Aprendizaje federado asíncrono               & 1/2 & Documental \\
    F6  & Modos de ejecución y despliegue              & 2/2 & Exp.\ + documental \\
    F7  & Integración con herramientas externas        & 2/2 & Exp.\ + documental \\
    F8  & Seguridad, privacidad y extensiones          & 1/2 & Documental \\
    F9  & Modificación de la arquitectura federada     & 1/2 & Documental \\
    F10 & Reproducibilidad y configuración             & 2/2 & Exp.\ + documental \\
    \midrule
    \multicolumn{2}{r}{\textbf{Total}} & \multicolumn{2}{l}{\textbf{17/20}} \\
    \bottomrule
  \end{tabular}
\end{table}

En conjunto, FedML alcanza una valoración de flexibilidad de 17/20. Su mayor fortaleza
experimental reside en la facilidad para sustituir componentes centrales del sistema
—entrenador y agregador— mediante herencia, y en haber ejecutado el abanico más amplio
de estrategias del estudio a través de la configuración. Su perfil documental es
asimismo notable, con un catálogo extenso de algoritmos y escenarios y un soporte de
seguridad orientado al benchmarking de ataques y defensas. Las cautelas son dos: las
variantes FedOpt completaron la ejecución pero no aprendieron sin un ajuste
adicional, y algunas capacidades avanzadas —asincronía, privacidad criptográfica y
esquemas alternativos— solo se verificaron documentalmente.

% =====================================================================
%  SECCIÓN AÑADIDA: Comparativa de flexibilidad entre frameworks
% =====================================================================

\section{Comparativa de flexibilidad}\label{sec:flex-comparativa}

Una vez evaluados los tres frameworks por separado, esta sección los contrasta sobre
el mismo conjunto de criterios. La Tabla~\ref{tab:flex-comparativa} reúne las
puntuaciones F1--F10 de Flower, NVIDIA FLARE y FedML. Los tres obtienen la misma
valoración, 17/20, lo que confirma que las tres herramientas son, en lo esencial,
comparablemente flexibles para el tipo de adaptaciones que exige este trabajo y que la
diferenciación entre ellas debe buscarse, más que en la cifra global, en el perfil
cualitativo de cada una.

\begin{table}[htbp]
  \centering
  \caption{Comparativa de la valoración de flexibilidad de los tres frameworks según
  los criterios F1--F10 definidos en la metodología.}
  \label{tab:flex-comparativa}
  \begin{tabular}{clccc}
    \toprule
    \textbf{Criterio} & \textbf{Descripción} & \textbf{Flower} & \textbf{NVIDIA FLARE} & \textbf{FedML} \\
    \midrule
    F1  & Modelo y dataset                  & 2/2 & 2/2 & 2/2 \\
    F2  & Entrenamiento local               & 2/2 & 2/2 & 2/2 \\
    F3  & Agregación personalizada          & 2/2 & 2/2 & 2/2 \\
    F4  & Algoritmos disponibles            & 2/2 & 2/2 & 2/2 \\
    F5  & Aprendizaje asíncrono             & 1/2 & 1/2 & 1/2 \\
    F6  & Modos de ejecución y despliegue   & 2/2 & 2/2 & 2/2 \\
    F7  & Integración externa               & 2/2 & 2/2 & 2/2 \\
    F8  & Seguridad y privacidad            & 1/2 & 1/2 & 1/2 \\
    F9  & Arquitectura federada             & 1/2 & 1/2 & 1/2 \\
    F10 & Reproducibilidad y configuración  & 2/2 & 2/2 & 2/2 \\
    \midrule
    \multicolumn{2}{r}{\textbf{Total}} & \textbf{17/20} & \textbf{17/20} & \textbf{17/20} \\
    \bottomrule
  \end{tabular}
\end{table}

La coincidencia en los criterios experimentales (F1--F4, F6, F7 y F10) refleja un
hecho de fondo: bajo un mismo escenario de modelo, datos y algoritmo base, los tres
frameworks permitieron integrar la ResNet-18 sobre CIFAR-10, editar el entrenamiento
local, sustituir la estrategia de agregación, registrar métricas y reproducir los
experimentos mediante ficheros y scripts. La igualdad también alcanza a los tres
criterios documentales: ninguno validó experimentalmente el aprendizaje asíncrono (F5),
la seguridad avanzada (F8) ni esquemas alternativos a la arquitectura clásica (F9), pero
los tres los documentan, por lo que reciben idéntica puntuación parcial. En F5, en
concreto, los tres declaran soporte —Flower mediante la \textit{Driver API} y
\textit{Flower Next}, NVIDIA FLARE mediante FedBuff en el borde y FedML en su escenario
jerárquico cross-silo—, si bien en todos los casos se trata de capacidades
experimentales, limitadas a entornos concretos o de bajo nivel.

Más allá de las cifras, el estudio revela un comportamiento transversal en las
estrategias adaptativas de la familia FedOpt. En los tres frameworks fue posible
configurar FedAdam y FedAdagrad, pero su ejecución resultó frágil con poca
configuración: en Flower divergían con los valores por defecto y necesitaron ajustar
la tasa del servidor y el factor de adaptabilidad; en NVIDIA FLARE, FedAdam no llegó a
completarse en GPU por un error de CUDA; y en FedML, FedAdam y FedAdagrad finalizaron
pero sin aprender. La conclusión es clara y común: la disponibilidad de una estrategia
no equivale a su estabilidad, que depende de un ajuste deliberado de hiperparámetros,
especialmente con pocas rondas.

Por último, cada framework muestra un perfil de flexibilidad propio que la puntuación,
por su granularidad, no captura del todo. Flower es el más directo para sustituir la
estrategia de agregación —basta con cambiar una clase en el servidor— y su documentación
revela una superficie amplia que abarca \textit{Mods}, privacidad diferencial, FedBN y,
sobre todo, un modelo de despliegue muy desarrollado hacia la nube (Google Cloud, Azure,
OpenShift e interconexión de clústeres); su contrapartida es que el soporte asíncrono
sigue siendo experimental y que el backend Ray arrastra una huella de memoria elevada.
NVIDIA FLARE ofrece la mayor profundidad en seguridad y privacidad —filtros de privacidad
diferencial, cifrado homomórfico y computación confidencial sobre entornos de ejecución
confiables— junto con un catálogo de recetas y entornos (\texttt{SimEnv}, \texttt{PocEnv},
\texttt{ProdEnv}) orientado a producción, a costa de una complejidad de configuración
mayor, pues alterar la agregación exige intervenir en componentes del servidor. FedML
destaca por la sustitución de entrenador y agregador mediante herencia, por haber
ejecutado el mayor número de estrategias, por ser el único cuyo modo concurrente se
completó, y por una integración estrecha con su plataforma de operaciones y la analítica
federada, con un soporte de seguridad orientado al \textit{benchmarking} de ataques y
defensas. La elección entre ellos, en términos de flexibilidad, depende por tanto de la
prioridad del proyecto: la simplicidad y el despliegue en la nube de Flower, la
profundidad en seguridad y la orientación a producción de NVIDIA FLARE, o la
extensibilidad por componentes y la experimentación algorítmica de FedML.

%  (dentro de \mainmatter, tras el capítulo de resultados de rendimiento)
% =====================================================================

\chapter{Resultados de flexibilidad}\label{chap:flexibilidad}

Este capítulo presenta el tercer eje de la comparación: la flexibilidad de
cada framework. A diferencia de la comparativa de rendimiento, que es
cuantitativa, y de la valoración de usabilidad, que recoge la percepción de uso,
la flexibilidad describe una propiedad arquitectónica: hasta qué punto la
herramienta permite adaptarse a escenarios distintos de los previstos por defecto.
Conviene recordar que este eje no debe confundirse con la participación parcial
de clientes medida en el experimento e3, que es una variación cuantitativa del
escenario; aquí el término flexibilidad se reserva para la capacidad de
modificar, sustituir o ampliar el comportamiento del framework.

La valoración se apoya en los diez criterios F1--F10 definidos en la
Sección~\ref{sec:flex-criterios} y se puntúa con la escala de tres niveles
(0, 1, 2) de la Tabla~\ref{tab:escala-flexibilidad}. Para cada framework se
distinguen dos planos de evidencia. El primero es la evidencia experimental:
cambios efectivamente realizados en el código, la configuración y los scripts, y
ejecuciones reales sobre el escenario común. El segundo es la evidencia
documental: funcionalidades avanzadas descritas en la documentación oficial pero
no implementadas en los experimentos del trabajo. Esta separación es deliberada y
evita presentar como probado aquello que solo se ha revisado en la documentación,
manteniendo así la comparabilidad del diseño experimental.

\section{Flexibilidad de Flower}\label{sec:flex-flower}

Flower es un framework agnóstico respecto a la biblioteca de aprendizaje
automático y con una separación clara entre la lógica del cliente y la del
servidor~\parencite{beutel2020flower}, rasgos que condicionan directamente su
flexibilidad. La evaluación experimental se centró en cinco capacidades: modificar
la configuración federada, sustituir la estrategia de agregación, integrar un
modelo y un conjunto de datos propios en PyTorch, registrar métricas personalizadas
y alternar el modo de ejecución. Todas las ejecuciones se realizaron con
\texttt{flwr} 1.29.0 sobre el backend de simulación Ray~\parencite{moritz2018ray},
con Python 3.12 (entorno \texttt{tfg-flower}), una GPU NVIDIA GeForce RTX~3050
Laptop de 6\,GB y el modelo ResNet-18 adaptado sobre CIFAR-10 descrito en el
capítulo de metodología.

\subsection{Adaptaciones para habilitar la evaluación}\label{subsec:flex-flower-adapt}

La primera evidencia de flexibilidad es la propia parametrización del escenario.
Para que Flower recibiera la configuración de federación en cada ejecución, se
ajustó el script de lanzamiento de modo que toda la configuración viajara en una
única cadena pasada a \texttt{flwr run}, como muestra el Código~\ref{cod:flower-fedconfig}.
Con ello fue posible variar el número de supernodos o clientes y la fracción de
GPU asignada a cada uno en función del modo de ejecución, sin tocar el código de
la aplicación. El parámetro \texttt{init-args-num-cpus=2} no es cosmético: se
incorporó para resolver un agotamiento de memoria RAM del anfitrión, una
limitación que se detalla en la Sección~\ref{subsec:flex-flower-lim}.

\begin{bashcode}[title={run\_flower\_experiments.sh}, label={cod:flower-fedconfig}]
--federation-config="num-supernodes=$clients \
  init-args-num-cpus=2 \
  client-resources-num-cpus=1 \
  client-resources-num-gpus=$gpu_share"
\end{bashcode}

Una segunda adaptación afectó a la declaración de la estrategia. Flower exige que
las claves que se sobrescriben mediante \texttt{-{}-run-config} estén previamente
declaradas en la sección \texttt{[tool.flwr.app.config]} del fichero
\texttt{pyproject.toml}; de lo contrario, la ejecución aborta con el mensaje
\mintinline{text}{Key 'strategy' is not present in the main dictionary}. Por ese
motivo se añadió la clave \texttt{strategy} a la configuración de la aplicación,
junto con el resto de hiperparámetros del escenario (Código~\ref{cod:flower-pyproject}).
Asimismo, el valor de \texttt{learning-rate} se fijó en 0,02 para que coincidiera
con el que el script aplica de forma efectiva, eliminando una posible fuente de
confusión entre el valor declarado y el realmente utilizado.

\begin{codigo}[title={pyproject.toml}, label={cod:flower-pyproject}]{toml}
[tool.flwr.app.config]
num-server-rounds = 3
fraction-train = 0.5
fraction-evaluate = 0.0
local-epochs = 1
learning-rate = 0.02
batch-size = 32
strategy = "fedavg"
\end{codigo}

\subsection{Selección dinámica de la estrategia de agregación}\label{subsec:flex-flower-strategy}

El núcleo de la evaluación experimental de flexibilidad fue la sustitución de la
estrategia de agregación sin alterar el resto del sistema. Para ello se modificó
el servidor de Flower de manera que la estrategia se seleccionara dinámicamente a
partir de la configuración de ejecución, leyendo la clave \texttt{strategy} de
\mintinline{python}{context.run_config}. El servidor instancia entonces la clase
correspondiente y, en el caso de las estrategias adaptativas, fija los
hiperparámetros del optimizador del servidor antes de iniciar el entrenamiento,
como recoge el Código~\ref{cod:flower-strategy}. Gracias a este diseño, el mismo
código de cliente, modelo y conjunto de datos puede ejecutarse con distintas
estrategias modificando únicamente la lógica del servidor.

\begin{pythoncode}[title={server\_app.py}, label={cod:flower-strategy}]
from flwr.serverapp.strategy import FedAvg, FedAdagrad, FedAdam

strategy_name = str(context.run_config.get("strategy", "fedavg")).lower()

if strategy_name == "fedadam":
    strategy_class = FedAdam
    adaptive_kwargs = dict(
        eta=float(context.run_config.get("server-lr", 0.01)),
        tau=float(context.run_config.get("tau", 0.01)),
    )
elif strategy_name == "fedadagrad":
    strategy_class = FedAdagrad
    adaptive_kwargs = dict(
        eta=float(context.run_config.get("server-lr", 0.01)),
        tau=float(context.run_config.get("tau", 0.01)),
    )
else:
    strategy_class = FedAvg
    adaptive_kwargs = {}

strategy = strategy_class(**common_kwargs, **adaptive_kwargs)
\end{pythoncode}

Las estrategias FedAdam y FedAdagrad pertenecen a la familia de optimización
federada adaptativa propuesta por \textcite{reddi2021adaptive}, que generaliza
FedAvg sustituyendo el promedio del servidor por un optimizador con momento y tasa
de aprendizaje propios. Con los valores por defecto de Flower
(\texttt{eta}~=~0,1 y \texttt{tau}~=~0,001), ambas estrategias divergieron
numéricamente en este escenario: la pérdida del modelo global creció hasta el
orden de $10^{9}$ mientras las pérdidas locales de entrenamiento descendían con
normalidad. El problema se localiza en el paso de agregación del servidor, donde
una tasa de aprendizaje alta combinada con un factor de adaptabilidad pequeño
produce primeros pasos desproporcionados antes de que el segundo momento se
estabilice. Reducir la tasa del servidor a \texttt{eta}~=~0,01 y elevar el factor
a \texttt{tau}~=~0,01 estabilizó el entrenamiento. Lejos de restar valor a la
evaluación, este ajuste evidencia precisamente la flexibilidad buscada: Flower
expone esos hiperparámetros de agregación y permite configurarlos desde el
servidor.

\subsection{Mini-experimento de flexibilidad}\label{subsec:flex-flower-mini}

Para comprobar la sustitución de estrategias de forma controlada se definió una
variante reducida del experimento e3, denominada \texttt{e3\_flex\_mini}, cuyos
parámetros recoge la Tabla~\ref{tab:flex-flower-config}. La finalidad de esta
prueba no era comparar el rendimiento de aprendizaje, sino demostrar que Flower
permite cambiar la estrategia federada y la configuración de participación sin
rehacer la implementación, reduciendo a la vez el coste temporal de las pruebas.
Se ejecutó con una única semilla (42) y una repetición, por lo que las cifras de
aprendizaje son auxiliares y no llevan estimación de variabilidad.

\begin{table}[htbp]
  \centering
  \caption{Configuración del mini-experimento de flexibilidad de Flower.}
  \label{tab:flex-flower-config}
  \begin{tabular}{lc}
    \toprule
    \textbf{Parámetro} & \textbf{Valor} \\
    \midrule
    Clientes totales & 8 \\
    Clientes por ronda & 4 \\
    Fracción de entrenamiento & 0,5 \\
    Rondas federadas & 3 \\
    Épocas locales & 1 \\
    Tamaño de lote & 32 \\
    Tasa de aprendizaje (cliente) & 0,02 \\
    \texttt{eta} servidor (FedAdam/FedAdagrad) & 0,01 \\
    \texttt{tau} (FedAdam/FedAdagrad) & 0,01 \\
    Modo de ejecución & Secuencial \\
    Repeticiones / semilla & 1 / 42 \\
    \bottomrule
  \end{tabular}
\end{table}

Los tres mini-experimentos terminaron correctamente, lo que confirma que el mismo
flujo se ejecutó con FedAvg, FedAdam y FedAdagrad manteniendo constantes el
conjunto de datos, el modelo, el número de clientes, los clientes por ronda, las
rondas, las épocas locales, el tamaño de lote y la tasa de aprendizaje del
cliente. La Tabla~\ref{tab:flex-flower-resultados} resume las métricas obtenidas.

\begin{table}[htbp]
  \centering
  \caption{Resultados del mini-experimento de flexibilidad de Flower con tres
  estrategias de agregación. Las cifras son auxiliares y no deben interpretarse
  como una comparación de rendimiento. La memoria de GPU corresponde al pico
  reservado por PyTorch al proceso del cliente, métrica comparable entre
  estrategias.}
  \label{tab:flex-flower-resultados}
  \begin{tabular}{lcccrcr}
    \toprule
    \textbf{Estrategia} & \textbf{Estado} & \textbf{Acc.} & \textbf{Loss} &
    \textbf{Dur. (s)} & \textbf{GPU proc. (MiB)} & \textbf{Energía (Wh)} \\
    \midrule
    FedAvg     & success & 0,4257 & 1,5265 & 132 & $\approx$289 & 1,272 \\
    FedAdam    & success & 0,1410 & 2,3401 & 124 & $\approx$289 & 1,252 \\
    FedAdagrad & success & 0,1000 & 18,7832 & 144 & $\approx$289 & 1,123 \\
    \bottomrule
  \end{tabular}
\end{table}

Estos resultados no permiten concluir que una estrategia sea superior a otra, por
dos motivos. En primer lugar, se trata de un mini-experimento de solo tres rondas
y una época local, insuficiente para que las estrategias adaptativas estabilicen
su segundo momento y converjan; con los hiperparámetros ajustados ya no divergen,
pero todavía no alcanzan el comportamiento de FedAvg. En segundo lugar, se ejecutó
con una única semilla, sin medida de variabilidad. Conviene además una
advertencia de medición: la memoria de GPU reportada por \texttt{nvidia-smi} a
nivel de sistema mostró un valor anómalo de 1623~MiB en la ejecución de
FedAdagrad frente a 757~MiB en las otras dos, atribuible a memoria ocupada por
otras aplicaciones del equipo durante esa corrida; el consumo real por proceso fue
prácticamente idéntico ($\approx$289~MiB) en las tres estrategias, por lo que para
cualquier comparación debe emplearse la métrica por proceso. La conclusión
experimental válida es que Flower permitió sustituir la estrategia federada entre
FedAvg, FedAdam y FedAdagrad sin modificar el cliente, el modelo ni el conjunto de
datos, demostrando su flexibilidad para alterar la lógica de agregación desde el
servidor y la configuración de ejecución.

\subsection{Integración de modelo, datos y métricas}\label{subsec:flex-flower-integracion}

Más allá de las estrategias, la evaluación confirmó que Flower integra sin
fricción un modelo y un conjunto de datos propios. El modelo es la ResNet-18
adaptada a CIFAR-10, construida en PyTorch sustituyendo la primera convolución por
una de núcleo 3$\times$3 y eliminando el \textit{maxpool} inicial, mientras que el
conjunto de datos se carga mediante \mintinline{python}{FederatedDataset} con un
particionador IID. El bucle de entrenamiento local es código propio del cliente:
emplea descenso de gradiente estocástico con momento y \textit{weight decay},
entropía cruzada como función de pérdida y un número de épocas, tasa y tamaño de
lote configurables. Además, el cliente registra métricas personalizadas, entre
ellas la pérdida de entrenamiento, el número de ejemplos, la duración y el pico de
memoria de GPU asignada por PyTorch, que el script de ejecución vuelca junto con
los registros completos en ficheros de resumen reutilizables. Esta combinación de
modelo propio, bucle de entrenamiento editable y métricas a medida sustenta una
valoración alta en los criterios de personalización, integración y
reproducibilidad.

\subsection{Capacidades avanzadas según la documentación}\label{subsec:flex-flower-doc}

Más allá de lo ejecutado, la documentación oficial de Flower (versión 1.31) describe
un conjunto amplio de técnicas de flexibilidad que se valoran como evidencia
documental. En el plano de la configuración y el control del servidor, se detallan
cuatro vías de personalización de las estrategias: ajustar los parámetros de una
estrategia existente (\texttt{fraction\_train}, \texttt{fraction\_evaluate},
\texttt{min\_available\_nodes}); inyectar funciones de devolución como
\texttt{evaluate\_fn} que el servidor ejecuta en momentos concretos; sobrescribir
métodos internos como \texttt{configure\_train}, \texttt{aggregate\_fit} o
\texttt{aggregate\_evaluate}; e implementar una estrategia nueva heredando de
\texttt{BaseStrategy}~\parencite{beutel2020flower}. La misma documentación muestra cómo
enviar configuraciones por ronda mediante \texttt{ConfigRecord} —por ejemplo, reducir
progresivamente la tasa de aprendizaje leyendo \texttt{server\_round}— y cómo agregar
métricas con criterios propios a través de \texttt{evaluate\_metrics\_aggr\_fn}, en
lugar del promedio ponderado por número de ejemplos. Los clientes, sin estado por
defecto, pueden conservar información entre rondas mediante \texttt{Context.state}, lo
que habilita técnicas como el \textit{checkpointing} en el cliente.

Una pieza distintiva son los \textit{Mods}, funciones que interceptan los mensajes
entre servidor y cliente para transformarlos sin alterar la lógica principal; se
encadenan en un orden definido y permiten, por ejemplo, recortar gradientes
(\texttt{fixedclipping\_mod}, \texttt{adaptiveclipping\_mod}) o limitar el tamaño de
los mensajes (\texttt{message\_size\_mod})~\parencite{beutel2020flower}. Sobre esta
base se construye la privacidad diferencial, documentada en modo preliminar en tres
variantes —central del lado del servidor, central del lado del cliente y local
mediante \texttt{LocalDpMod}—, todas con recorte y ruido gaussiano configurables, a las
que se suma la agregación segura. La documentación también ejemplifica esquemas no
estándar, como FedBN —que excluye los parámetros de normalización por lotes del
intercambio para preservar las estadísticas locales de cada cliente—, el aprendizaje
federado vertical y la integración con XGBoost.

La documentación recoge asimismo soporte para el entrenamiento asíncrono, si bien de
carácter experimental y fuera de las guías principales. Las notas de versión describen
una \textit{Driver API} experimental (versión 1.2) que ejecuta un servidor programable,
asíncrono y multi-inquilino, y la API \textit{Flower Next} (versión 1.8), un conjunto de
primitivas de bajo nivel que permite personalizar el bucle de entrenamiento e incluye,
de forma explícita, soporte para aprendizaje federado asíncrono y esquemas cíclicos,
con un \texttt{ServerApp} capaz de permanecer activo de forma
indefinida~\parencite{beutel2020flower}. Por su naturaleza en desarrollo y de bajo
nivel, esta capacidad se valora como soporte parcial y no como una funcionalidad
directa y estable. Todas estas capacidades se reconocen como existentes, pero no se
contabilizan como demostradas empíricamente en este trabajo.

\subsection{Ejecución, despliegue y operación según la documentación}\label{subsec:flex-flower-deploy}

La flexibilidad de Flower se extiende a su modelo de ejecución y despliegue. En local,
los comandos \texttt{flwr run} y \texttt{flwr list} levantan automáticamente un
\textit{SuperLink} que gestiona las ejecuciones, mientras que las simulaciones se
lanzan sobre el \textit{Simulation Runtime}, cuyo número de \textit{SuperNodes} y
recursos por cliente se fijan de forma persistente o se sobrescriben por ejecución;
la documentación contempla incluso simulaciones sobre un clúster Ray
multi-anfitrión~\parencite{beutel2020flower}. Para despliegues distribuidos, la
federación se compone de un \textit{SuperLink} que coordina y varios
\textit{SuperNodes} que ejecutan los clientes, modelo que la documentación lleva hasta
entornos de producción en la nube —Google Cloud sobre Kubernetes, Microsoft Azure y
Red Hat OpenShift— e incluso a la interconexión de varios clústeres mediante Skupper.

Este recorrido se acompaña de mecanismos de seguridad operativa que también cuentan
para la flexibilidad: cifrado de las comunicaciones mediante TLS con certificados
propios, autenticación de \textit{SuperNodes} basada en firmas y lista blanca de
claves públicas, autenticación de cuentas de usuario mediante OpenID Connect con
autorización fina y un registro de auditoría en formato JSON que traza qué actor
ejecuta cada acción~\parencite{beutel2020flower}. Por último, la interfaz de línea de
comandos puede emitir su salida en formato JSON mediante la opción
\texttt{-{}-format json} para integrarse en flujos automatizados, y las herramientas de
gestión de federaciones permiten inspeccionar nodos, ejecuciones y estado. Ninguna de
estas capacidades se ejercitó en los experimentos, pero refuerzan la valoración de
Flower en los modos de ejecución y despliegue y en la integración con herramientas
externas.

\subsection{Limitaciones observadas}\label{subsec:flex-flower-lim}

El hallazgo más relevante de la evaluación afecta a la memoria RAM del anfitrión y
no a la GPU. El motor de simulación de Flower delega en Ray, que al iniciarse
detecta los procesadores del equipo (veinte en este caso) y prelanza un proceso
\textit{worker} ocioso por cada uno; cada worker carga el intérprete y las
bibliotecas (en torno a 0,28~GB), de modo que el conjunto de workers, sumado al
proceso de simulación y al resto del sistema, agotaba los aproximadamente 14,6~GB
de RAM disponibles. Ray respondía terminando procesos con un error de memoria
agotada, lo que ocurría incluso en modo secuencial, donde a priori solo hay un
cliente entrenando a la vez. La solución fue limitar el número de procesadores que
ve la \textit{Simulation Runtime} mediante \texttt{init-args-num-cpus=2}, lo que
reduce drásticamente la huella de RAM. Se trata de un resultado comparativo de
interés: el backend basado en Ray presenta una huella de memoria de anfitrión
notablemente mayor que el simulador de NVIDIA FLARE o el backend de proceso único
de FedML, que se ejecutan en un único proceso, por lo que en equipos con poca RAM
la simulación de Flower obliga a acotar los recursos manualmente.

La ejecución concurrente, probada también mediante Ray con una fracción de GPU por
cliente de 0,5, quedó condicionada por esta misma limitación de memoria, en línea
con lo descrito en el capítulo de resultados de rendimiento. Por su parte, la
divergencia de las estrategias adaptativas con los valores por defecto, ya
comentada, es una limitación de uso y no del framework: las estrategias de la
familia FedOpt son sensibles a la tasa de aprendizaje del servidor y requieren un
ajuste deliberado, especialmente con pocas rondas~\parencite{reddi2021adaptive}.
Ninguna de estas limitaciones impidió completar la evaluación, pero todas matizan
la valoración final.

\subsection{Valoración de la flexibilidad de Flower}\label{subsec:flex-flower-score}

La Tabla~\ref{tab:flex-flower-score} resume la valoración de Flower sobre los diez
criterios F1--F10, distinguiendo el tipo de evidencia que sustenta cada
puntuación. Los criterios de personalización del modelo y los datos (F1), del
entrenamiento local (F2), de la estrategia de agregación (F3) y de los algoritmos
federados disponibles (F4) se valoran con la máxima puntuación a partir de
evidencia experimental directa, al igual que los modos de ejecución y despliegue
(F6), la integración con herramientas externas (F7) y la reproducibilidad (F10).
El soporte de aprendizaje asíncrono (F5), la seguridad y privacidad (F8) y la
modificación de la arquitectura federada (F9) reciben puntuación parcial por estar
respaldados solo por documentación; en el caso de F5, mediante interfaces de bajo
nivel todavía en desarrollo (la \textit{Driver API} y \textit{Flower Next}), por lo
que se reconoce el soporte sin elevarlo a la máxima puntuación.

\begin{table}[htbp]
  \centering
  \caption{Valoración de la flexibilidad de Flower según los criterios F1--F10.
  El tipo de evidencia distingue las capacidades ejecutadas en los experimentos
  (experimental) de las descritas en la documentación oficial (documental).}
  \label{tab:flex-flower-score}
  \begin{tabular}{clcl}
    \toprule
    \textbf{Criterio} & \textbf{Descripción} & \textbf{Punt.} & \textbf{Evidencia} \\
    \midrule
    F1  & Personalización del modelo y el dataset      & 2/2 & Experimental \\
    F2  & Personalización del entrenamiento local      & 2/2 & Experimental \\
    F3  & Estrategias de agregación personalizadas     & 2/2 & Experimental \\
    F4  & Algoritmos federados disponibles             & 2/2 & Experimental \\
    F5  & Aprendizaje federado asíncrono               & 1/2 & Documental \\
    F6  & Modos de ejecución y despliegue              & 2/2 & Exp.\ + documental \\
    F7  & Integración con herramientas externas        & 2/2 & Exp.\ + documental \\
    F8  & Seguridad, privacidad y extensiones          & 1/2 & Documental \\
    F9  & Modificación de la arquitectura federada     & 1/2 & Documental \\
    F10 & Reproducibilidad y configuración             & 2/2 & Experimental \\
    \midrule
    \multicolumn{2}{r}{\textbf{Total}} & \multicolumn{2}{l}{\textbf{17/20}} \\
    \bottomrule
  \end{tabular}
\end{table}

En conjunto, Flower obtiene una valoración alta en flexibilidad (17/20). La
evidencia experimental muestra una herramienta cómoda de adaptar:
permitió integrar un modelo y un conjunto de datos propios en PyTorch, editar el
bucle de entrenamiento local, sustituir la estrategia de agregación entre FedAvg,
FedAdam y FedAdagrad desde el servidor —incluyendo el ajuste de sus
hiperparámetros— y automatizar la recogida de métricas. La valoración no alcanza
el máximo porque varias capacidades avanzadas —el aprendizaje asíncrono, la
privacidad diferencial, la agregación segura o los esquemas no horizontales— solo se
verificaron documentalmente, y porque el soporte asíncrono se apoya en interfaces
todavía experimentales; a ello se añade que la viabilidad práctica del modo
concurrente quedó condicionada por la elevada huella de memoria del backend Ray en
el equipo empleado. Estas conclusiones se contrastan con las de NVIDIA FLARE y FedML
en la comparativa final del capítulo.

% =====================================================================
%  SECCIÓN AÑADIDA: Flexibilidad de NVIDIA FLARE
% =====================================================================

\section{Flexibilidad de NVIDIA FLARE}\label{sec:flex-nvflare}

NVIDIA FLARE se presenta como un SDK de aprendizaje federado agnóstico al dominio
y compatible con flujos de trabajo de aprendizaje automático ya existentes, con
una orientación que abarca desde la simulación local hasta el despliegue en
producción y en el borde~\parencite{roth2022nvflare,nvflareWelcome}. La evaluación
experimental de su flexibilidad se apoya en un paquete de cinco ejecuciones de la
variante reducida del experimento e3, la misma \texttt{e3\_flex\_mini} empleada con
Flower (Subsección~\ref{subsec:flex-flower-mini}), con ocho clientes totales,
cuatro por ronda, tres rondas, una época local, tamaño de lote 32, tasa de
aprendizaje 0,02 y semilla 42, en modo secuencial con un único hilo del simulador.
Conviene precisar el alcance: el paquete inspeccionado contiene una selección de
suites concretas y no la totalidad del histórico de ejecuciones, por lo que las
conclusiones experimentales se refieren únicamente a los archivos incluidos en él.

\subsection{Diseño modular: jobs, componentes y workflows}\label{subsec:flex-nvflare-modular}

La flexibilidad de NVIDIA FLARE descansa en su modelo de configuración modular. La
unidad de ejecución es el \textit{job}, que describe por completo un experimento
mediante ficheros de configuración de servidor y de cliente
(\texttt{config\_fed\_server.conf} y \texttt{config\_fed\_client.conf}), en los que
cada elemento del sistema se declara como un componente con su ruta de clase y sus
argumentos de construcción~\parencite{nvflareComponents}. Al desplegar el job, el
servidor y los clientes analizan esos ficheros e instancian dinámicamente los
componentes declarados, registrándolos en el bucle de eventos del framework; cada
componente se identifica con un \texttt{component\_id} que permite referenciarlo
desde otros componentes o workflows~\parencite{nvflareComponents}. Este diseño es
el que permite sustituir un agregador, un persistor, un \textit{executor}, un
generador de objetos compartibles (\textit{Shareable}) o un filtro modificando la
configuración, sin reescribir la aplicación completa.

Sobre esa base de bajo nivel, la documentación ofrece además una capa de
\textit{recipes} que empaqueta algoritmos y workflows habituales —FedAvg, FedProx,
FedOpt, SCAFFOLD, aprendizaje cíclico, XGBoost, Sklearn, estadísticas federadas,
evaluación cruzada, PSI, integración con Flower y swarm learning, entre
otros~\parencite{nvflareRecipes}. Esta dualidad, configuración granular por
componentes y configuración de alto nivel por recetas, sustenta una valoración
elevada en personalización y reproducibilidad, y constituye el punto de partida de
las adaptaciones realizadas para probar estrategias alternativas.

\subsection{Adaptaciones para soportar las estrategias}\label{subsec:flex-nvflare-adapt}

El objetivo de la prueba fue comprobar si un mismo job podía adaptarse para
ejecutar FedAvg, FedProx y variantes FedOpt del servidor (FedAdagrad y FedAdam),
manteniendo constante la estructura general del experimento. Para ello se
introdujeron tres cambios de naturaleza distinta. En primer lugar, la selección de
estrategia se trasladó al cliente, que admite los valores \texttt{fedavg},
\texttt{fedprox}, \texttt{fedadam} y \texttt{fedadagrad} mediante un argumento de
línea de órdenes que, por defecto, toma su valor de la variable de entorno
\texttt{TFG\_FL\_STRATEGY}. En segundo lugar, FedProx se implementó en el
entrenamiento local: cuando la estrategia es \texttt{fedprox} y el coeficiente
proximal es positivo, el entrenamiento añade un término que penaliza la desviación
de los pesos locales respecto al modelo global recibido al inicio de la tarea,
manteniendo el descenso de gradiente estocástico como optimizador
local~\parencite{nvflareAlgorithms}.

El tercer cambio afecta al servidor y es el más relevante para la flexibilidad de
agregación. Para las variantes FedOpt, el script de ejecución parchea el fichero de
configuración del servidor sustituyendo el generador de compartibles por defecto,
\mintinline{text}{FullModelShareableGenerator}, por
\mintinline{text}{PTFedOptModelShareableGenerator}, y añade un componente de modelo
fuente requerido por este generador, como ilustra el
Código~\ref{cod:nvflare-fedopt}. Este generador actualiza el modelo global mediante
un optimizador de PyTorch —y, opcionalmente, un planificador de tasa de
aprendizaje— cuyos argumentos incluyen \texttt{optimizer\_args},
\texttt{lr\_scheduler\_args}, \texttt{source\_model} y
\texttt{device}~\parencite{nvflareFedopt}. En consecuencia, FedAdam y FedAdagrad no
son recetas independientes, sino configuraciones de FedOpt que se obtienen
seleccionando \mintinline{text}{torch.optim.Adam} o
\mintinline{text}{torch.optim.Adagrad} como optimizador del servidor.

\begin{codigo}[title={Fragmento representativo de config\_fed\_server.conf (FedOpt)}, label={cod:nvflare-fedopt}]{json}
"shareable_generator": {
  "path": "nvflare.app_opt.pt.fedopt.PTFedOptModelShareableGenerator",
  "args": {
    "source_model": "model",
    "optimizer_args": { "path": "torch.optim.Adagrad" }
  }
},
"components": [
  { "id": "model", "path": "<ruta de la ResNet-18 adaptada>" }
]
\end{codigo}

Una precaución metodológica relevante surgió con el dispositivo de ejecución. En
todos los ficheros de configuración del cliente aparece \texttt{-{}-device cuda},
pero en la ejecución de FedAdam en CPU el propio cliente seleccionó CPU porque
\mintinline{python}{torch.cuda.is_available()} devolvió falso al lanzarse con
\texttt{CUDA\_VISIBLE\_DEVICES} vacío. Por ello, el dispositivo realmente empleado
debe leerse del registro de ejecución y no del fichero de configuración, criterio
que se ha seguido al elaborar la Tabla~\ref{tab:flex-nvflare-resultados}.

\subsection{Mini-experimento de flexibilidad}\label{subsec:flex-nvflare-mini}

La Tabla~\ref{tab:flex-nvflare-resultados} resume las cinco ejecuciones. FedAvg
sirvió de referencia y se ejecutó correctamente con el workflow \textit{Scatter and
Gather}, en el que el servidor distribuye los pesos iniciales, los clientes
entrenan localmente y devuelven sus pesos, y el sistema los agrega por promedio
ponderado~\parencite{nvflareAlgorithms}. A partir de esa base, FedProx finalizó las
tres rondas y los registros confirmaron \texttt{strategy=fedprox} con un coeficiente
proximal \texttt{fedprox\_mu=0,001} activo durante el entrenamiento. FedAdagrad
también completó las tres rondas como FedOpt real —y no como mera etiqueta— porque
la configuración del servidor incorporaba el generador
\mintinline{text}{PTFedOptModelShareableGenerator} y el registro mostraba las
actualizaciones del servidor con \mintinline{text}{torch.optim.Adagrad}.

\begin{table}[htbp]
  \centering
  \caption{Resultados del mini-experimento de flexibilidad de NVIDIA FLARE. Las
  cifras de aprendizaje son auxiliares (tres rondas, una época, una semilla) y no
  constituyen una comparación de rendimiento. El dispositivo real procede del
  registro de ejecución, no del fichero de configuración del cliente.}
  \label{tab:flex-nvflare-resultados}
  \begin{tabular}{llccccrr}
    \toprule
    \textbf{Estrategia} & \textbf{Estado} & \textbf{Disp.} & \textbf{FedOpt} &
    \textbf{Rondas} & \textbf{Dur. (s)} & \textbf{Acc.} & \textbf{Loss} \\
    \midrule
    FedAvg      & success & \texttt{cuda:0} & No  & 3 & 205  & 0,4178 & 1,5601 \\
    FedProx     & success & \texttt{cuda:0} & No  & 3 & 225  & 0,4045 & 1,5876 \\
    FedAdagrad  & success & \texttt{cuda:0} & Sí  & 3 & 199  & 0,1402 & 2,3785 \\
    FedAdam     & failed  & \texttt{cuda:0} & Sí  & 0 & 77   & ---    & ---    \\
    FedAdam     & success & \texttt{cpu}    & Sí  & 3 & 1280 & 0,2586 & 3,4439 \\
    \bottomrule
  \end{tabular}
\end{table}

El caso de FedAdam exige una lectura cuidadosa. La estrategia llegó a configurarse
realmente mediante FedOpt con \mintinline{text}{torch.optim.Adam}, pero su ejecución
en GPU abortó durante el paso de actualización del servidor
(\mintinline{python}{torch.optim.Adam.step()}) con un error de inicialización de
CUDA, registrado como \texttt{FATAL\_SYSTEM\_ERROR}. Para aislar el problema, la
prueba se repitió en CPU, donde FedAdam completó las tres rondas; esta ejecución
valida FedAdam como estrategia configurable, pero su duración (1280~s) es muy
superior precisamente por entrenarse en CPU, de modo que no debe mezclarse con la
comparativa de rendimiento en GPU del capítulo correspondiente. En coherencia con
ello, la fila de la ejecución fallida no reporta precisión ni pérdida, ya que no
completó ninguna ronda.

Estos resultados no permiten ordenar las estrategias por calidad, por los mismos
motivos señalados para Flower: tres rondas, una época local y una sola semilla son
insuficientes para que las variantes adaptativas converjan. Su valor es
demostrar que NVIDIA FLARE permite adaptar tanto el entrenamiento local como la
lógica de actualización del servidor mediante configuración y componentes. Importa
además una cautela de interpretación: la validez de FedAdagrad y FedAdam no se
sustenta en el nombre de la ejecución, sino en la presencia del generador FedOpt y
del optimizador del servidor en los registros. Del mismo modo, el valor
\texttt{fedprox\_mu=0,001} que aparece en la configuración no implica que FedProx se
aplique a todas las estrategias: en FedAdagrad y FedAdam el entrenamiento imprime
\texttt{fedprox\_mu=0,0}, porque el término proximal solo se activa cuando la
estrategia es FedProx.

De esta experiencia se desprende una observación comparativa de interés. Frente a
Flower, donde la estrategia se sustituía cambiando una única clase en el servidor,
NVIDIA FLARE exigió modificar componentes del servidor, añadir un modelo fuente para
FedOpt y comprobar en los registros que la estrategia configurada era realmente la
que se ejecutaba. La flexibilidad de NVIDIA FLARE es, por tanto, amplia, pero su
ejercicio práctico conlleva una complejidad de configuración mayor.

\subsection{Capacidades avanzadas según la documentación}\label{subsec:flex-nvflare-doc}

Como en el caso de Flower, varias capacidades relevantes no se ejecutaron en los
experimentos y se valoran únicamente como evidencia documental. La documentación
oficial las articula en torno a las \textit{Job Recipes}, una API de alto nivel que
encapsula un trabajo federado exponiendo solo los parámetros clave —\texttt{min\_clients},
\texttt{num\_rounds}, modelo y \textit{script} de entrenamiento— y ocultando la
complejidad de controladores y \textit{workflows}~\parencite{nvflareRecipes}. Estas
recetas se ejecutan sobre tres entornos intercambiables: \texttt{SimEnv} simula los
clientes en procesos locales, \texttt{PocEnv} los lanza como procesos separados en la
misma máquina y \texttt{ProdEnv} despliega el trabajo en una instalación de
producción; un mismo trabajo puede así trasladarse de la simulación al despliegue sin
cambiar el código~\parencite{nvflareRecipes}. El catálogo de recetas cubre FedAvg,
FedProx, FedOpt —donde el optimizador del servidor se fija con un argumento que indica
su ruta de clase e hiperparámetros, por ejemplo \texttt{torch.optim.SGD} con tasa y
momento—, SCAFFOLD y el aprendizaje cíclico, además de XGBoost federado en sus
variantes horizontal, \textit{bagging} y vertical~\parencite{nvflareAlgorithms,nvflareFedopt}.

En cuanto al aprendizaje asíncrono, la receta \texttt{EdgeFedBuffRecipe}, orientada a
dispositivos de borde, admite entrenamiento síncrono y asíncrono: en modo síncrono el
servidor espera todas las actualizaciones antes de generar un nuevo modelo, mientras
que en modo asíncrono lo genera en cuanto recibe la primera, lo que reduce el tiempo
total cuando los dispositivos presentan latencias heterogéneas~\parencite{nvflareEdge,nvflareWelcome}.
La receta expone parámetros como el número máximo de versiones de modelo activas o el
tamaño de selección de dispositivos por ronda. No obstante, la propia documentación
advierte de que la asincronía general aún no está disponible para todas las recetas y
se concentra en el escenario de borde, por lo que su valoración se mantiene como
soporte documental parcial.

En materia de seguridad y privacidad, la documentación es notablemente extensa.
NVIDIA FLARE implementa la privacidad mediante un mecanismo general de filtros que
procesan los objetos compartibles antes o después de la comunicación, configurables en
la receta como \texttt{task\_data\_filters} o \texttt{task\_result\_filters}. Soporta
privacidad diferencial a nivel de muestra mediante la integración de DP-SGD con Opacus
—controlada por los parámetros \texttt{epsilon} y \texttt{delta}— y filtros de DP local
como \texttt{PercentilePrivacy}, que comparte solo las diferencias de pesos por encima
de un percentil, \texttt{SVTPrivacy}, basado en la técnica del vector disperso con ruido
de Laplace, y \texttt{StatisticsPrivacyFilter} para estadísticas
federadas~\parencite{nvflareDP}. A un nivel superior, la solución de computación
confidencial ejecuta la agregación dentro de un entorno de ejecución confiable
(\textit{Trusted Execution Environment}) sobre hardware como AMD SEV-SNP o Intel TDX,
protege el modelo y el código del cliente en máquinas virtuales confidenciales y
permite combinar privacidad diferencial con cifrado homomórfico o agregación
segura~\parencite{nvflareSecurity,nvflareConfidential}.

Por último, NVIDIA FLARE documenta esquemas descentralizados mediante los
\textit{workflows} controlados por clientes, construidos sobre las clases base
\texttt{ServerSideController} y \texttt{ClientSideController} para escenarios donde el
servidor no es plenamente confiable~\parencite{nvflareSwarm}. Entre ellos destacan el
aprendizaje cíclico, en el que cada cliente entrena con el modelo recibido del anterior
según un orden fijo o aleatorio~\parencite{nvflareCyclic}; el aprendizaje \textit{swarm},
que designa un cliente agregador distinto por ronda y mejora la resiliencia frente a
fallos del servidor~\parencite{nvflareSwarm}; y la evaluación cruzada entre sitios, que
permite a los clientes validar modelos de otros mediante comunicación entre pares sin
exponerlos al servidor. A esto se suman el aprendizaje vertical con XGBoost y las
estadísticas federadas, así como las utilidades para integrar el seguimiento de
experimentos con TensorBoard, MLflow o Weights \& Biases sobre cualquier
receta~\parencite{nvflareTracking,nvflareRecipes}. La documentación reconoce, además,
dos cautelas: la API de recetas se encuentra en estado de \textit{technical preview} y
no cubre todavía todos los algoritmos, y la configuración de la computación confidencial
exige infraestructura especializada.

\subsection{Limitaciones observadas}\label{subsec:flex-nvflare-lim}

La limitación más concreta de la evaluación fue el fallo de FedAdam en GPU por un
error de inicialización de CUDA durante la actualización FedOpt del servidor. Debe
interpretarse como una incidencia técnica del entorno, no como una carencia del
framework, puesto que la misma configuración completó las tres rondas en CPU; aun
así, impide presentar FedAdam como ejecución correcta en GPU. La segunda limitación
es la complejidad práctica ya señalada: a diferencia de Flower, alterar la lógica de
agregación obligó a intervenir en componentes del servidor y a validar
cuidadosamente los registros, lo que incrementa la probabilidad de errores de
configuración. Por último, el modo concurrente quedó condicionado por los recursos
del equipo, en línea con lo expuesto en el capítulo de resultados de rendimiento;
aunque la documentación describe un sistema de ejecución multi-job con gestión de
recursos, esa capacidad no pudo validarse de forma estable en el hardware empleado.

\subsection{Valoración de la flexibilidad de NVIDIA FLARE}\label{subsec:flex-nvflare-score}

La Tabla~\ref{tab:flex-nvflare-score} recoge la valoración sobre los criterios
F1--F10, con el mismo convenio aplicado a Flower: puntuación máxima cuando existe
evidencia experimental directa, puntuación parcial cuando la capacidad solo consta
en la documentación y puntuación nula cuando no se identifica soporte. NVIDIA FLARE
obtiene la máxima puntuación en personalización del modelo y los datos (F1), del
entrenamiento local (F2), de la estrategia de agregación (F3) y de los algoritmos
federados disponibles (F4), todos respaldados por las ejecuciones, así como en los
modos de ejecución y despliegue (F6), la integración con herramientas externas (F7)
y la reproducibilidad (F10). La diferencia respecto a Flower aparece en el
aprendizaje asíncrono (F5): mientras que en Flower no se identificó soporte, NVIDIA
FLARE lo documenta de forma explícita mediante FedBuff, lo que justifica una
puntuación parcial. La seguridad y privacidad (F8) y la modificación de la
arquitectura federada (F9) reciben también puntuación parcial por estar respaldadas
solo documentalmente.

\begin{table}[htbp]
  \centering
  \caption{Valoración de la flexibilidad de NVIDIA FLARE según los criterios
  F1--F10. El tipo de evidencia distingue las capacidades ejecutadas en los
  experimentos de las descritas en la documentación oficial.}
  \label{tab:flex-nvflare-score}
  \begin{tabular}{clcl}
    \toprule
    \textbf{Criterio} & \textbf{Descripción} & \textbf{Punt.} & \textbf{Evidencia} \\
    \midrule
    F1  & Personalización del modelo y el dataset      & 2/2 & Exp.\ + documental \\
    F2  & Personalización del entrenamiento local      & 2/2 & Experimental \\
    F3  & Estrategias de agregación personalizadas     & 2/2 & Experimental \\
    F4  & Algoritmos federados disponibles             & 2/2 & Exp.\ + documental \\
    F5  & Aprendizaje federado asíncrono               & 1/2 & Documental \\
    F6  & Modos de ejecución y despliegue              & 2/2 & Exp.\ + documental \\
    F7  & Integración con herramientas externas        & 2/2 & Exp.\ + documental \\
    F8  & Seguridad, privacidad y extensiones          & 1/2 & Documental \\
    F9  & Modificación de la arquitectura federada     & 1/2 & Documental \\
    F10 & Reproducibilidad y configuración             & 2/2 & Exp.\ + documental \\
    \midrule
    \multicolumn{2}{r}{\textbf{Total}} & \multicolumn{2}{l}{\textbf{17/20}} \\
    \bottomrule
  \end{tabular}
\end{table}

En conjunto, NVIDIA FLARE alcanza una valoración de flexibilidad de 17/20. La
evidencia experimental
demuestra que permite adaptar el entrenamiento local —incluido el término proximal
de FedProx— y, sobre todo, la lógica de agregación del servidor mediante FedOpt, con
FedAvg, FedProx y FedAdagrad ejecutados correctamente y FedAdam validado en CPU. A
ello se suma una de las superficies de capacidades documentadas más amplias del
estudio, especialmente sólida en seguridad y privacidad, que abarca aprendizaje
asíncrono, workflows cíclicos y \textit{swarm}, privacidad
diferencial, cifrado homomórfico y computación confidencial. El contrapunto es una
complejidad de configuración superior a la de Flower y el hecho de que esas
capacidades avanzadas solo se verificaron documentalmente; además, la ejecución de
FedAdam en GPU falló por un error de CUDA y el modo concurrente quedó limitado por el
hardware. Esta lectura se contrasta con la de los otros frameworks en la comparativa
final del capítulo.

% =====================================================================
%  SECCIÓN AÑADIDA: Flexibilidad de FedML
% =====================================================================

\section{Flexibilidad de FedML}\label{sec:flex-fedml}

FedML es una biblioteca de investigación y banco de pruebas para aprendizaje
federado~\parencite{he2020fedml}, documentada actualmente dentro de la plataforma
TensorOpera aunque la librería y los ejemplos siguen empleando el paquete
\texttt{fedml}. Su documentación organiza el framework en varios escenarios
—simulación, cross-silo y cross-device— y ofrece un catálogo amplio de algoritmos
de referencia. La evaluación de su flexibilidad combina la implementación propia
realizada para la comparativa, que extiende las abstracciones de FedML, con una
prueba reducida de sustitución de estrategias. Conviene una precisión de alcance:
por el coste computacional, la prueba reducida de FedML se ejecutó con cuatro
clientes totales, dos por ronda y dos rondas, una configuración algo menor que la
empleada con Flower y NVIDIA FLARE; las métricas de aprendizaje son, por tanto, una
comprobación funcional y no una comparación de rendimiento.

\subsection{Configuración y componentes personalizados}\label{subsec:flex-fedml-config}

La configuración de FedML se apoya en un fichero \texttt{fedml\_config.yaml}
estructurado en secciones (\texttt{common\_args}, \texttt{data\_args},
\texttt{model\_args}, \texttt{train\_args}, \texttt{device\_args},
\texttt{comm\_args} y \texttt{tracking\_args}), donde se declaran los parámetros
federados como \texttt{federated\_optimizer}, \texttt{client\_num\_in\_total},
\texttt{client\_num\_per\_round}, \texttt{comm\_round}, \texttt{epochs},
\texttt{batch\_size}, \texttt{client\_optimizer} y \texttt{learning\_rate}~\parencite{fedmlSimApi}.
Esta organización permitió variar clientes, rondas, épocas e hiperparámetros entre
ejecuciones sin reescribir la aplicación.

Más allá del fichero de configuración, la flexibilidad de FedML en este trabajo se
sustenta en evidencia experimental directa: la implementación de la comparativa
define un entrenador y un agregador propios que extienden las abstracciones
\mintinline{python}{ClientTrainer} y \mintinline{python}{ServerAggregator} del
framework. El entrenador \texttt{FedMLCifar10Trainer} encapsula el bucle de
entrenamiento local sobre la ResNet-18 adaptada —descenso de gradiente estocástico,
entropía cruzada y métricas propias—, mientras que el agregador
\texttt{FedMLCifar10Aggregator} gobierna la evaluación del modelo global; ambos se
conectan al flujo mediante \texttt{FedMLRunner}. La documentación respalda esta vía,
al indicar que el usuario puede definir modelo, datos, optimizadores, funciones de
pérdida y lógica de entrenamiento arbitrarios, así como dataloaders propios que
respeten la estructura esperada por el simulador~\parencite{fedmlTrainer,fedmlDataloader}.
Esta capacidad de sustituir componentes centrales mediante herencia sostiene una
valoración alta en personalización del entrenamiento local, de la agregación y de la
integración.

\subsection{Sustitución de la estrategia federada}\label{subsec:flex-fedml-strategy}

La prueba de flexibilidad consistió en ejecutar cuatro estrategias modificando
únicamente la configuración, manteniendo constantes el conjunto de datos, el
modelo y el resto de hiperparámetros. FedAvg y FedProx se seleccionan asignando
\texttt{federated\_optimizer} a \texttt{"FedAvg"} o \texttt{"FedProx"}, mientras que
FedAdam y FedAdagrad se obtienen mediante \texttt{federated\_optimizer:\ "FedOpt"}
combinado con un optimizador de servidor \texttt{Adam} o \texttt{Adagrad}, tal como
ilustra el Código~\ref{cod:fedml-config}. Durante las primeras pruebas, las
variantes FedOpt fallaban porque FedML esperaba el atributo \texttt{ci} en los
argumentos de configuración; tras añadir \texttt{ci:\ 0} a la configuración común,
ambas pudieron completar la ejecución. La selección del algoritmo desde la
configuración está respaldada por la documentación, que controla el optimizador
federado a través del mismo campo~\parencite{fedmlSimApi}.

\begin{yamlcode}[title={Fragmento de fedml\_config.yaml (FedAdam vía FedOpt)}, label={cod:fedml-config}]
train_args:
  federated_optimizer: "FedOpt"
  server_optimizer: "Adam"
  ci: 0
  client_num_in_total: 4
  client_num_per_round: 2
  comm_round: 2
  epochs: 1
  batch_size: 32
  learning_rate: 0.02
\end{yamlcode}

\subsection{Mini-experimento de flexibilidad}\label{subsec:flex-fedml-mini}

La Tabla~\ref{tab:flex-fedml-resultados} resume las cuatro ejecuciones, todas
finalizadas con estado \texttt{success}. FedAvg y FedProx concluyeron con métricas
coherentes para una prueba de dos rondas, con precisiones del 43,01\,\% y el
42,60\,\% respectivamente, lo que confirma el soporte experimental de una estrategia
alternativa a FedAvg con comportamiento correcto. FedAdam y FedAdagrad también
completaron la ejecución, lo que demuestra que fue posible activar FedOpt con
optimizadores de servidor Adam y Adagrad; sin embargo, sus métricas finales fueron
claramente deficientes, con una precisión próxima al 10\,\% —equivalente al azar en
un problema de diez clases— y pérdidas muy elevadas.

\begin{table}[htbp]
  \centering
  \caption{Resultados del mini-experimento de flexibilidad de FedML. Las cifras de
  aprendizaje son auxiliares (dos rondas, una época, una semilla) y constituyen una
  comprobación funcional, no una comparación de rendimiento entre estrategias.}
  \label{tab:flex-fedml-resultados}
  \begin{tabular}{llrrrrr}
    \toprule
    \textbf{Estrategia} & \textbf{Estado} & \textbf{Dur. (s)} & \textbf{Acc.} &
    \textbf{Loss} & \textbf{GPU (MiB)} & \textbf{Util. (\%)} \\
    \midrule
    FedAvg     & success & 118 & 0,4301 & 1,5248   & 416 & 75,29 \\
    FedProx    & success & 123 & 0,4260 & 1,5137   & 484 & 79,01 \\
    FedAdam    & success & 117 & 0,0998 & 809,7697 & 416 & 76,09 \\
    FedAdagrad & success & 116 & 0,0998 & 191,2948 & 416 & 76,85 \\
    \bottomrule
  \end{tabular}
\end{table}

La lectura correcta de estos datos es que FedML permite cambiar la estrategia
federada mediante configuración, pero que la mera disponibilidad de una estrategia
no garantiza su estabilidad ni un buen rendimiento. Las variantes FedOpt
ejecutaron de principio a fin, de modo que la flexibilidad queda demostrada, pero su
falta de aprendizaje en esta prueba indica que requieren un ajuste específico de
hiperparámetros —como la tasa de aprendizaje del servidor o un mayor número de
rondas— para resultar competitivas. Este patrón es coherente con lo observado en los
otros frameworks: en Flower las estrategias adaptativas divergían con los valores por
defecto y exigieron ajustar \texttt{eta} y \texttt{tau}, y en NVIDIA FLARE FedAdam no
llegó a completarse en GPU. La fragilidad de las variantes FedOpt con poca
configuración es, por tanto, un comportamiento transversal y no una carencia
particular de FedML.

\subsection{Capacidades avanzadas según la documentación}\label{subsec:flex-fedml-doc}

La documentación de FedML respalda una superficie de capacidades más amplia que la
probada experimentalmente, organizada en varias capas. En el plano de uso, ofrece una
interfaz de línea de comandos —integrada en la plataforma TensorOpera— con órdenes para
autenticación y gestión de dispositivos (\texttt{fedml login}, \texttt{fedml device
bind}), lanzamiento y gestión de trabajos y clústeres (\texttt{fedml launch},
\texttt{fedml cluster}, \texttt{fedml run}), y gestión y despliegue de modelos como
servicios de inferencia (\texttt{fedml model})~\parencite{fedmlExamples}. En el plano de
programación, expone una API de Python con una llamada de una línea
(\texttt{fedml.run\_simulation()}) y una API más explícita de varias líneas que
inicializa dispositivo, datos y modelo y ejecuta la federación mediante
\texttt{SimulatorMPI} o las clases \texttt{Client}/\texttt{Server} envueltas por
\texttt{FedMLRunner}~\parencite{fedmlSimApi}.

En cuanto al catálogo algorítmico, el simulador Parrot enumera implementaciones de
referencia como FedOpt, FedNova, FedGKT, aprendizaje federado descentralizado, vertical
y jerárquico, FedNAS y \textit{Split Learning}, seleccionables cambiando el campo
\texttt{federated\_optimizer}~\parencite{fedmlSimOverview,fedmlSimApi}, y el módulo de
seguridad declara compatibilidad con optimizadores como FedAvg, FedSGD, FedOpt, FedProx,
FedGKT y FedNova~\parencite{fedmlSecurity}. A este catálogo se añaden la analítica
federada, que calcula estadísticas distribuidas como promedios o percentiles mediante un
ejecutor específico sin compartir datos crudos, y un proyecto independiente orientado al
ajuste federado de grandes modelos de lenguaje, que reutiliza las mismas API y la
plataforma de operaciones~\parencite{fedmlExamples}. En cuanto al aprendizaje asíncrono,
FedML Octopus documenta un escenario jerárquico heterogéneo en el que cada silo puede
ejecutar entrenamiento distribuido con varias GPU y orquestarse mediante optimización
federada asíncrona o síncrona~\parencite{fedmlCrossSilo}; se trata, no obstante, de un
soporte menos específico que el de NVIDIA FLARE, ya que no se documenta una receta directa
equivalente a FedBuff, aprendizaje cíclico o \textit{swarm}.

En privacidad y seguridad, el perfil de FedML difiere del de NVIDIA FLARE. Su
fortaleza documental está en la simulación de ataques y defensas mediante
\texttt{FedMLAttacker} y \texttt{FedMLDefender}, con ataques de tipo bizantino,
\textit{label flipping} y reemplazo de modelo, y defensas como recorte de norma, Krum,
mediana geométrica o FoolsGold, activables desde el fichero de configuración con
\texttt{enable\_attack} y \texttt{enable\_defense}~\parencite{fedmlSecurity}, además de un
ejemplo de agregación segura basada en computación multiparte~\parencite{fedmlExamples}.
En cambio, en las páginas revisadas no figura una integración de privacidad diferencial o
cifrado homomórfico comparable a la de NVIDIA FLARE, por lo que su valoración en este
criterio se sostiene en el \textit{benchmarking} de robustez más que en la privacidad
criptográfica. Por último, para integración y trazabilidad, FedML documenta el registro
de métricas, artefactos y modelos mediante \texttt{fedml.log()}, la captura automática de
métricas de hardware y GPU, los backends de comunicación MPI, gRPC, PyTorch RPC y MQTT+S3,
y la integración con la plataforma de operaciones de TensorOpera para gestionar agentes,
experimentos y despliegues~\parencite{fedmlTracking,fedmlExamples}.

\subsection{Limitaciones observadas}\label{subsec:flex-fedml-lim}

La limitación más visible es que las variantes FedOpt finalizaron sin aprender bajo
la configuración empleada, lo que confirma que su soporte es real pero condicionado a
un ajuste de hiperparámetros que queda fuera del alcance de esta prueba de
flexibilidad. A ello se añade el detalle de configuración del atributo \texttt{ci},
que tuvo que fijarse manualmente para que FedOpt funcionara, y el hecho de que la
prueba reducida de FedML emplea una configuración menor que la de los otros dos
frameworks, por lo que sus métricas auxiliares no son directamente comparables. En
cuanto a la ejecución, el modo concurrente de FedML mediante el backend MPI fue el
único que llegó a completarse entre los tres frameworks, según se documenta en el
capítulo de resultados de rendimiento, si bien se mostró sensible al consumo de
memoria en las configuraciones más grandes.

\subsection{Valoración de la flexibilidad de FedML}\label{subsec:flex-fedml-score}

La Tabla~\ref{tab:flex-fedml-score} recoge la valoración de FedML sobre los criterios
F1--F10, con el convenio ya aplicado a los otros frameworks. FedML obtiene la máxima
puntuación en personalización del modelo y los datos (F1), del entrenamiento local
(F2) y de la agregación (F3), respaldadas por la implementación de entrenador y
agregador propios; en los algoritmos federados disponibles (F4), por haber ejecutado
cuatro estrategias mediante configuración; en los modos de ejecución (F6), por ser el
único framework cuyo modo concurrente se completó; y en integración (F7) y
reproducibilidad (F10). El aprendizaje asíncrono (F5), la seguridad y privacidad (F8)
y la modificación de la arquitectura federada (F9) reciben puntuación parcial por
estar respaldados únicamente por la documentación.

\begin{table}[htbp]
  \centering
  \caption{Valoración de la flexibilidad de FedML según los criterios F1--F10. El
  tipo de evidencia distingue las capacidades ejecutadas en los experimentos de las
  descritas en la documentación oficial.}
  \label{tab:flex-fedml-score}
  \begin{tabular}{clcl}
    \toprule
    \textbf{Criterio} & \textbf{Descripción} & \textbf{Punt.} & \textbf{Evidencia} \\
    \midrule
    F1  & Personalización del modelo y el dataset      & 2/2 & Exp.\ + documental \\
    F2  & Personalización del entrenamiento local      & 2/2 & Exp.\ + documental \\
    F3  & Estrategias de agregación personalizadas     & 2/2 & Experimental \\
    F4  & Algoritmos federados disponibles             & 2/2 & Exp.\ + documental \\
    F5  & Aprendizaje federado asíncrono               & 1/2 & Documental \\
    F6  & Modos de ejecución y despliegue              & 2/2 & Exp.\ + documental \\
    F7  & Integración con herramientas externas        & 2/2 & Exp.\ + documental \\
    F8  & Seguridad, privacidad y extensiones          & 1/2 & Documental \\
    F9  & Modificación de la arquitectura federada     & 1/2 & Documental \\
    F10 & Reproducibilidad y configuración             & 2/2 & Exp.\ + documental \\
    \midrule
    \multicolumn{2}{r}{\textbf{Total}} & \multicolumn{2}{l}{\textbf{17/20}} \\
    \bottomrule
  \end{tabular}
\end{table}

En conjunto, FedML alcanza una valoración de flexibilidad de 17/20. Su mayor fortaleza
experimental reside en la facilidad para sustituir componentes centrales del sistema
—entrenador y agregador— mediante herencia, y en haber ejecutado el abanico más amplio
de estrategias del estudio a través de la configuración. Su perfil documental es
asimismo notable, con un catálogo extenso de algoritmos y escenarios y un soporte de
seguridad orientado al benchmarking de ataques y defensas. Las cautelas son dos: las
variantes FedOpt completaron la ejecución pero no aprendieron sin un ajuste
adicional, y algunas capacidades avanzadas —asincronía, privacidad criptográfica y
esquemas alternativos— solo se verificaron documentalmente.

% =====================================================================
%  SECCIÓN AÑADIDA: Comparativa de flexibilidad entre frameworks
% =====================================================================

\section{Comparativa de flexibilidad}\label{sec:flex-comparativa}

Una vez evaluados los tres frameworks por separado, esta sección los contrasta sobre
el mismo conjunto de criterios. La Tabla~\ref{tab:flex-comparativa} reúne las
puntuaciones F1--F10 de Flower, NVIDIA FLARE y FedML. Los tres obtienen la misma
valoración, 17/20, lo que confirma que las tres herramientas son, en lo esencial,
comparablemente flexibles para el tipo de adaptaciones que exige este trabajo y que la
diferenciación entre ellas debe buscarse, más que en la cifra global, en el perfil
cualitativo de cada una.

\begin{table}[htbp]
  \centering
  \caption{Comparativa de la valoración de flexibilidad de los tres frameworks según
  los criterios F1--F10 definidos en la metodología.}
  \label{tab:flex-comparativa}
  \begin{tabular}{clccc}
    \toprule
    \textbf{Criterio} & \textbf{Descripción} & \textbf{Flower} & \textbf{NVIDIA FLARE} & \textbf{FedML} \\
    \midrule
    F1  & Modelo y dataset                  & 2/2 & 2/2 & 2/2 \\
    F2  & Entrenamiento local               & 2/2 & 2/2 & 2/2 \\
    F3  & Agregación personalizada          & 2/2 & 2/2 & 2/2 \\
    F4  & Algoritmos disponibles            & 2/2 & 2/2 & 2/2 \\
    F5  & Aprendizaje asíncrono             & 1/2 & 1/2 & 1/2 \\
    F6  & Modos de ejecución y despliegue   & 2/2 & 2/2 & 2/2 \\
    F7  & Integración externa               & 2/2 & 2/2 & 2/2 \\
    F8  & Seguridad y privacidad            & 1/2 & 1/2 & 1/2 \\
    F9  & Arquitectura federada             & 1/2 & 1/2 & 1/2 \\
    F10 & Reproducibilidad y configuración  & 2/2 & 2/2 & 2/2 \\
    \midrule
    \multicolumn{2}{r}{\textbf{Total}} & \textbf{17/20} & \textbf{17/20} & \textbf{17/20} \\
    \bottomrule
  \end{tabular}
\end{table}

La coincidencia en los criterios experimentales (F1--F4, F6, F7 y F10) refleja un
hecho de fondo: bajo un mismo escenario de modelo, datos y algoritmo base, los tres
frameworks permitieron integrar la ResNet-18 sobre CIFAR-10, editar el entrenamiento
local, sustituir la estrategia de agregación, registrar métricas y reproducir los
experimentos mediante ficheros y scripts. La igualdad también alcanza a los tres
criterios documentales: ninguno validó experimentalmente el aprendizaje asíncrono (F5),
la seguridad avanzada (F8) ni esquemas alternativos a la arquitectura clásica (F9), pero
los tres los documentan, por lo que reciben idéntica puntuación parcial. En F5, en
concreto, los tres declaran soporte —Flower mediante la \textit{Driver API} y
\textit{Flower Next}, NVIDIA FLARE mediante FedBuff en el borde y FedML en su escenario
jerárquico cross-silo—, si bien en todos los casos se trata de capacidades
experimentales, limitadas a entornos concretos o de bajo nivel.

Más allá de las cifras, el estudio revela un comportamiento transversal en las
estrategias adaptativas de la familia FedOpt. En los tres frameworks fue posible
configurar FedAdam y FedAdagrad, pero su ejecución resultó frágil con poca
configuración: en Flower divergían con los valores por defecto y necesitaron ajustar
la tasa del servidor y el factor de adaptabilidad; en NVIDIA FLARE, FedAdam no llegó a
completarse en GPU por un error de CUDA; y en FedML, FedAdam y FedAdagrad finalizaron
pero sin aprender. La conclusión es clara y común: la disponibilidad de una estrategia
no equivale a su estabilidad, que depende de un ajuste deliberado de hiperparámetros,
especialmente con pocas rondas.

Por último, cada framework muestra un perfil de flexibilidad propio que la puntuación,
por su granularidad, no captura del todo. Flower es el más directo para sustituir la
estrategia de agregación —basta con cambiar una clase en el servidor— y su documentación
revela una superficie amplia que abarca \textit{Mods}, privacidad diferencial, FedBN y,
sobre todo, un modelo de despliegue muy desarrollado hacia la nube (Google Cloud, Azure,
OpenShift e interconexión de clústeres); su contrapartida es que el soporte asíncrono
sigue siendo experimental y que el backend Ray arrastra una huella de memoria elevada.
NVIDIA FLARE ofrece la mayor profundidad en seguridad y privacidad —filtros de privacidad
diferencial, cifrado homomórfico y computación confidencial sobre entornos de ejecución
confiables— junto con un catálogo de recetas y entornos (\texttt{SimEnv}, \texttt{PocEnv},
\texttt{ProdEnv}) orientado a producción, a costa de una complejidad de configuración
mayor, pues alterar la agregación exige intervenir en componentes del servidor. FedML
destaca por la sustitución de entrenador y agregador mediante herencia, por haber
ejecutado el mayor número de estrategias, por ser el único cuyo modo concurrente se
completó, y por una integración estrecha con su plataforma de operaciones y la analítica
federada, con un soporte de seguridad orientado al \textit{benchmarking} de ataques y
defensas. La elección entre ellos, en términos de flexibilidad, depende por tanto de la
prioridad del proyecto: la simplicidad y el despliegue en la nube de Flower, la
profundidad en seguridad y la orientación a producción de NVIDIA FLARE, o la
extensibilidad por componentes y la experimentación algorítmica de FedML.

%  (dentro de \mainmatter, tras el capítulo de resultados de rendimiento)
% =====================================================================

\chapter{Resultados de flexibilidad}\label{chap:flexibilidad}

Este capítulo presenta el tercer eje de la comparación: la flexibilidad de
cada framework. A diferencia de la comparativa de rendimiento, que es
cuantitativa, y de la valoración de usabilidad, que recoge la percepción de uso,
la flexibilidad describe una propiedad arquitectónica: hasta qué punto la
herramienta permite adaptarse a escenarios distintos de los previstos por defecto.
Conviene recordar que este eje no debe confundirse con la participación parcial
de clientes medida en el experimento e3, que es una variación cuantitativa del
escenario; aquí el término flexibilidad se reserva para la capacidad de
modificar, sustituir o ampliar el comportamiento del framework.

La valoración se apoya en los diez criterios F1--F10 definidos en la
Sección~\ref{sec:flex-criterios} y se puntúa con la escala de tres niveles
(0, 1, 2) de la Tabla~\ref{tab:escala-flexibilidad}. Para cada framework se
distinguen dos planos de evidencia. El primero es la evidencia experimental:
cambios efectivamente realizados en el código, la configuración y los scripts, y
ejecuciones reales sobre el escenario común. El segundo es la evidencia
documental: funcionalidades avanzadas descritas en la documentación oficial pero
no implementadas en los experimentos del trabajo. Esta separación es deliberada y
evita presentar como probado aquello que solo se ha revisado en la documentación,
manteniendo así la comparabilidad del diseño experimental.

\section{Flexibilidad de Flower}\label{sec:flex-flower}

Flower es un framework agnóstico respecto a la biblioteca de aprendizaje
automático y con una separación clara entre la lógica del cliente y la del
servidor~\parencite{beutel2020flower}, rasgos que condicionan directamente su
flexibilidad. La evaluación experimental se centró en cinco capacidades: modificar
la configuración federada, sustituir la estrategia de agregación, integrar un
modelo y un conjunto de datos propios en PyTorch, registrar métricas personalizadas
y alternar el modo de ejecución. Todas las ejecuciones se realizaron con
\texttt{flwr} 1.29.0 sobre el backend de simulación Ray~\parencite{moritz2018ray},
con Python 3.12 (entorno \texttt{tfg-flower}), una GPU NVIDIA GeForce RTX~3050
Laptop de 6\,GB y el modelo ResNet-18 adaptado sobre CIFAR-10 descrito en el
capítulo de metodología.

\subsection{Adaptaciones para habilitar la evaluación}\label{subsec:flex-flower-adapt}

La primera evidencia de flexibilidad es la propia parametrización del escenario.
Para que Flower recibiera la configuración de federación en cada ejecución, se
ajustó el script de lanzamiento de modo que toda la configuración viajara en una
única cadena pasada a \texttt{flwr run}, como muestra el Código~\ref{cod:flower-fedconfig}.
Con ello fue posible variar el número de supernodos o clientes y la fracción de
GPU asignada a cada uno en función del modo de ejecución, sin tocar el código de
la aplicación. El parámetro \texttt{init-args-num-cpus=2} no es cosmético: se
incorporó para resolver un agotamiento de memoria RAM del anfitrión, una
limitación que se detalla en la Sección~\ref{subsec:flex-flower-lim}.

\begin{bashcode}[title={run\_flower\_experiments.sh}, label={cod:flower-fedconfig}]
--federation-config="num-supernodes=$clients \
  init-args-num-cpus=2 \
  client-resources-num-cpus=1 \
  client-resources-num-gpus=$gpu_share"
\end{bashcode}

Una segunda adaptación afectó a la declaración de la estrategia. Flower exige que
las claves que se sobrescriben mediante \texttt{-{}-run-config} estén previamente
declaradas en la sección \texttt{[tool.flwr.app.config]} del fichero
\texttt{pyproject.toml}; de lo contrario, la ejecución aborta con el mensaje
\mintinline{text}{Key 'strategy' is not present in the main dictionary}. Por ese
motivo se añadió la clave \texttt{strategy} a la configuración de la aplicación,
junto con el resto de hiperparámetros del escenario (Código~\ref{cod:flower-pyproject}).
Asimismo, el valor de \texttt{learning-rate} se fijó en 0,02 para que coincidiera
con el que el script aplica de forma efectiva, eliminando una posible fuente de
confusión entre el valor declarado y el realmente utilizado.

\begin{codigo}[title={pyproject.toml}, label={cod:flower-pyproject}]{toml}
[tool.flwr.app.config]
num-server-rounds = 3
fraction-train = 0.5
fraction-evaluate = 0.0
local-epochs = 1
learning-rate = 0.02
batch-size = 32
strategy = "fedavg"
\end{codigo}

\subsection{Selección dinámica de la estrategia de agregación}\label{subsec:flex-flower-strategy}

El núcleo de la evaluación experimental de flexibilidad fue la sustitución de la
estrategia de agregación sin alterar el resto del sistema. Para ello se modificó
el servidor de Flower de manera que la estrategia se seleccionara dinámicamente a
partir de la configuración de ejecución, leyendo la clave \texttt{strategy} de
\mintinline{python}{context.run_config}. El servidor instancia entonces la clase
correspondiente y, en el caso de las estrategias adaptativas, fija los
hiperparámetros del optimizador del servidor antes de iniciar el entrenamiento,
como recoge el Código~\ref{cod:flower-strategy}. Gracias a este diseño, el mismo
código de cliente, modelo y conjunto de datos puede ejecutarse con distintas
estrategias modificando únicamente la lógica del servidor.

\begin{pythoncode}[title={server\_app.py}, label={cod:flower-strategy}]
from flwr.serverapp.strategy import FedAvg, FedAdagrad, FedAdam

strategy_name = str(context.run_config.get("strategy", "fedavg")).lower()

if strategy_name == "fedadam":
    strategy_class = FedAdam
    adaptive_kwargs = dict(
        eta=float(context.run_config.get("server-lr", 0.01)),
        tau=float(context.run_config.get("tau", 0.01)),
    )
elif strategy_name == "fedadagrad":
    strategy_class = FedAdagrad
    adaptive_kwargs = dict(
        eta=float(context.run_config.get("server-lr", 0.01)),
        tau=float(context.run_config.get("tau", 0.01)),
    )
else:
    strategy_class = FedAvg
    adaptive_kwargs = {}

strategy = strategy_class(**common_kwargs, **adaptive_kwargs)
\end{pythoncode}

Las estrategias FedAdam y FedAdagrad pertenecen a la familia de optimización
federada adaptativa propuesta por \textcite{reddi2021adaptive}, que generaliza
FedAvg sustituyendo el promedio del servidor por un optimizador con momento y tasa
de aprendizaje propios. Con los valores por defecto de Flower
(\texttt{eta}~=~0,1 y \texttt{tau}~=~0,001), ambas estrategias divergieron
numéricamente en este escenario: la pérdida del modelo global creció hasta el
orden de $10^{9}$ mientras las pérdidas locales de entrenamiento descendían con
normalidad. El problema se localiza en el paso de agregación del servidor, donde
una tasa de aprendizaje alta combinada con un factor de adaptabilidad pequeño
produce primeros pasos desproporcionados antes de que el segundo momento se
estabilice. Reducir la tasa del servidor a \texttt{eta}~=~0,01 y elevar el factor
a \texttt{tau}~=~0,01 estabilizó el entrenamiento. Lejos de restar valor a la
evaluación, este ajuste evidencia precisamente la flexibilidad buscada: Flower
expone esos hiperparámetros de agregación y permite configurarlos desde el
servidor.

\subsection{Mini-experimento de flexibilidad}\label{subsec:flex-flower-mini}

Para comprobar la sustitución de estrategias de forma controlada se definió una
variante reducida del experimento e3, denominada \texttt{e3\_flex\_mini}, cuyos
parámetros recoge la Tabla~\ref{tab:flex-flower-config}. La finalidad de esta
prueba no era comparar el rendimiento de aprendizaje, sino demostrar que Flower
permite cambiar la estrategia federada y la configuración de participación sin
rehacer la implementación, reduciendo a la vez el coste temporal de las pruebas.
Se ejecutó con una única semilla (42) y una repetición, por lo que las cifras de
aprendizaje son auxiliares y no llevan estimación de variabilidad.

\begin{table}[htbp]
  \centering
  \caption{Configuración del mini-experimento de flexibilidad de Flower.}
  \label{tab:flex-flower-config}
  \begin{tabular}{lc}
    \toprule
    \textbf{Parámetro} & \textbf{Valor} \\
    \midrule
    Clientes totales & 8 \\
    Clientes por ronda & 4 \\
    Fracción de entrenamiento & 0,5 \\
    Rondas federadas & 3 \\
    Épocas locales & 1 \\
    Tamaño de lote & 32 \\
    Tasa de aprendizaje (cliente) & 0,02 \\
    \texttt{eta} servidor (FedAdam/FedAdagrad) & 0,01 \\
    \texttt{tau} (FedAdam/FedAdagrad) & 0,01 \\
    Modo de ejecución & Secuencial \\
    Repeticiones / semilla & 1 / 42 \\
    \bottomrule
  \end{tabular}
\end{table}

Los tres mini-experimentos terminaron correctamente, lo que confirma que el mismo
flujo se ejecutó con FedAvg, FedAdam y FedAdagrad manteniendo constantes el
conjunto de datos, el modelo, el número de clientes, los clientes por ronda, las
rondas, las épocas locales, el tamaño de lote y la tasa de aprendizaje del
cliente. La Tabla~\ref{tab:flex-flower-resultados} resume las métricas obtenidas.

\begin{table}[htbp]
  \centering
  \caption{Resultados del mini-experimento de flexibilidad de Flower con tres
  estrategias de agregación. Las cifras son auxiliares y no deben interpretarse
  como una comparación de rendimiento. La memoria de GPU corresponde al pico
  reservado por PyTorch al proceso del cliente, métrica comparable entre
  estrategias.}
  \label{tab:flex-flower-resultados}
  \begin{tabular}{lcccrcr}
    \toprule
    \textbf{Estrategia} & \textbf{Estado} & \textbf{Acc.} & \textbf{Loss} &
    \textbf{Dur. (s)} & \textbf{GPU proc. (MiB)} & \textbf{Energía (Wh)} \\
    \midrule
    FedAvg     & success & 0,4257 & 1,5265 & 132 & $\approx$289 & 1,272 \\
    FedAdam    & success & 0,1410 & 2,3401 & 124 & $\approx$289 & 1,252 \\
    FedAdagrad & success & 0,1000 & 18,7832 & 144 & $\approx$289 & 1,123 \\
    \bottomrule
  \end{tabular}
\end{table}

Estos resultados no permiten concluir que una estrategia sea superior a otra, por
dos motivos. En primer lugar, se trata de un mini-experimento de solo tres rondas
y una época local, insuficiente para que las estrategias adaptativas estabilicen
su segundo momento y converjan; con los hiperparámetros ajustados ya no divergen,
pero todavía no alcanzan el comportamiento de FedAvg. En segundo lugar, se ejecutó
con una única semilla, sin medida de variabilidad. Conviene además una
advertencia de medición: la memoria de GPU reportada por \texttt{nvidia-smi} a
nivel de sistema mostró un valor anómalo de 1623~MiB en la ejecución de
FedAdagrad frente a 757~MiB en las otras dos, atribuible a memoria ocupada por
otras aplicaciones del equipo durante esa corrida; el consumo real por proceso fue
prácticamente idéntico ($\approx$289~MiB) en las tres estrategias, por lo que para
cualquier comparación debe emplearse la métrica por proceso. La conclusión
experimental válida es que Flower permitió sustituir la estrategia federada entre
FedAvg, FedAdam y FedAdagrad sin modificar el cliente, el modelo ni el conjunto de
datos, demostrando su flexibilidad para alterar la lógica de agregación desde el
servidor y la configuración de ejecución.

\subsection{Integración de modelo, datos y métricas}\label{subsec:flex-flower-integracion}

Más allá de las estrategias, la evaluación confirmó que Flower integra sin
fricción un modelo y un conjunto de datos propios. El modelo es la ResNet-18
adaptada a CIFAR-10, construida en PyTorch sustituyendo la primera convolución por
una de núcleo 3$\times$3 y eliminando el \textit{maxpool} inicial, mientras que el
conjunto de datos se carga mediante \mintinline{python}{FederatedDataset} con un
particionador IID. El bucle de entrenamiento local es código propio del cliente:
emplea descenso de gradiente estocástico con momento y \textit{weight decay},
entropía cruzada como función de pérdida y un número de épocas, tasa y tamaño de
lote configurables. Además, el cliente registra métricas personalizadas, entre
ellas la pérdida de entrenamiento, el número de ejemplos, la duración y el pico de
memoria de GPU asignada por PyTorch, que el script de ejecución vuelca junto con
los registros completos en ficheros de resumen reutilizables. Esta combinación de
modelo propio, bucle de entrenamiento editable y métricas a medida sustenta una
valoración alta en los criterios de personalización, integración y
reproducibilidad.

\subsection{Capacidades avanzadas según la documentación}\label{subsec:flex-flower-doc}

Más allá de lo ejecutado, la documentación oficial de Flower (versión 1.31) describe
un conjunto amplio de técnicas de flexibilidad que se valoran como evidencia
documental. En el plano de la configuración y el control del servidor, se detallan
cuatro vías de personalización de las estrategias: ajustar los parámetros de una
estrategia existente (\texttt{fraction\_train}, \texttt{fraction\_evaluate},
\texttt{min\_available\_nodes}); inyectar funciones de devolución como
\texttt{evaluate\_fn} que el servidor ejecuta en momentos concretos; sobrescribir
métodos internos como \texttt{configure\_train}, \texttt{aggregate\_fit} o
\texttt{aggregate\_evaluate}; e implementar una estrategia nueva heredando de
\texttt{BaseStrategy}~\parencite{beutel2020flower}. La misma documentación muestra cómo
enviar configuraciones por ronda mediante \texttt{ConfigRecord} —por ejemplo, reducir
progresivamente la tasa de aprendizaje leyendo \texttt{server\_round}— y cómo agregar
métricas con criterios propios a través de \texttt{evaluate\_metrics\_aggr\_fn}, en
lugar del promedio ponderado por número de ejemplos. Los clientes, sin estado por
defecto, pueden conservar información entre rondas mediante \texttt{Context.state}, lo
que habilita técnicas como el \textit{checkpointing} en el cliente.

Una pieza distintiva son los \textit{Mods}, funciones que interceptan los mensajes
entre servidor y cliente para transformarlos sin alterar la lógica principal; se
encadenan en un orden definido y permiten, por ejemplo, recortar gradientes
(\texttt{fixedclipping\_mod}, \texttt{adaptiveclipping\_mod}) o limitar el tamaño de
los mensajes (\texttt{message\_size\_mod})~\parencite{beutel2020flower}. Sobre esta
base se construye la privacidad diferencial, documentada en modo preliminar en tres
variantes —central del lado del servidor, central del lado del cliente y local
mediante \texttt{LocalDpMod}—, todas con recorte y ruido gaussiano configurables, a las
que se suma la agregación segura. La documentación también ejemplifica esquemas no
estándar, como FedBN —que excluye los parámetros de normalización por lotes del
intercambio para preservar las estadísticas locales de cada cliente—, el aprendizaje
federado vertical y la integración con XGBoost.

La documentación recoge asimismo soporte para el entrenamiento asíncrono, si bien de
carácter experimental y fuera de las guías principales. Las notas de versión describen
una \textit{Driver API} experimental (versión 1.2) que ejecuta un servidor programable,
asíncrono y multi-inquilino, y la API \textit{Flower Next} (versión 1.8), un conjunto de
primitivas de bajo nivel que permite personalizar el bucle de entrenamiento e incluye,
de forma explícita, soporte para aprendizaje federado asíncrono y esquemas cíclicos,
con un \texttt{ServerApp} capaz de permanecer activo de forma
indefinida~\parencite{beutel2020flower}. Por su naturaleza en desarrollo y de bajo
nivel, esta capacidad se valora como soporte parcial y no como una funcionalidad
directa y estable. Todas estas capacidades se reconocen como existentes, pero no se
contabilizan como demostradas empíricamente en este trabajo.

\subsection{Ejecución, despliegue y operación según la documentación}\label{subsec:flex-flower-deploy}

La flexibilidad de Flower se extiende a su modelo de ejecución y despliegue. En local,
los comandos \texttt{flwr run} y \texttt{flwr list} levantan automáticamente un
\textit{SuperLink} que gestiona las ejecuciones, mientras que las simulaciones se
lanzan sobre el \textit{Simulation Runtime}, cuyo número de \textit{SuperNodes} y
recursos por cliente se fijan de forma persistente o se sobrescriben por ejecución;
la documentación contempla incluso simulaciones sobre un clúster Ray
multi-anfitrión~\parencite{beutel2020flower}. Para despliegues distribuidos, la
federación se compone de un \textit{SuperLink} que coordina y varios
\textit{SuperNodes} que ejecutan los clientes, modelo que la documentación lleva hasta
entornos de producción en la nube —Google Cloud sobre Kubernetes, Microsoft Azure y
Red Hat OpenShift— e incluso a la interconexión de varios clústeres mediante Skupper.

Este recorrido se acompaña de mecanismos de seguridad operativa que también cuentan
para la flexibilidad: cifrado de las comunicaciones mediante TLS con certificados
propios, autenticación de \textit{SuperNodes} basada en firmas y lista blanca de
claves públicas, autenticación de cuentas de usuario mediante OpenID Connect con
autorización fina y un registro de auditoría en formato JSON que traza qué actor
ejecuta cada acción~\parencite{beutel2020flower}. Por último, la interfaz de línea de
comandos puede emitir su salida en formato JSON mediante la opción
\texttt{-{}-format json} para integrarse en flujos automatizados, y las herramientas de
gestión de federaciones permiten inspeccionar nodos, ejecuciones y estado. Ninguna de
estas capacidades se ejercitó en los experimentos, pero refuerzan la valoración de
Flower en los modos de ejecución y despliegue y en la integración con herramientas
externas.

\subsection{Limitaciones observadas}\label{subsec:flex-flower-lim}

El hallazgo más relevante de la evaluación afecta a la memoria RAM del anfitrión y
no a la GPU. El motor de simulación de Flower delega en Ray, que al iniciarse
detecta los procesadores del equipo (veinte en este caso) y prelanza un proceso
\textit{worker} ocioso por cada uno; cada worker carga el intérprete y las
bibliotecas (en torno a 0,28~GB), de modo que el conjunto de workers, sumado al
proceso de simulación y al resto del sistema, agotaba los aproximadamente 14,6~GB
de RAM disponibles. Ray respondía terminando procesos con un error de memoria
agotada, lo que ocurría incluso en modo secuencial, donde a priori solo hay un
cliente entrenando a la vez. La solución fue limitar el número de procesadores que
ve la \textit{Simulation Runtime} mediante \texttt{init-args-num-cpus=2}, lo que
reduce drásticamente la huella de RAM. Se trata de un resultado comparativo de
interés: el backend basado en Ray presenta una huella de memoria de anfitrión
notablemente mayor que el simulador de NVIDIA FLARE o el backend de proceso único
de FedML, que se ejecutan en un único proceso, por lo que en equipos con poca RAM
la simulación de Flower obliga a acotar los recursos manualmente.

La ejecución concurrente, probada también mediante Ray con una fracción de GPU por
cliente de 0,5, quedó condicionada por esta misma limitación de memoria, en línea
con lo descrito en el capítulo de resultados de rendimiento. Por su parte, la
divergencia de las estrategias adaptativas con los valores por defecto, ya
comentada, es una limitación de uso y no del framework: las estrategias de la
familia FedOpt son sensibles a la tasa de aprendizaje del servidor y requieren un
ajuste deliberado, especialmente con pocas rondas~\parencite{reddi2021adaptive}.
Ninguna de estas limitaciones impidió completar la evaluación, pero todas matizan
la valoración final.

\subsection{Valoración de la flexibilidad de Flower}\label{subsec:flex-flower-score}

La Tabla~\ref{tab:flex-flower-score} resume la valoración de Flower sobre los diez
criterios F1--F10, distinguiendo el tipo de evidencia que sustenta cada
puntuación. Los criterios de personalización del modelo y los datos (F1), del
entrenamiento local (F2), de la estrategia de agregación (F3) y de los algoritmos
federados disponibles (F4) se valoran con la máxima puntuación a partir de
evidencia experimental directa, al igual que los modos de ejecución y despliegue
(F6), la integración con herramientas externas (F7) y la reproducibilidad (F10).
El soporte de aprendizaje asíncrono (F5), la seguridad y privacidad (F8) y la
modificación de la arquitectura federada (F9) reciben puntuación parcial por estar
respaldados solo por documentación; en el caso de F5, mediante interfaces de bajo
nivel todavía en desarrollo (la \textit{Driver API} y \textit{Flower Next}), por lo
que se reconoce el soporte sin elevarlo a la máxima puntuación.

\begin{table}[htbp]
  \centering
  \caption{Valoración de la flexibilidad de Flower según los criterios F1--F10.
  El tipo de evidencia distingue las capacidades ejecutadas en los experimentos
  (experimental) de las descritas en la documentación oficial (documental).}
  \label{tab:flex-flower-score}
  \begin{tabular}{clcl}
    \toprule
    \textbf{Criterio} & \textbf{Descripción} & \textbf{Punt.} & \textbf{Evidencia} \\
    \midrule
    F1  & Personalización del modelo y el dataset      & 2/2 & Experimental \\
    F2  & Personalización del entrenamiento local      & 2/2 & Experimental \\
    F3  & Estrategias de agregación personalizadas     & 2/2 & Experimental \\
    F4  & Algoritmos federados disponibles             & 2/2 & Experimental \\
    F5  & Aprendizaje federado asíncrono               & 1/2 & Documental \\
    F6  & Modos de ejecución y despliegue              & 2/2 & Exp.\ + documental \\
    F7  & Integración con herramientas externas        & 2/2 & Exp.\ + documental \\
    F8  & Seguridad, privacidad y extensiones          & 1/2 & Documental \\
    F9  & Modificación de la arquitectura federada     & 1/2 & Documental \\
    F10 & Reproducibilidad y configuración             & 2/2 & Experimental \\
    \midrule
    \multicolumn{2}{r}{\textbf{Total}} & \multicolumn{2}{l}{\textbf{17/20}} \\
    \bottomrule
  \end{tabular}
\end{table}

En conjunto, Flower obtiene una valoración alta en flexibilidad (17/20). La
evidencia experimental muestra una herramienta cómoda de adaptar:
permitió integrar un modelo y un conjunto de datos propios en PyTorch, editar el
bucle de entrenamiento local, sustituir la estrategia de agregación entre FedAvg,
FedAdam y FedAdagrad desde el servidor —incluyendo el ajuste de sus
hiperparámetros— y automatizar la recogida de métricas. La valoración no alcanza
el máximo porque varias capacidades avanzadas —el aprendizaje asíncrono, la
privacidad diferencial, la agregación segura o los esquemas no horizontales— solo se
verificaron documentalmente, y porque el soporte asíncrono se apoya en interfaces
todavía experimentales; a ello se añade que la viabilidad práctica del modo
concurrente quedó condicionada por la elevada huella de memoria del backend Ray en
el equipo empleado. Estas conclusiones se contrastan con las de NVIDIA FLARE y FedML
en la comparativa final del capítulo.

% =====================================================================
%  SECCIÓN AÑADIDA: Flexibilidad de NVIDIA FLARE
% =====================================================================

\section{Flexibilidad de NVIDIA FLARE}\label{sec:flex-nvflare}

NVIDIA FLARE se presenta como un SDK de aprendizaje federado agnóstico al dominio
y compatible con flujos de trabajo de aprendizaje automático ya existentes, con
una orientación que abarca desde la simulación local hasta el despliegue en
producción y en el borde~\parencite{roth2022nvflare,nvflareWelcome}. La evaluación
experimental de su flexibilidad se apoya en un paquete de cinco ejecuciones de la
variante reducida del experimento e3, la misma \texttt{e3\_flex\_mini} empleada con
Flower (Subsección~\ref{subsec:flex-flower-mini}), con ocho clientes totales,
cuatro por ronda, tres rondas, una época local, tamaño de lote 32, tasa de
aprendizaje 0,02 y semilla 42, en modo secuencial con un único hilo del simulador.
Conviene precisar el alcance: el paquete inspeccionado contiene una selección de
suites concretas y no la totalidad del histórico de ejecuciones, por lo que las
conclusiones experimentales se refieren únicamente a los archivos incluidos en él.

\subsection{Diseño modular: jobs, componentes y workflows}\label{subsec:flex-nvflare-modular}

La flexibilidad de NVIDIA FLARE descansa en su modelo de configuración modular. La
unidad de ejecución es el \textit{job}, que describe por completo un experimento
mediante ficheros de configuración de servidor y de cliente
(\texttt{config\_fed\_server.conf} y \texttt{config\_fed\_client.conf}), en los que
cada elemento del sistema se declara como un componente con su ruta de clase y sus
argumentos de construcción~\parencite{nvflareComponents}. Al desplegar el job, el
servidor y los clientes analizan esos ficheros e instancian dinámicamente los
componentes declarados, registrándolos en el bucle de eventos del framework; cada
componente se identifica con un \texttt{component\_id} que permite referenciarlo
desde otros componentes o workflows~\parencite{nvflareComponents}. Este diseño es
el que permite sustituir un agregador, un persistor, un \textit{executor}, un
generador de objetos compartibles (\textit{Shareable}) o un filtro modificando la
configuración, sin reescribir la aplicación completa.

Sobre esa base de bajo nivel, la documentación ofrece además una capa de
\textit{recipes} que empaqueta algoritmos y workflows habituales —FedAvg, FedProx,
FedOpt, SCAFFOLD, aprendizaje cíclico, XGBoost, Sklearn, estadísticas federadas,
evaluación cruzada, PSI, integración con Flower y swarm learning, entre
otros~\parencite{nvflareRecipes}. Esta dualidad, configuración granular por
componentes y configuración de alto nivel por recetas, sustenta una valoración
elevada en personalización y reproducibilidad, y constituye el punto de partida de
las adaptaciones realizadas para probar estrategias alternativas.

\subsection{Adaptaciones para soportar las estrategias}\label{subsec:flex-nvflare-adapt}

El objetivo de la prueba fue comprobar si un mismo job podía adaptarse para
ejecutar FedAvg, FedProx y variantes FedOpt del servidor (FedAdagrad y FedAdam),
manteniendo constante la estructura general del experimento. Para ello se
introdujeron tres cambios de naturaleza distinta. En primer lugar, la selección de
estrategia se trasladó al cliente, que admite los valores \texttt{fedavg},
\texttt{fedprox}, \texttt{fedadam} y \texttt{fedadagrad} mediante un argumento de
línea de órdenes que, por defecto, toma su valor de la variable de entorno
\texttt{TFG\_FL\_STRATEGY}. En segundo lugar, FedProx se implementó en el
entrenamiento local: cuando la estrategia es \texttt{fedprox} y el coeficiente
proximal es positivo, el entrenamiento añade un término que penaliza la desviación
de los pesos locales respecto al modelo global recibido al inicio de la tarea,
manteniendo el descenso de gradiente estocástico como optimizador
local~\parencite{nvflareAlgorithms}.

El tercer cambio afecta al servidor y es el más relevante para la flexibilidad de
agregación. Para las variantes FedOpt, el script de ejecución parchea el fichero de
configuración del servidor sustituyendo el generador de compartibles por defecto,
\mintinline{text}{FullModelShareableGenerator}, por
\mintinline{text}{PTFedOptModelShareableGenerator}, y añade un componente de modelo
fuente requerido por este generador, como ilustra el
Código~\ref{cod:nvflare-fedopt}. Este generador actualiza el modelo global mediante
un optimizador de PyTorch —y, opcionalmente, un planificador de tasa de
aprendizaje— cuyos argumentos incluyen \texttt{optimizer\_args},
\texttt{lr\_scheduler\_args}, \texttt{source\_model} y
\texttt{device}~\parencite{nvflareFedopt}. En consecuencia, FedAdam y FedAdagrad no
son recetas independientes, sino configuraciones de FedOpt que se obtienen
seleccionando \mintinline{text}{torch.optim.Adam} o
\mintinline{text}{torch.optim.Adagrad} como optimizador del servidor.

\begin{codigo}[title={Fragmento representativo de config\_fed\_server.conf (FedOpt)}, label={cod:nvflare-fedopt}]{json}
"shareable_generator": {
  "path": "nvflare.app_opt.pt.fedopt.PTFedOptModelShareableGenerator",
  "args": {
    "source_model": "model",
    "optimizer_args": { "path": "torch.optim.Adagrad" }
  }
},
"components": [
  { "id": "model", "path": "<ruta de la ResNet-18 adaptada>" }
]
\end{codigo}

Una precaución metodológica relevante surgió con el dispositivo de ejecución. En
todos los ficheros de configuración del cliente aparece \texttt{-{}-device cuda},
pero en la ejecución de FedAdam en CPU el propio cliente seleccionó CPU porque
\mintinline{python}{torch.cuda.is_available()} devolvió falso al lanzarse con
\texttt{CUDA\_VISIBLE\_DEVICES} vacío. Por ello, el dispositivo realmente empleado
debe leerse del registro de ejecución y no del fichero de configuración, criterio
que se ha seguido al elaborar la Tabla~\ref{tab:flex-nvflare-resultados}.

\subsection{Mini-experimento de flexibilidad}\label{subsec:flex-nvflare-mini}

La Tabla~\ref{tab:flex-nvflare-resultados} resume las cinco ejecuciones. FedAvg
sirvió de referencia y se ejecutó correctamente con el workflow \textit{Scatter and
Gather}, en el que el servidor distribuye los pesos iniciales, los clientes
entrenan localmente y devuelven sus pesos, y el sistema los agrega por promedio
ponderado~\parencite{nvflareAlgorithms}. A partir de esa base, FedProx finalizó las
tres rondas y los registros confirmaron \texttt{strategy=fedprox} con un coeficiente
proximal \texttt{fedprox\_mu=0,001} activo durante el entrenamiento. FedAdagrad
también completó las tres rondas como FedOpt real —y no como mera etiqueta— porque
la configuración del servidor incorporaba el generador
\mintinline{text}{PTFedOptModelShareableGenerator} y el registro mostraba las
actualizaciones del servidor con \mintinline{text}{torch.optim.Adagrad}.

\begin{table}[htbp]
  \centering
  \caption{Resultados del mini-experimento de flexibilidad de NVIDIA FLARE. Las
  cifras de aprendizaje son auxiliares (tres rondas, una época, una semilla) y no
  constituyen una comparación de rendimiento. El dispositivo real procede del
  registro de ejecución, no del fichero de configuración del cliente.}
  \label{tab:flex-nvflare-resultados}
  \begin{tabular}{llccccrr}
    \toprule
    \textbf{Estrategia} & \textbf{Estado} & \textbf{Disp.} & \textbf{FedOpt} &
    \textbf{Rondas} & \textbf{Dur. (s)} & \textbf{Acc.} & \textbf{Loss} \\
    \midrule
    FedAvg      & success & \texttt{cuda:0} & No  & 3 & 205  & 0,4178 & 1,5601 \\
    FedProx     & success & \texttt{cuda:0} & No  & 3 & 225  & 0,4045 & 1,5876 \\
    FedAdagrad  & success & \texttt{cuda:0} & Sí  & 3 & 199  & 0,1402 & 2,3785 \\
    FedAdam     & failed  & \texttt{cuda:0} & Sí  & 0 & 77   & ---    & ---    \\
    FedAdam     & success & \texttt{cpu}    & Sí  & 3 & 1280 & 0,2586 & 3,4439 \\
    \bottomrule
  \end{tabular}
\end{table}

El caso de FedAdam exige una lectura cuidadosa. La estrategia llegó a configurarse
realmente mediante FedOpt con \mintinline{text}{torch.optim.Adam}, pero su ejecución
en GPU abortó durante el paso de actualización del servidor
(\mintinline{python}{torch.optim.Adam.step()}) con un error de inicialización de
CUDA, registrado como \texttt{FATAL\_SYSTEM\_ERROR}. Para aislar el problema, la
prueba se repitió en CPU, donde FedAdam completó las tres rondas; esta ejecución
valida FedAdam como estrategia configurable, pero su duración (1280~s) es muy
superior precisamente por entrenarse en CPU, de modo que no debe mezclarse con la
comparativa de rendimiento en GPU del capítulo correspondiente. En coherencia con
ello, la fila de la ejecución fallida no reporta precisión ni pérdida, ya que no
completó ninguna ronda.

Estos resultados no permiten ordenar las estrategias por calidad, por los mismos
motivos señalados para Flower: tres rondas, una época local y una sola semilla son
insuficientes para que las variantes adaptativas converjan. Su valor es
demostrar que NVIDIA FLARE permite adaptar tanto el entrenamiento local como la
lógica de actualización del servidor mediante configuración y componentes. Importa
además una cautela de interpretación: la validez de FedAdagrad y FedAdam no se
sustenta en el nombre de la ejecución, sino en la presencia del generador FedOpt y
del optimizador del servidor en los registros. Del mismo modo, el valor
\texttt{fedprox\_mu=0,001} que aparece en la configuración no implica que FedProx se
aplique a todas las estrategias: en FedAdagrad y FedAdam el entrenamiento imprime
\texttt{fedprox\_mu=0,0}, porque el término proximal solo se activa cuando la
estrategia es FedProx.

De esta experiencia se desprende una observación comparativa de interés. Frente a
Flower, donde la estrategia se sustituía cambiando una única clase en el servidor,
NVIDIA FLARE exigió modificar componentes del servidor, añadir un modelo fuente para
FedOpt y comprobar en los registros que la estrategia configurada era realmente la
que se ejecutaba. La flexibilidad de NVIDIA FLARE es, por tanto, amplia, pero su
ejercicio práctico conlleva una complejidad de configuración mayor.

\subsection{Capacidades avanzadas según la documentación}\label{subsec:flex-nvflare-doc}

Como en el caso de Flower, varias capacidades relevantes no se ejecutaron en los
experimentos y se valoran únicamente como evidencia documental. La documentación
oficial las articula en torno a las \textit{Job Recipes}, una API de alto nivel que
encapsula un trabajo federado exponiendo solo los parámetros clave —\texttt{min\_clients},
\texttt{num\_rounds}, modelo y \textit{script} de entrenamiento— y ocultando la
complejidad de controladores y \textit{workflows}~\parencite{nvflareRecipes}. Estas
recetas se ejecutan sobre tres entornos intercambiables: \texttt{SimEnv} simula los
clientes en procesos locales, \texttt{PocEnv} los lanza como procesos separados en la
misma máquina y \texttt{ProdEnv} despliega el trabajo en una instalación de
producción; un mismo trabajo puede así trasladarse de la simulación al despliegue sin
cambiar el código~\parencite{nvflareRecipes}. El catálogo de recetas cubre FedAvg,
FedProx, FedOpt —donde el optimizador del servidor se fija con un argumento que indica
su ruta de clase e hiperparámetros, por ejemplo \texttt{torch.optim.SGD} con tasa y
momento—, SCAFFOLD y el aprendizaje cíclico, además de XGBoost federado en sus
variantes horizontal, \textit{bagging} y vertical~\parencite{nvflareAlgorithms,nvflareFedopt}.

En cuanto al aprendizaje asíncrono, la receta \texttt{EdgeFedBuffRecipe}, orientada a
dispositivos de borde, admite entrenamiento síncrono y asíncrono: en modo síncrono el
servidor espera todas las actualizaciones antes de generar un nuevo modelo, mientras
que en modo asíncrono lo genera en cuanto recibe la primera, lo que reduce el tiempo
total cuando los dispositivos presentan latencias heterogéneas~\parencite{nvflareEdge,nvflareWelcome}.
La receta expone parámetros como el número máximo de versiones de modelo activas o el
tamaño de selección de dispositivos por ronda. No obstante, la propia documentación
advierte de que la asincronía general aún no está disponible para todas las recetas y
se concentra en el escenario de borde, por lo que su valoración se mantiene como
soporte documental parcial.

En materia de seguridad y privacidad, la documentación es notablemente extensa.
NVIDIA FLARE implementa la privacidad mediante un mecanismo general de filtros que
procesan los objetos compartibles antes o después de la comunicación, configurables en
la receta como \texttt{task\_data\_filters} o \texttt{task\_result\_filters}. Soporta
privacidad diferencial a nivel de muestra mediante la integración de DP-SGD con Opacus
—controlada por los parámetros \texttt{epsilon} y \texttt{delta}— y filtros de DP local
como \texttt{PercentilePrivacy}, que comparte solo las diferencias de pesos por encima
de un percentil, \texttt{SVTPrivacy}, basado en la técnica del vector disperso con ruido
de Laplace, y \texttt{StatisticsPrivacyFilter} para estadísticas
federadas~\parencite{nvflareDP}. A un nivel superior, la solución de computación
confidencial ejecuta la agregación dentro de un entorno de ejecución confiable
(\textit{Trusted Execution Environment}) sobre hardware como AMD SEV-SNP o Intel TDX,
protege el modelo y el código del cliente en máquinas virtuales confidenciales y
permite combinar privacidad diferencial con cifrado homomórfico o agregación
segura~\parencite{nvflareSecurity,nvflareConfidential}.

Por último, NVIDIA FLARE documenta esquemas descentralizados mediante los
\textit{workflows} controlados por clientes, construidos sobre las clases base
\texttt{ServerSideController} y \texttt{ClientSideController} para escenarios donde el
servidor no es plenamente confiable~\parencite{nvflareSwarm}. Entre ellos destacan el
aprendizaje cíclico, en el que cada cliente entrena con el modelo recibido del anterior
según un orden fijo o aleatorio~\parencite{nvflareCyclic}; el aprendizaje \textit{swarm},
que designa un cliente agregador distinto por ronda y mejora la resiliencia frente a
fallos del servidor~\parencite{nvflareSwarm}; y la evaluación cruzada entre sitios, que
permite a los clientes validar modelos de otros mediante comunicación entre pares sin
exponerlos al servidor. A esto se suman el aprendizaje vertical con XGBoost y las
estadísticas federadas, así como las utilidades para integrar el seguimiento de
experimentos con TensorBoard, MLflow o Weights \& Biases sobre cualquier
receta~\parencite{nvflareTracking,nvflareRecipes}. La documentación reconoce, además,
dos cautelas: la API de recetas se encuentra en estado de \textit{technical preview} y
no cubre todavía todos los algoritmos, y la configuración de la computación confidencial
exige infraestructura especializada.

\subsection{Limitaciones observadas}\label{subsec:flex-nvflare-lim}

La limitación más concreta de la evaluación fue el fallo de FedAdam en GPU por un
error de inicialización de CUDA durante la actualización FedOpt del servidor. Debe
interpretarse como una incidencia técnica del entorno, no como una carencia del
framework, puesto que la misma configuración completó las tres rondas en CPU; aun
así, impide presentar FedAdam como ejecución correcta en GPU. La segunda limitación
es la complejidad práctica ya señalada: a diferencia de Flower, alterar la lógica de
agregación obligó a intervenir en componentes del servidor y a validar
cuidadosamente los registros, lo que incrementa la probabilidad de errores de
configuración. Por último, el modo concurrente quedó condicionado por los recursos
del equipo, en línea con lo expuesto en el capítulo de resultados de rendimiento;
aunque la documentación describe un sistema de ejecución multi-job con gestión de
recursos, esa capacidad no pudo validarse de forma estable en el hardware empleado.

\subsection{Valoración de la flexibilidad de NVIDIA FLARE}\label{subsec:flex-nvflare-score}

La Tabla~\ref{tab:flex-nvflare-score} recoge la valoración sobre los criterios
F1--F10, con el mismo convenio aplicado a Flower: puntuación máxima cuando existe
evidencia experimental directa, puntuación parcial cuando la capacidad solo consta
en la documentación y puntuación nula cuando no se identifica soporte. NVIDIA FLARE
obtiene la máxima puntuación en personalización del modelo y los datos (F1), del
entrenamiento local (F2), de la estrategia de agregación (F3) y de los algoritmos
federados disponibles (F4), todos respaldados por las ejecuciones, así como en los
modos de ejecución y despliegue (F6), la integración con herramientas externas (F7)
y la reproducibilidad (F10). La diferencia respecto a Flower aparece en el
aprendizaje asíncrono (F5): mientras que en Flower no se identificó soporte, NVIDIA
FLARE lo documenta de forma explícita mediante FedBuff, lo que justifica una
puntuación parcial. La seguridad y privacidad (F8) y la modificación de la
arquitectura federada (F9) reciben también puntuación parcial por estar respaldadas
solo documentalmente.

\begin{table}[htbp]
  \centering
  \caption{Valoración de la flexibilidad de NVIDIA FLARE según los criterios
  F1--F10. El tipo de evidencia distingue las capacidades ejecutadas en los
  experimentos de las descritas en la documentación oficial.}
  \label{tab:flex-nvflare-score}
  \begin{tabular}{clcl}
    \toprule
    \textbf{Criterio} & \textbf{Descripción} & \textbf{Punt.} & \textbf{Evidencia} \\
    \midrule
    F1  & Personalización del modelo y el dataset      & 2/2 & Exp.\ + documental \\
    F2  & Personalización del entrenamiento local      & 2/2 & Experimental \\
    F3  & Estrategias de agregación personalizadas     & 2/2 & Experimental \\
    F4  & Algoritmos federados disponibles             & 2/2 & Exp.\ + documental \\
    F5  & Aprendizaje federado asíncrono               & 1/2 & Documental \\
    F6  & Modos de ejecución y despliegue              & 2/2 & Exp.\ + documental \\
    F7  & Integración con herramientas externas        & 2/2 & Exp.\ + documental \\
    F8  & Seguridad, privacidad y extensiones          & 1/2 & Documental \\
    F9  & Modificación de la arquitectura federada     & 1/2 & Documental \\
    F10 & Reproducibilidad y configuración             & 2/2 & Exp.\ + documental \\
    \midrule
    \multicolumn{2}{r}{\textbf{Total}} & \multicolumn{2}{l}{\textbf{17/20}} \\
    \bottomrule
  \end{tabular}
\end{table}

En conjunto, NVIDIA FLARE alcanza una valoración de flexibilidad de 17/20. La
evidencia experimental
demuestra que permite adaptar el entrenamiento local —incluido el término proximal
de FedProx— y, sobre todo, la lógica de agregación del servidor mediante FedOpt, con
FedAvg, FedProx y FedAdagrad ejecutados correctamente y FedAdam validado en CPU. A
ello se suma una de las superficies de capacidades documentadas más amplias del
estudio, especialmente sólida en seguridad y privacidad, que abarca aprendizaje
asíncrono, workflows cíclicos y \textit{swarm}, privacidad
diferencial, cifrado homomórfico y computación confidencial. El contrapunto es una
complejidad de configuración superior a la de Flower y el hecho de que esas
capacidades avanzadas solo se verificaron documentalmente; además, la ejecución de
FedAdam en GPU falló por un error de CUDA y el modo concurrente quedó limitado por el
hardware. Esta lectura se contrasta con la de los otros frameworks en la comparativa
final del capítulo.

% =====================================================================
%  SECCIÓN AÑADIDA: Flexibilidad de FedML
% =====================================================================

\section{Flexibilidad de FedML}\label{sec:flex-fedml}

FedML es una biblioteca de investigación y banco de pruebas para aprendizaje
federado~\parencite{he2020fedml}, documentada actualmente dentro de la plataforma
TensorOpera aunque la librería y los ejemplos siguen empleando el paquete
\texttt{fedml}. Su documentación organiza el framework en varios escenarios
—simulación, cross-silo y cross-device— y ofrece un catálogo amplio de algoritmos
de referencia. La evaluación de su flexibilidad combina la implementación propia
realizada para la comparativa, que extiende las abstracciones de FedML, con una
prueba reducida de sustitución de estrategias. Conviene una precisión de alcance:
por el coste computacional, la prueba reducida de FedML se ejecutó con cuatro
clientes totales, dos por ronda y dos rondas, una configuración algo menor que la
empleada con Flower y NVIDIA FLARE; las métricas de aprendizaje son, por tanto, una
comprobación funcional y no una comparación de rendimiento.

\subsection{Configuración y componentes personalizados}\label{subsec:flex-fedml-config}

La configuración de FedML se apoya en un fichero \texttt{fedml\_config.yaml}
estructurado en secciones (\texttt{common\_args}, \texttt{data\_args},
\texttt{model\_args}, \texttt{train\_args}, \texttt{device\_args},
\texttt{comm\_args} y \texttt{tracking\_args}), donde se declaran los parámetros
federados como \texttt{federated\_optimizer}, \texttt{client\_num\_in\_total},
\texttt{client\_num\_per\_round}, \texttt{comm\_round}, \texttt{epochs},
\texttt{batch\_size}, \texttt{client\_optimizer} y \texttt{learning\_rate}~\parencite{fedmlSimApi}.
Esta organización permitió variar clientes, rondas, épocas e hiperparámetros entre
ejecuciones sin reescribir la aplicación.

Más allá del fichero de configuración, la flexibilidad de FedML en este trabajo se
sustenta en evidencia experimental directa: la implementación de la comparativa
define un entrenador y un agregador propios que extienden las abstracciones
\mintinline{python}{ClientTrainer} y \mintinline{python}{ServerAggregator} del
framework. El entrenador \texttt{FedMLCifar10Trainer} encapsula el bucle de
entrenamiento local sobre la ResNet-18 adaptada —descenso de gradiente estocástico,
entropía cruzada y métricas propias—, mientras que el agregador
\texttt{FedMLCifar10Aggregator} gobierna la evaluación del modelo global; ambos se
conectan al flujo mediante \texttt{FedMLRunner}. La documentación respalda esta vía,
al indicar que el usuario puede definir modelo, datos, optimizadores, funciones de
pérdida y lógica de entrenamiento arbitrarios, así como dataloaders propios que
respeten la estructura esperada por el simulador~\parencite{fedmlTrainer,fedmlDataloader}.
Esta capacidad de sustituir componentes centrales mediante herencia sostiene una
valoración alta en personalización del entrenamiento local, de la agregación y de la
integración.

\subsection{Sustitución de la estrategia federada}\label{subsec:flex-fedml-strategy}

La prueba de flexibilidad consistió en ejecutar cuatro estrategias modificando
únicamente la configuración, manteniendo constantes el conjunto de datos, el
modelo y el resto de hiperparámetros. FedAvg y FedProx se seleccionan asignando
\texttt{federated\_optimizer} a \texttt{"FedAvg"} o \texttt{"FedProx"}, mientras que
FedAdam y FedAdagrad se obtienen mediante \texttt{federated\_optimizer:\ "FedOpt"}
combinado con un optimizador de servidor \texttt{Adam} o \texttt{Adagrad}, tal como
ilustra el Código~\ref{cod:fedml-config}. Durante las primeras pruebas, las
variantes FedOpt fallaban porque FedML esperaba el atributo \texttt{ci} en los
argumentos de configuración; tras añadir \texttt{ci:\ 0} a la configuración común,
ambas pudieron completar la ejecución. La selección del algoritmo desde la
configuración está respaldada por la documentación, que controla el optimizador
federado a través del mismo campo~\parencite{fedmlSimApi}.

\begin{yamlcode}[title={Fragmento de fedml\_config.yaml (FedAdam vía FedOpt)}, label={cod:fedml-config}]
train_args:
  federated_optimizer: "FedOpt"
  server_optimizer: "Adam"
  ci: 0
  client_num_in_total: 4
  client_num_per_round: 2
  comm_round: 2
  epochs: 1
  batch_size: 32
  learning_rate: 0.02
\end{yamlcode}

\subsection{Mini-experimento de flexibilidad}\label{subsec:flex-fedml-mini}

La Tabla~\ref{tab:flex-fedml-resultados} resume las cuatro ejecuciones, todas
finalizadas con estado \texttt{success}. FedAvg y FedProx concluyeron con métricas
coherentes para una prueba de dos rondas, con precisiones del 43,01\,\% y el
42,60\,\% respectivamente, lo que confirma el soporte experimental de una estrategia
alternativa a FedAvg con comportamiento correcto. FedAdam y FedAdagrad también
completaron la ejecución, lo que demuestra que fue posible activar FedOpt con
optimizadores de servidor Adam y Adagrad; sin embargo, sus métricas finales fueron
claramente deficientes, con una precisión próxima al 10\,\% —equivalente al azar en
un problema de diez clases— y pérdidas muy elevadas.

\begin{table}[htbp]
  \centering
  \caption{Resultados del mini-experimento de flexibilidad de FedML. Las cifras de
  aprendizaje son auxiliares (dos rondas, una época, una semilla) y constituyen una
  comprobación funcional, no una comparación de rendimiento entre estrategias.}
  \label{tab:flex-fedml-resultados}
  \begin{tabular}{llrrrrr}
    \toprule
    \textbf{Estrategia} & \textbf{Estado} & \textbf{Dur. (s)} & \textbf{Acc.} &
    \textbf{Loss} & \textbf{GPU (MiB)} & \textbf{Util. (\%)} \\
    \midrule
    FedAvg     & success & 118 & 0,4301 & 1,5248   & 416 & 75,29 \\
    FedProx    & success & 123 & 0,4260 & 1,5137   & 484 & 79,01 \\
    FedAdam    & success & 117 & 0,0998 & 809,7697 & 416 & 76,09 \\
    FedAdagrad & success & 116 & 0,0998 & 191,2948 & 416 & 76,85 \\
    \bottomrule
  \end{tabular}
\end{table}

La lectura correcta de estos datos es que FedML permite cambiar la estrategia
federada mediante configuración, pero que la mera disponibilidad de una estrategia
no garantiza su estabilidad ni un buen rendimiento. Las variantes FedOpt
ejecutaron de principio a fin, de modo que la flexibilidad queda demostrada, pero su
falta de aprendizaje en esta prueba indica que requieren un ajuste específico de
hiperparámetros —como la tasa de aprendizaje del servidor o un mayor número de
rondas— para resultar competitivas. Este patrón es coherente con lo observado en los
otros frameworks: en Flower las estrategias adaptativas divergían con los valores por
defecto y exigieron ajustar \texttt{eta} y \texttt{tau}, y en NVIDIA FLARE FedAdam no
llegó a completarse en GPU. La fragilidad de las variantes FedOpt con poca
configuración es, por tanto, un comportamiento transversal y no una carencia
particular de FedML.

\subsection{Capacidades avanzadas según la documentación}\label{subsec:flex-fedml-doc}

La documentación de FedML respalda una superficie de capacidades más amplia que la
probada experimentalmente, organizada en varias capas. En el plano de uso, ofrece una
interfaz de línea de comandos —integrada en la plataforma TensorOpera— con órdenes para
autenticación y gestión de dispositivos (\texttt{fedml login}, \texttt{fedml device
bind}), lanzamiento y gestión de trabajos y clústeres (\texttt{fedml launch},
\texttt{fedml cluster}, \texttt{fedml run}), y gestión y despliegue de modelos como
servicios de inferencia (\texttt{fedml model})~\parencite{fedmlExamples}. En el plano de
programación, expone una API de Python con una llamada de una línea
(\texttt{fedml.run\_simulation()}) y una API más explícita de varias líneas que
inicializa dispositivo, datos y modelo y ejecuta la federación mediante
\texttt{SimulatorMPI} o las clases \texttt{Client}/\texttt{Server} envueltas por
\texttt{FedMLRunner}~\parencite{fedmlSimApi}.

En cuanto al catálogo algorítmico, el simulador Parrot enumera implementaciones de
referencia como FedOpt, FedNova, FedGKT, aprendizaje federado descentralizado, vertical
y jerárquico, FedNAS y \textit{Split Learning}, seleccionables cambiando el campo
\texttt{federated\_optimizer}~\parencite{fedmlSimOverview,fedmlSimApi}, y el módulo de
seguridad declara compatibilidad con optimizadores como FedAvg, FedSGD, FedOpt, FedProx,
FedGKT y FedNova~\parencite{fedmlSecurity}. A este catálogo se añaden la analítica
federada, que calcula estadísticas distribuidas como promedios o percentiles mediante un
ejecutor específico sin compartir datos crudos, y un proyecto independiente orientado al
ajuste federado de grandes modelos de lenguaje, que reutiliza las mismas API y la
plataforma de operaciones~\parencite{fedmlExamples}. En cuanto al aprendizaje asíncrono,
FedML Octopus documenta un escenario jerárquico heterogéneo en el que cada silo puede
ejecutar entrenamiento distribuido con varias GPU y orquestarse mediante optimización
federada asíncrona o síncrona~\parencite{fedmlCrossSilo}; se trata, no obstante, de un
soporte menos específico que el de NVIDIA FLARE, ya que no se documenta una receta directa
equivalente a FedBuff, aprendizaje cíclico o \textit{swarm}.

En privacidad y seguridad, el perfil de FedML difiere del de NVIDIA FLARE. Su
fortaleza documental está en la simulación de ataques y defensas mediante
\texttt{FedMLAttacker} y \texttt{FedMLDefender}, con ataques de tipo bizantino,
\textit{label flipping} y reemplazo de modelo, y defensas como recorte de norma, Krum,
mediana geométrica o FoolsGold, activables desde el fichero de configuración con
\texttt{enable\_attack} y \texttt{enable\_defense}~\parencite{fedmlSecurity}, además de un
ejemplo de agregación segura basada en computación multiparte~\parencite{fedmlExamples}.
En cambio, en las páginas revisadas no figura una integración de privacidad diferencial o
cifrado homomórfico comparable a la de NVIDIA FLARE, por lo que su valoración en este
criterio se sostiene en el \textit{benchmarking} de robustez más que en la privacidad
criptográfica. Por último, para integración y trazabilidad, FedML documenta el registro
de métricas, artefactos y modelos mediante \texttt{fedml.log()}, la captura automática de
métricas de hardware y GPU, los backends de comunicación MPI, gRPC, PyTorch RPC y MQTT+S3,
y la integración con la plataforma de operaciones de TensorOpera para gestionar agentes,
experimentos y despliegues~\parencite{fedmlTracking,fedmlExamples}.

\subsection{Limitaciones observadas}\label{subsec:flex-fedml-lim}

La limitación más visible es que las variantes FedOpt finalizaron sin aprender bajo
la configuración empleada, lo que confirma que su soporte es real pero condicionado a
un ajuste de hiperparámetros que queda fuera del alcance de esta prueba de
flexibilidad. A ello se añade el detalle de configuración del atributo \texttt{ci},
que tuvo que fijarse manualmente para que FedOpt funcionara, y el hecho de que la
prueba reducida de FedML emplea una configuración menor que la de los otros dos
frameworks, por lo que sus métricas auxiliares no son directamente comparables. En
cuanto a la ejecución, el modo concurrente de FedML mediante el backend MPI fue el
único que llegó a completarse entre los tres frameworks, según se documenta en el
capítulo de resultados de rendimiento, si bien se mostró sensible al consumo de
memoria en las configuraciones más grandes.

\subsection{Valoración de la flexibilidad de FedML}\label{subsec:flex-fedml-score}

La Tabla~\ref{tab:flex-fedml-score} recoge la valoración de FedML sobre los criterios
F1--F10, con el convenio ya aplicado a los otros frameworks. FedML obtiene la máxima
puntuación en personalización del modelo y los datos (F1), del entrenamiento local
(F2) y de la agregación (F3), respaldadas por la implementación de entrenador y
agregador propios; en los algoritmos federados disponibles (F4), por haber ejecutado
cuatro estrategias mediante configuración; en los modos de ejecución (F6), por ser el
único framework cuyo modo concurrente se completó; y en integración (F7) y
reproducibilidad (F10). El aprendizaje asíncrono (F5), la seguridad y privacidad (F8)
y la modificación de la arquitectura federada (F9) reciben puntuación parcial por
estar respaldados únicamente por la documentación.

\begin{table}[htbp]
  \centering
  \caption{Valoración de la flexibilidad de FedML según los criterios F1--F10. El
  tipo de evidencia distingue las capacidades ejecutadas en los experimentos de las
  descritas en la documentación oficial.}
  \label{tab:flex-fedml-score}
  \begin{tabular}{clcl}
    \toprule
    \textbf{Criterio} & \textbf{Descripción} & \textbf{Punt.} & \textbf{Evidencia} \\
    \midrule
    F1  & Personalización del modelo y el dataset      & 2/2 & Exp.\ + documental \\
    F2  & Personalización del entrenamiento local      & 2/2 & Exp.\ + documental \\
    F3  & Estrategias de agregación personalizadas     & 2/2 & Experimental \\
    F4  & Algoritmos federados disponibles             & 2/2 & Exp.\ + documental \\
    F5  & Aprendizaje federado asíncrono               & 1/2 & Documental \\
    F6  & Modos de ejecución y despliegue              & 2/2 & Exp.\ + documental \\
    F7  & Integración con herramientas externas        & 2/2 & Exp.\ + documental \\
    F8  & Seguridad, privacidad y extensiones          & 1/2 & Documental \\
    F9  & Modificación de la arquitectura federada     & 1/2 & Documental \\
    F10 & Reproducibilidad y configuración             & 2/2 & Exp.\ + documental \\
    \midrule
    \multicolumn{2}{r}{\textbf{Total}} & \multicolumn{2}{l}{\textbf{17/20}} \\
    \bottomrule
  \end{tabular}
\end{table}

En conjunto, FedML alcanza una valoración de flexibilidad de 17/20. Su mayor fortaleza
experimental reside en la facilidad para sustituir componentes centrales del sistema
—entrenador y agregador— mediante herencia, y en haber ejecutado el abanico más amplio
de estrategias del estudio a través de la configuración. Su perfil documental es
asimismo notable, con un catálogo extenso de algoritmos y escenarios y un soporte de
seguridad orientado al benchmarking de ataques y defensas. Las cautelas son dos: las
variantes FedOpt completaron la ejecución pero no aprendieron sin un ajuste
adicional, y algunas capacidades avanzadas —asincronía, privacidad criptográfica y
esquemas alternativos— solo se verificaron documentalmente.

% =====================================================================
%  SECCIÓN AÑADIDA: Comparativa de flexibilidad entre frameworks
% =====================================================================

\section{Comparativa de flexibilidad}\label{sec:flex-comparativa}

Una vez evaluados los tres frameworks por separado, esta sección los contrasta sobre
el mismo conjunto de criterios. La Tabla~\ref{tab:flex-comparativa} reúne las
puntuaciones F1--F10 de Flower, NVIDIA FLARE y FedML. Los tres obtienen la misma
valoración, 17/20, lo que confirma que las tres herramientas son, en lo esencial,
comparablemente flexibles para el tipo de adaptaciones que exige este trabajo y que la
diferenciación entre ellas debe buscarse, más que en la cifra global, en el perfil
cualitativo de cada una.

\begin{table}[htbp]
  \centering
  \caption{Comparativa de la valoración de flexibilidad de los tres frameworks según
  los criterios F1--F10 definidos en la metodología.}
  \label{tab:flex-comparativa}
  \begin{tabular}{clccc}
    \toprule
    \textbf{Criterio} & \textbf{Descripción} & \textbf{Flower} & \textbf{NVIDIA FLARE} & \textbf{FedML} \\
    \midrule
    F1  & Modelo y dataset                  & 2/2 & 2/2 & 2/2 \\
    F2  & Entrenamiento local               & 2/2 & 2/2 & 2/2 \\
    F3  & Agregación personalizada          & 2/2 & 2/2 & 2/2 \\
    F4  & Algoritmos disponibles            & 2/2 & 2/2 & 2/2 \\
    F5  & Aprendizaje asíncrono             & 1/2 & 1/2 & 1/2 \\
    F6  & Modos de ejecución y despliegue   & 2/2 & 2/2 & 2/2 \\
    F7  & Integración externa               & 2/2 & 2/2 & 2/2 \\
    F8  & Seguridad y privacidad            & 1/2 & 1/2 & 1/2 \\
    F9  & Arquitectura federada             & 1/2 & 1/2 & 1/2 \\
    F10 & Reproducibilidad y configuración  & 2/2 & 2/2 & 2/2 \\
    \midrule
    \multicolumn{2}{r}{\textbf{Total}} & \textbf{17/20} & \textbf{17/20} & \textbf{17/20} \\
    \bottomrule
  \end{tabular}
\end{table}

La coincidencia en los criterios experimentales (F1--F4, F6, F7 y F10) refleja un
hecho de fondo: bajo un mismo escenario de modelo, datos y algoritmo base, los tres
frameworks permitieron integrar la ResNet-18 sobre CIFAR-10, editar el entrenamiento
local, sustituir la estrategia de agregación, registrar métricas y reproducir los
experimentos mediante ficheros y scripts. La igualdad también alcanza a los tres
criterios documentales: ninguno validó experimentalmente el aprendizaje asíncrono (F5),
la seguridad avanzada (F8) ni esquemas alternativos a la arquitectura clásica (F9), pero
los tres los documentan, por lo que reciben idéntica puntuación parcial. En F5, en
concreto, los tres declaran soporte —Flower mediante la \textit{Driver API} y
\textit{Flower Next}, NVIDIA FLARE mediante FedBuff en el borde y FedML en su escenario
jerárquico cross-silo—, si bien en todos los casos se trata de capacidades
experimentales, limitadas a entornos concretos o de bajo nivel.

Más allá de las cifras, el estudio revela un comportamiento transversal en las
estrategias adaptativas de la familia FedOpt. En los tres frameworks fue posible
configurar FedAdam y FedAdagrad, pero su ejecución resultó frágil con poca
configuración: en Flower divergían con los valores por defecto y necesitaron ajustar
la tasa del servidor y el factor de adaptabilidad; en NVIDIA FLARE, FedAdam no llegó a
completarse en GPU por un error de CUDA; y en FedML, FedAdam y FedAdagrad finalizaron
pero sin aprender. La conclusión es clara y común: la disponibilidad de una estrategia
no equivale a su estabilidad, que depende de un ajuste deliberado de hiperparámetros,
especialmente con pocas rondas.

Por último, cada framework muestra un perfil de flexibilidad propio que la puntuación,
por su granularidad, no captura del todo. Flower es el más directo para sustituir la
estrategia de agregación —basta con cambiar una clase en el servidor— y su documentación
revela una superficie amplia que abarca \textit{Mods}, privacidad diferencial, FedBN y,
sobre todo, un modelo de despliegue muy desarrollado hacia la nube (Google Cloud, Azure,
OpenShift e interconexión de clústeres); su contrapartida es que el soporte asíncrono
sigue siendo experimental y que el backend Ray arrastra una huella de memoria elevada.
NVIDIA FLARE ofrece la mayor profundidad en seguridad y privacidad —filtros de privacidad
diferencial, cifrado homomórfico y computación confidencial sobre entornos de ejecución
confiables— junto con un catálogo de recetas y entornos (\texttt{SimEnv}, \texttt{PocEnv},
\texttt{ProdEnv}) orientado a producción, a costa de una complejidad de configuración
mayor, pues alterar la agregación exige intervenir en componentes del servidor. FedML
destaca por la sustitución de entrenador y agregador mediante herencia, por haber
ejecutado el mayor número de estrategias, por ser el único cuyo modo concurrente se
completó, y por una integración estrecha con su plataforma de operaciones y la analítica
federada, con un soporte de seguridad orientado al \textit{benchmarking} de ataques y
defensas. La elección entre ellos, en términos de flexibilidad, depende por tanto de la
prioridad del proyecto: la simplicidad y el despliegue en la nube de Flower, la
profundidad en seguridad y la orientación a producción de NVIDIA FLARE, o la
extensibilidad por componentes y la experimentación algorítmica de FedML.

%  (dentro de \mainmatter, tras el capítulo de resultados de rendimiento)
% =====================================================================

\chapter{Resultados de flexibilidad}\label{chap:flexibilidad}

Este capítulo presenta el tercer eje de la comparación: la flexibilidad de
cada framework. A diferencia de la comparativa de rendimiento, que es
cuantitativa, y de la valoración de usabilidad, que recoge la percepción de uso,
la flexibilidad describe una propiedad arquitectónica: hasta qué punto la
herramienta permite adaptarse a escenarios distintos de los previstos por defecto.
Conviene recordar que este eje no debe confundirse con la participación parcial
de clientes medida en el experimento e3, que es una variación cuantitativa del
escenario; aquí el término flexibilidad se reserva para la capacidad de
modificar, sustituir o ampliar el comportamiento del framework.

La valoración se apoya en los diez criterios F1--F10 definidos en la
Sección~\ref{sec:flex-criterios} y se puntúa con la escala de tres niveles
(0, 1, 2) de la Tabla~\ref{tab:escala-flexibilidad}. Para cada framework se
distinguen dos planos de evidencia. El primero es la evidencia experimental:
cambios efectivamente realizados en el código, la configuración y los scripts, y
ejecuciones reales sobre el escenario común. El segundo es la evidencia
documental: funcionalidades avanzadas descritas en la documentación oficial pero
no implementadas en los experimentos del trabajo. Esta separación es deliberada y
evita presentar como probado aquello que solo se ha revisado en la documentación,
manteniendo así la comparabilidad del diseño experimental.

\section{Flexibilidad de Flower}\label{sec:flex-flower}

Flower es un framework agnóstico respecto a la biblioteca de aprendizaje
automático y con una separación clara entre la lógica del cliente y la del
servidor~\parencite{beutel2020flower}, rasgos que condicionan directamente su
flexibilidad. La evaluación experimental se centró en cinco capacidades: modificar
la configuración federada, sustituir la estrategia de agregación, integrar un
modelo y un conjunto de datos propios en PyTorch, registrar métricas personalizadas
y alternar el modo de ejecución. Todas las ejecuciones se realizaron con
\texttt{flwr} 1.29.0 sobre el backend de simulación Ray~\parencite{moritz2018ray},
con Python 3.12 (entorno \texttt{tfg-flower}), una GPU NVIDIA GeForce RTX~3050
Laptop de 6\,GB y el modelo ResNet-18 adaptado sobre CIFAR-10 descrito en el
capítulo de metodología.

\subsection{Adaptaciones para habilitar la evaluación}\label{subsec:flex-flower-adapt}

La primera evidencia de flexibilidad es la propia parametrización del escenario.
Para que Flower recibiera la configuración de federación en cada ejecución, se
ajustó el script de lanzamiento de modo que toda la configuración viajara en una
única cadena pasada a \texttt{flwr run}, como muestra el Código~\ref{cod:flower-fedconfig}.
Con ello fue posible variar el número de supernodos o clientes y la fracción de
GPU asignada a cada uno en función del modo de ejecución, sin tocar el código de
la aplicación. El parámetro \texttt{init-args-num-cpus=2} no es cosmético: se
incorporó para resolver un agotamiento de memoria RAM del anfitrión, una
limitación que se detalla en la Sección~\ref{subsec:flex-flower-lim}.

\begin{bashcode}[title={run\_flower\_experiments.sh}, label={cod:flower-fedconfig}]
--federation-config="num-supernodes=$clients \
  init-args-num-cpus=2 \
  client-resources-num-cpus=1 \
  client-resources-num-gpus=$gpu_share"
\end{bashcode}

Una segunda adaptación afectó a la declaración de la estrategia. Flower exige que
las claves que se sobrescriben mediante \texttt{-{}-run-config} estén previamente
declaradas en la sección \texttt{[tool.flwr.app.config]} del fichero
\texttt{pyproject.toml}; de lo contrario, la ejecución aborta con el mensaje
\mintinline{text}{Key 'strategy' is not present in the main dictionary}. Por ese
motivo se añadió la clave \texttt{strategy} a la configuración de la aplicación,
junto con el resto de hiperparámetros del escenario (Código~\ref{cod:flower-pyproject}).
Asimismo, el valor de \texttt{learning-rate} se fijó en 0,02 para que coincidiera
con el que el script aplica de forma efectiva, eliminando una posible fuente de
confusión entre el valor declarado y el realmente utilizado.

\begin{codigo}[title={pyproject.toml}, label={cod:flower-pyproject}]{toml}
[tool.flwr.app.config]
num-server-rounds = 3
fraction-train = 0.5
fraction-evaluate = 0.0
local-epochs = 1
learning-rate = 0.02
batch-size = 32
strategy = "fedavg"
\end{codigo}

\subsection{Selección dinámica de la estrategia de agregación}\label{subsec:flex-flower-strategy}

El núcleo de la evaluación experimental de flexibilidad fue la sustitución de la
estrategia de agregación sin alterar el resto del sistema. Para ello se modificó
el servidor de Flower de manera que la estrategia se seleccionara dinámicamente a
partir de la configuración de ejecución, leyendo la clave \texttt{strategy} de
\mintinline{python}{context.run_config}. El servidor instancia entonces la clase
correspondiente y, en el caso de las estrategias adaptativas, fija los
hiperparámetros del optimizador del servidor antes de iniciar el entrenamiento,
como recoge el Código~\ref{cod:flower-strategy}. Gracias a este diseño, el mismo
código de cliente, modelo y conjunto de datos puede ejecutarse con distintas
estrategias modificando únicamente la lógica del servidor.

\begin{pythoncode}[title={server\_app.py}, label={cod:flower-strategy}]
from flwr.serverapp.strategy import FedAvg, FedAdagrad, FedAdam

strategy_name = str(context.run_config.get("strategy", "fedavg")).lower()

if strategy_name == "fedadam":
    strategy_class = FedAdam
    adaptive_kwargs = dict(
        eta=float(context.run_config.get("server-lr", 0.01)),
        tau=float(context.run_config.get("tau", 0.01)),
    )
elif strategy_name == "fedadagrad":
    strategy_class = FedAdagrad
    adaptive_kwargs = dict(
        eta=float(context.run_config.get("server-lr", 0.01)),
        tau=float(context.run_config.get("tau", 0.01)),
    )
else:
    strategy_class = FedAvg
    adaptive_kwargs = {}

strategy = strategy_class(**common_kwargs, **adaptive_kwargs)
\end{pythoncode}

Las estrategias FedAdam y FedAdagrad pertenecen a la familia de optimización
federada adaptativa propuesta por \textcite{reddi2021adaptive}, que generaliza
FedAvg sustituyendo el promedio del servidor por un optimizador con momento y tasa
de aprendizaje propios. Con los valores por defecto de Flower
(\texttt{eta}~=~0,1 y \texttt{tau}~=~0,001), ambas estrategias divergieron
numéricamente en este escenario: la pérdida del modelo global creció hasta el
orden de $10^{9}$ mientras las pérdidas locales de entrenamiento descendían con
normalidad. El problema se localiza en el paso de agregación del servidor, donde
una tasa de aprendizaje alta combinada con un factor de adaptabilidad pequeño
produce primeros pasos desproporcionados antes de que el segundo momento se
estabilice. Reducir la tasa del servidor a \texttt{eta}~=~0,01 y elevar el factor
a \texttt{tau}~=~0,01 estabilizó el entrenamiento. Lejos de restar valor a la
evaluación, este ajuste evidencia precisamente la flexibilidad buscada: Flower
expone esos hiperparámetros de agregación y permite configurarlos desde el
servidor.

\subsection{Mini-experimento de flexibilidad}\label{subsec:flex-flower-mini}

Para comprobar la sustitución de estrategias de forma controlada se definió una
variante reducida del experimento e3, denominada \texttt{e3\_flex\_mini}, cuyos
parámetros recoge la Tabla~\ref{tab:flex-flower-config}. La finalidad de esta
prueba no era comparar el rendimiento de aprendizaje, sino demostrar que Flower
permite cambiar la estrategia federada y la configuración de participación sin
rehacer la implementación, reduciendo a la vez el coste temporal de las pruebas.
Se ejecutó con una única semilla (42) y una repetición, por lo que las cifras de
aprendizaje son auxiliares y no llevan estimación de variabilidad.

\begin{table}[htbp]
  \centering
  \caption{Configuración del mini-experimento de flexibilidad de Flower.}
  \label{tab:flex-flower-config}
  \begin{tabular}{lc}
    \toprule
    \textbf{Parámetro} & \textbf{Valor} \\
    \midrule
    Clientes totales & 8 \\
    Clientes por ronda & 4 \\
    Fracción de entrenamiento & 0,5 \\
    Rondas federadas & 3 \\
    Épocas locales & 1 \\
    Tamaño de lote & 32 \\
    Tasa de aprendizaje (cliente) & 0,02 \\
    \texttt{eta} servidor (FedAdam/FedAdagrad) & 0,01 \\
    \texttt{tau} (FedAdam/FedAdagrad) & 0,01 \\
    Modo de ejecución & Secuencial \\
    Repeticiones / semilla & 1 / 42 \\
    \bottomrule
  \end{tabular}
\end{table}

Los tres mini-experimentos terminaron correctamente, lo que confirma que el mismo
flujo se ejecutó con FedAvg, FedAdam y FedAdagrad manteniendo constantes el
conjunto de datos, el modelo, el número de clientes, los clientes por ronda, las
rondas, las épocas locales, el tamaño de lote y la tasa de aprendizaje del
cliente. La Tabla~\ref{tab:flex-flower-resultados} resume las métricas obtenidas.

\begin{table}[htbp]
  \centering
  \caption{Resultados del mini-experimento de flexibilidad de Flower con tres
  estrategias de agregación. Las cifras son auxiliares y no deben interpretarse
  como una comparación de rendimiento. La memoria de GPU corresponde al pico
  reservado por PyTorch al proceso del cliente, métrica comparable entre
  estrategias.}
  \label{tab:flex-flower-resultados}
  \begin{tabular}{lcccrcr}
    \toprule
    \textbf{Estrategia} & \textbf{Estado} & \textbf{Acc.} & \textbf{Loss} &
    \textbf{Dur. (s)} & \textbf{GPU proc. (MiB)} & \textbf{Energía (Wh)} \\
    \midrule
    FedAvg     & success & 0,4257 & 1,5265 & 132 & $\approx$289 & 1,272 \\
    FedAdam    & success & 0,1410 & 2,3401 & 124 & $\approx$289 & 1,252 \\
    FedAdagrad & success & 0,1000 & 18,7832 & 144 & $\approx$289 & 1,123 \\
    \bottomrule
  \end{tabular}
\end{table}

Estos resultados no permiten concluir que una estrategia sea superior a otra, por
dos motivos. En primer lugar, se trata de un mini-experimento de solo tres rondas
y una época local, insuficiente para que las estrategias adaptativas estabilicen
su segundo momento y converjan; con los hiperparámetros ajustados ya no divergen,
pero todavía no alcanzan el comportamiento de FedAvg. En segundo lugar, se ejecutó
con una única semilla, sin medida de variabilidad. Conviene además una
advertencia de medición: la memoria de GPU reportada por \texttt{nvidia-smi} a
nivel de sistema mostró un valor anómalo de 1623~MiB en la ejecución de
FedAdagrad frente a 757~MiB en las otras dos, atribuible a memoria ocupada por
otras aplicaciones del equipo durante esa corrida; el consumo real por proceso fue
prácticamente idéntico ($\approx$289~MiB) en las tres estrategias, por lo que para
cualquier comparación debe emplearse la métrica por proceso. La conclusión
experimental válida es que Flower permitió sustituir la estrategia federada entre
FedAvg, FedAdam y FedAdagrad sin modificar el cliente, el modelo ni el conjunto de
datos, demostrando su flexibilidad para alterar la lógica de agregación desde el
servidor y la configuración de ejecución.

\subsection{Integración de modelo, datos y métricas}\label{subsec:flex-flower-integracion}

Más allá de las estrategias, la evaluación confirmó que Flower integra sin
fricción un modelo y un conjunto de datos propios. El modelo es la ResNet-18
adaptada a CIFAR-10, construida en PyTorch sustituyendo la primera convolución por
una de núcleo 3$\times$3 y eliminando el \textit{maxpool} inicial, mientras que el
conjunto de datos se carga mediante \mintinline{python}{FederatedDataset} con un
particionador IID. El bucle de entrenamiento local es código propio del cliente:
emplea descenso de gradiente estocástico con momento y \textit{weight decay},
entropía cruzada como función de pérdida y un número de épocas, tasa y tamaño de
lote configurables. Además, el cliente registra métricas personalizadas, entre
ellas la pérdida de entrenamiento, el número de ejemplos, la duración y el pico de
memoria de GPU asignada por PyTorch, que el script de ejecución vuelca junto con
los registros completos en ficheros de resumen reutilizables. Esta combinación de
modelo propio, bucle de entrenamiento editable y métricas a medida sustenta una
valoración alta en los criterios de personalización, integración y
reproducibilidad.

\subsection{Capacidades avanzadas según la documentación}\label{subsec:flex-flower-doc}

Más allá de lo ejecutado, la documentación oficial de Flower (versión 1.31) describe
un conjunto amplio de técnicas de flexibilidad que se valoran como evidencia
documental. En el plano de la configuración y el control del servidor, se detallan
cuatro vías de personalización de las estrategias: ajustar los parámetros de una
estrategia existente (\texttt{fraction\_train}, \texttt{fraction\_evaluate},
\texttt{min\_available\_nodes}); inyectar funciones de devolución como
\texttt{evaluate\_fn} que el servidor ejecuta en momentos concretos; sobrescribir
métodos internos como \texttt{configure\_train}, \texttt{aggregate\_fit} o
\texttt{aggregate\_evaluate}; e implementar una estrategia nueva heredando de
\texttt{BaseStrategy}~\parencite{beutel2020flower}. La misma documentación muestra cómo
enviar configuraciones por ronda mediante \texttt{ConfigRecord} —por ejemplo, reducir
progresivamente la tasa de aprendizaje leyendo \texttt{server\_round}— y cómo agregar
métricas con criterios propios a través de \texttt{evaluate\_metrics\_aggr\_fn}, en
lugar del promedio ponderado por número de ejemplos. Los clientes, sin estado por
defecto, pueden conservar información entre rondas mediante \texttt{Context.state}, lo
que habilita técnicas como el \textit{checkpointing} en el cliente.

Una pieza distintiva son los \textit{Mods}, funciones que interceptan los mensajes
entre servidor y cliente para transformarlos sin alterar la lógica principal; se
encadenan en un orden definido y permiten, por ejemplo, recortar gradientes
(\texttt{fixedclipping\_mod}, \texttt{adaptiveclipping\_mod}) o limitar el tamaño de
los mensajes (\texttt{message\_size\_mod})~\parencite{beutel2020flower}. Sobre esta
base se construye la privacidad diferencial, documentada en modo preliminar en tres
variantes —central del lado del servidor, central del lado del cliente y local
mediante \texttt{LocalDpMod}—, todas con recorte y ruido gaussiano configurables, a las
que se suma la agregación segura. La documentación también ejemplifica esquemas no
estándar, como FedBN —que excluye los parámetros de normalización por lotes del
intercambio para preservar las estadísticas locales de cada cliente—, el aprendizaje
federado vertical y la integración con XGBoost.

La documentación recoge asimismo soporte para el entrenamiento asíncrono, si bien de
carácter experimental y fuera de las guías principales. Las notas de versión describen
una \textit{Driver API} experimental (versión 1.2) que ejecuta un servidor programable,
asíncrono y multi-inquilino, y la API \textit{Flower Next} (versión 1.8), un conjunto de
primitivas de bajo nivel que permite personalizar el bucle de entrenamiento e incluye,
de forma explícita, soporte para aprendizaje federado asíncrono y esquemas cíclicos,
con un \texttt{ServerApp} capaz de permanecer activo de forma
indefinida~\parencite{beutel2020flower}. Por su naturaleza en desarrollo y de bajo
nivel, esta capacidad se valora como soporte parcial y no como una funcionalidad
directa y estable. Todas estas capacidades se reconocen como existentes, pero no se
contabilizan como demostradas empíricamente en este trabajo.

\subsection{Ejecución, despliegue y operación según la documentación}\label{subsec:flex-flower-deploy}

La flexibilidad de Flower se extiende a su modelo de ejecución y despliegue. En local,
los comandos \texttt{flwr run} y \texttt{flwr list} levantan automáticamente un
\textit{SuperLink} que gestiona las ejecuciones, mientras que las simulaciones se
lanzan sobre el \textit{Simulation Runtime}, cuyo número de \textit{SuperNodes} y
recursos por cliente se fijan de forma persistente o se sobrescriben por ejecución;
la documentación contempla incluso simulaciones sobre un clúster Ray
multi-anfitrión~\parencite{beutel2020flower}. Para despliegues distribuidos, la
federación se compone de un \textit{SuperLink} que coordina y varios
\textit{SuperNodes} que ejecutan los clientes, modelo que la documentación lleva hasta
entornos de producción en la nube —Google Cloud sobre Kubernetes, Microsoft Azure y
Red Hat OpenShift— e incluso a la interconexión de varios clústeres mediante Skupper.

Este recorrido se acompaña de mecanismos de seguridad operativa que también cuentan
para la flexibilidad: cifrado de las comunicaciones mediante TLS con certificados
propios, autenticación de \textit{SuperNodes} basada en firmas y lista blanca de
claves públicas, autenticación de cuentas de usuario mediante OpenID Connect con
autorización fina y un registro de auditoría en formato JSON que traza qué actor
ejecuta cada acción~\parencite{beutel2020flower}. Por último, la interfaz de línea de
comandos puede emitir su salida en formato JSON mediante la opción
\texttt{-{}-format json} para integrarse en flujos automatizados, y las herramientas de
gestión de federaciones permiten inspeccionar nodos, ejecuciones y estado. Ninguna de
estas capacidades se ejercitó en los experimentos, pero refuerzan la valoración de
Flower en los modos de ejecución y despliegue y en la integración con herramientas
externas.

\subsection{Limitaciones observadas}\label{subsec:flex-flower-lim}

El hallazgo más relevante de la evaluación afecta a la memoria RAM del anfitrión y
no a la GPU. El motor de simulación de Flower delega en Ray, que al iniciarse
detecta los procesadores del equipo (veinte en este caso) y prelanza un proceso
\textit{worker} ocioso por cada uno; cada worker carga el intérprete y las
bibliotecas (en torno a 0,28~GB), de modo que el conjunto de workers, sumado al
proceso de simulación y al resto del sistema, agotaba los aproximadamente 14,6~GB
de RAM disponibles. Ray respondía terminando procesos con un error de memoria
agotada, lo que ocurría incluso en modo secuencial, donde a priori solo hay un
cliente entrenando a la vez. La solución fue limitar el número de procesadores que
ve la \textit{Simulation Runtime} mediante \texttt{init-args-num-cpus=2}, lo que
reduce drásticamente la huella de RAM. Se trata de un resultado comparativo de
interés: el backend basado en Ray presenta una huella de memoria de anfitrión
notablemente mayor que el simulador de NVIDIA FLARE o el backend de proceso único
de FedML, que se ejecutan en un único proceso, por lo que en equipos con poca RAM
la simulación de Flower obliga a acotar los recursos manualmente.

La ejecución concurrente, probada también mediante Ray con una fracción de GPU por
cliente de 0,5, quedó condicionada por esta misma limitación de memoria, en línea
con lo descrito en el capítulo de resultados de rendimiento. Por su parte, la
divergencia de las estrategias adaptativas con los valores por defecto, ya
comentada, es una limitación de uso y no del framework: las estrategias de la
familia FedOpt son sensibles a la tasa de aprendizaje del servidor y requieren un
ajuste deliberado, especialmente con pocas rondas~\parencite{reddi2021adaptive}.
Ninguna de estas limitaciones impidió completar la evaluación, pero todas matizan
la valoración final.

\subsection{Valoración de la flexibilidad de Flower}\label{subsec:flex-flower-score}

La Tabla~\ref{tab:flex-flower-score} resume la valoración de Flower sobre los diez
criterios F1--F10, distinguiendo el tipo de evidencia que sustenta cada
puntuación. Los criterios de personalización del modelo y los datos (F1), del
entrenamiento local (F2), de la estrategia de agregación (F3) y de los algoritmos
federados disponibles (F4) se valoran con la máxima puntuación a partir de
evidencia experimental directa, al igual que los modos de ejecución y despliegue
(F6), la integración con herramientas externas (F7) y la reproducibilidad (F10).
El soporte de aprendizaje asíncrono (F5), la seguridad y privacidad (F8) y la
modificación de la arquitectura federada (F9) reciben puntuación parcial por estar
respaldados solo por documentación; en el caso de F5, mediante interfaces de bajo
nivel todavía en desarrollo (la \textit{Driver API} y \textit{Flower Next}), por lo
que se reconoce el soporte sin elevarlo a la máxima puntuación.

\begin{table}[htbp]
  \centering
  \caption{Valoración de la flexibilidad de Flower según los criterios F1--F10.
  El tipo de evidencia distingue las capacidades ejecutadas en los experimentos
  (experimental) de las descritas en la documentación oficial (documental).}
  \label{tab:flex-flower-score}
  \begin{tabular}{clcl}
    \toprule
    \textbf{Criterio} & \textbf{Descripción} & \textbf{Punt.} & \textbf{Evidencia} \\
    \midrule
    F1  & Personalización del modelo y el dataset      & 2/2 & Experimental \\
    F2  & Personalización del entrenamiento local      & 2/2 & Experimental \\
    F3  & Estrategias de agregación personalizadas     & 2/2 & Experimental \\
    F4  & Algoritmos federados disponibles             & 2/2 & Experimental \\
    F5  & Aprendizaje federado asíncrono               & 1/2 & Documental \\
    F6  & Modos de ejecución y despliegue              & 2/2 & Exp.\ + documental \\
    F7  & Integración con herramientas externas        & 2/2 & Exp.\ + documental \\
    F8  & Seguridad, privacidad y extensiones          & 1/2 & Documental \\
    F9  & Modificación de la arquitectura federada     & 1/2 & Documental \\
    F10 & Reproducibilidad y configuración             & 2/2 & Experimental \\
    \midrule
    \multicolumn{2}{r}{\textbf{Total}} & \multicolumn{2}{l}{\textbf{17/20}} \\
    \bottomrule
  \end{tabular}
\end{table}

En conjunto, Flower obtiene una valoración alta en flexibilidad (17/20). La
evidencia experimental muestra una herramienta cómoda de adaptar:
permitió integrar un modelo y un conjunto de datos propios en PyTorch, editar el
bucle de entrenamiento local, sustituir la estrategia de agregación entre FedAvg,
FedAdam y FedAdagrad desde el servidor —incluyendo el ajuste de sus
hiperparámetros— y automatizar la recogida de métricas. La valoración no alcanza
el máximo porque varias capacidades avanzadas —el aprendizaje asíncrono, la
privacidad diferencial, la agregación segura o los esquemas no horizontales— solo se
verificaron documentalmente, y porque el soporte asíncrono se apoya en interfaces
todavía experimentales; a ello se añade que la viabilidad práctica del modo
concurrente quedó condicionada por la elevada huella de memoria del backend Ray en
el equipo empleado. Estas conclusiones se contrastan con las de NVIDIA FLARE y FedML
en la comparativa final del capítulo.

% =====================================================================
%  SECCIÓN AÑADIDA: Flexibilidad de NVIDIA FLARE
% =====================================================================

\section{Flexibilidad de NVIDIA FLARE}\label{sec:flex-nvflare}

NVIDIA FLARE se presenta como un SDK de aprendizaje federado agnóstico al dominio
y compatible con flujos de trabajo de aprendizaje automático ya existentes, con
una orientación que abarca desde la simulación local hasta el despliegue en
producción y en el borde~\parencite{roth2022nvflare,nvflareWelcome}. La evaluación
experimental de su flexibilidad se apoya en un paquete de cinco ejecuciones de la
variante reducida del experimento e3, la misma \texttt{e3\_flex\_mini} empleada con
Flower (Subsección~\ref{subsec:flex-flower-mini}), con ocho clientes totales,
cuatro por ronda, tres rondas, una época local, tamaño de lote 32, tasa de
aprendizaje 0,02 y semilla 42, en modo secuencial con un único hilo del simulador.
Conviene precisar el alcance: el paquete inspeccionado contiene una selección de
suites concretas y no la totalidad del histórico de ejecuciones, por lo que las
conclusiones experimentales se refieren únicamente a los archivos incluidos en él.

\subsection{Diseño modular: jobs, componentes y workflows}\label{subsec:flex-nvflare-modular}

La flexibilidad de NVIDIA FLARE descansa en su modelo de configuración modular. La
unidad de ejecución es el \textit{job}, que describe por completo un experimento
mediante ficheros de configuración de servidor y de cliente
(\texttt{config\_fed\_server.conf} y \texttt{config\_fed\_client.conf}), en los que
cada elemento del sistema se declara como un componente con su ruta de clase y sus
argumentos de construcción~\parencite{nvflareComponents}. Al desplegar el job, el
servidor y los clientes analizan esos ficheros e instancian dinámicamente los
componentes declarados, registrándolos en el bucle de eventos del framework; cada
componente se identifica con un \texttt{component\_id} que permite referenciarlo
desde otros componentes o workflows~\parencite{nvflareComponents}. Este diseño es
el que permite sustituir un agregador, un persistor, un \textit{executor}, un
generador de objetos compartibles (\textit{Shareable}) o un filtro modificando la
configuración, sin reescribir la aplicación completa.

Sobre esa base de bajo nivel, la documentación ofrece además una capa de
\textit{recipes} que empaqueta algoritmos y workflows habituales —FedAvg, FedProx,
FedOpt, SCAFFOLD, aprendizaje cíclico, XGBoost, Sklearn, estadísticas federadas,
evaluación cruzada, PSI, integración con Flower y swarm learning, entre
otros~\parencite{nvflareRecipes}. Esta dualidad, configuración granular por
componentes y configuración de alto nivel por recetas, sustenta una valoración
elevada en personalización y reproducibilidad, y constituye el punto de partida de
las adaptaciones realizadas para probar estrategias alternativas.

\subsection{Adaptaciones para soportar las estrategias}\label{subsec:flex-nvflare-adapt}

El objetivo de la prueba fue comprobar si un mismo job podía adaptarse para
ejecutar FedAvg, FedProx y variantes FedOpt del servidor (FedAdagrad y FedAdam),
manteniendo constante la estructura general del experimento. Para ello se
introdujeron tres cambios de naturaleza distinta. En primer lugar, la selección de
estrategia se trasladó al cliente, que admite los valores \texttt{fedavg},
\texttt{fedprox}, \texttt{fedadam} y \texttt{fedadagrad} mediante un argumento de
línea de órdenes que, por defecto, toma su valor de la variable de entorno
\texttt{TFG\_FL\_STRATEGY}. En segundo lugar, FedProx se implementó en el
entrenamiento local: cuando la estrategia es \texttt{fedprox} y el coeficiente
proximal es positivo, el entrenamiento añade un término que penaliza la desviación
de los pesos locales respecto al modelo global recibido al inicio de la tarea,
manteniendo el descenso de gradiente estocástico como optimizador
local~\parencite{nvflareAlgorithms}.

El tercer cambio afecta al servidor y es el más relevante para la flexibilidad de
agregación. Para las variantes FedOpt, el script de ejecución parchea el fichero de
configuración del servidor sustituyendo el generador de compartibles por defecto,
\mintinline{text}{FullModelShareableGenerator}, por
\mintinline{text}{PTFedOptModelShareableGenerator}, y añade un componente de modelo
fuente requerido por este generador, como ilustra el
Código~\ref{cod:nvflare-fedopt}. Este generador actualiza el modelo global mediante
un optimizador de PyTorch —y, opcionalmente, un planificador de tasa de
aprendizaje— cuyos argumentos incluyen \texttt{optimizer\_args},
\texttt{lr\_scheduler\_args}, \texttt{source\_model} y
\texttt{device}~\parencite{nvflareFedopt}. En consecuencia, FedAdam y FedAdagrad no
son recetas independientes, sino configuraciones de FedOpt que se obtienen
seleccionando \mintinline{text}{torch.optim.Adam} o
\mintinline{text}{torch.optim.Adagrad} como optimizador del servidor.

\begin{codigo}[title={Fragmento representativo de config\_fed\_server.conf (FedOpt)}, label={cod:nvflare-fedopt}]{json}
"shareable_generator": {
  "path": "nvflare.app_opt.pt.fedopt.PTFedOptModelShareableGenerator",
  "args": {
    "source_model": "model",
    "optimizer_args": { "path": "torch.optim.Adagrad" }
  }
},
"components": [
  { "id": "model", "path": "<ruta de la ResNet-18 adaptada>" }
]
\end{codigo}

Una precaución metodológica relevante surgió con el dispositivo de ejecución. En
todos los ficheros de configuración del cliente aparece \texttt{-{}-device cuda},
pero en la ejecución de FedAdam en CPU el propio cliente seleccionó CPU porque
\mintinline{python}{torch.cuda.is_available()} devolvió falso al lanzarse con
\texttt{CUDA\_VISIBLE\_DEVICES} vacío. Por ello, el dispositivo realmente empleado
debe leerse del registro de ejecución y no del fichero de configuración, criterio
que se ha seguido al elaborar la Tabla~\ref{tab:flex-nvflare-resultados}.

\subsection{Mini-experimento de flexibilidad}\label{subsec:flex-nvflare-mini}

La Tabla~\ref{tab:flex-nvflare-resultados} resume las cinco ejecuciones. FedAvg
sirvió de referencia y se ejecutó correctamente con el workflow \textit{Scatter and
Gather}, en el que el servidor distribuye los pesos iniciales, los clientes
entrenan localmente y devuelven sus pesos, y el sistema los agrega por promedio
ponderado~\parencite{nvflareAlgorithms}. A partir de esa base, FedProx finalizó las
tres rondas y los registros confirmaron \texttt{strategy=fedprox} con un coeficiente
proximal \texttt{fedprox\_mu=0,001} activo durante el entrenamiento. FedAdagrad
también completó las tres rondas como FedOpt real —y no como mera etiqueta— porque
la configuración del servidor incorporaba el generador
\mintinline{text}{PTFedOptModelShareableGenerator} y el registro mostraba las
actualizaciones del servidor con \mintinline{text}{torch.optim.Adagrad}.

\begin{table}[htbp]
  \centering
  \caption{Resultados del mini-experimento de flexibilidad de NVIDIA FLARE. Las
  cifras de aprendizaje son auxiliares (tres rondas, una época, una semilla) y no
  constituyen una comparación de rendimiento. El dispositivo real procede del
  registro de ejecución, no del fichero de configuración del cliente.}
  \label{tab:flex-nvflare-resultados}
  \begin{tabular}{llccccrr}
    \toprule
    \textbf{Estrategia} & \textbf{Estado} & \textbf{Disp.} & \textbf{FedOpt} &
    \textbf{Rondas} & \textbf{Dur. (s)} & \textbf{Acc.} & \textbf{Loss} \\
    \midrule
    FedAvg      & success & \texttt{cuda:0} & No  & 3 & 205  & 0,4178 & 1,5601 \\
    FedProx     & success & \texttt{cuda:0} & No  & 3 & 225  & 0,4045 & 1,5876 \\
    FedAdagrad  & success & \texttt{cuda:0} & Sí  & 3 & 199  & 0,1402 & 2,3785 \\
    FedAdam     & failed  & \texttt{cuda:0} & Sí  & 0 & 77   & ---    & ---    \\
    FedAdam     & success & \texttt{cpu}    & Sí  & 3 & 1280 & 0,2586 & 3,4439 \\
    \bottomrule
  \end{tabular}
\end{table}

El caso de FedAdam exige una lectura cuidadosa. La estrategia llegó a configurarse
realmente mediante FedOpt con \mintinline{text}{torch.optim.Adam}, pero su ejecución
en GPU abortó durante el paso de actualización del servidor
(\mintinline{python}{torch.optim.Adam.step()}) con un error de inicialización de
CUDA, registrado como \texttt{FATAL\_SYSTEM\_ERROR}. Para aislar el problema, la
prueba se repitió en CPU, donde FedAdam completó las tres rondas; esta ejecución
valida FedAdam como estrategia configurable, pero su duración (1280~s) es muy
superior precisamente por entrenarse en CPU, de modo que no debe mezclarse con la
comparativa de rendimiento en GPU del capítulo correspondiente. En coherencia con
ello, la fila de la ejecución fallida no reporta precisión ni pérdida, ya que no
completó ninguna ronda.

Estos resultados no permiten ordenar las estrategias por calidad, por los mismos
motivos señalados para Flower: tres rondas, una época local y una sola semilla son
insuficientes para que las variantes adaptativas converjan. Su valor es
demostrar que NVIDIA FLARE permite adaptar tanto el entrenamiento local como la
lógica de actualización del servidor mediante configuración y componentes. Importa
además una cautela de interpretación: la validez de FedAdagrad y FedAdam no se
sustenta en el nombre de la ejecución, sino en la presencia del generador FedOpt y
del optimizador del servidor en los registros. Del mismo modo, el valor
\texttt{fedprox\_mu=0,001} que aparece en la configuración no implica que FedProx se
aplique a todas las estrategias: en FedAdagrad y FedAdam el entrenamiento imprime
\texttt{fedprox\_mu=0,0}, porque el término proximal solo se activa cuando la
estrategia es FedProx.

De esta experiencia se desprende una observación comparativa de interés. Frente a
Flower, donde la estrategia se sustituía cambiando una única clase en el servidor,
NVIDIA FLARE exigió modificar componentes del servidor, añadir un modelo fuente para
FedOpt y comprobar en los registros que la estrategia configurada era realmente la
que se ejecutaba. La flexibilidad de NVIDIA FLARE es, por tanto, amplia, pero su
ejercicio práctico conlleva una complejidad de configuración mayor.

\subsection{Capacidades avanzadas según la documentación}\label{subsec:flex-nvflare-doc}

Como en el caso de Flower, varias capacidades relevantes no se ejecutaron en los
experimentos y se valoran únicamente como evidencia documental. La documentación
oficial las articula en torno a las \textit{Job Recipes}, una API de alto nivel que
encapsula un trabajo federado exponiendo solo los parámetros clave —\texttt{min\_clients},
\texttt{num\_rounds}, modelo y \textit{script} de entrenamiento— y ocultando la
complejidad de controladores y \textit{workflows}~\parencite{nvflareRecipes}. Estas
recetas se ejecutan sobre tres entornos intercambiables: \texttt{SimEnv} simula los
clientes en procesos locales, \texttt{PocEnv} los lanza como procesos separados en la
misma máquina y \texttt{ProdEnv} despliega el trabajo en una instalación de
producción; un mismo trabajo puede así trasladarse de la simulación al despliegue sin
cambiar el código~\parencite{nvflareRecipes}. El catálogo de recetas cubre FedAvg,
FedProx, FedOpt —donde el optimizador del servidor se fija con un argumento que indica
su ruta de clase e hiperparámetros, por ejemplo \texttt{torch.optim.SGD} con tasa y
momento—, SCAFFOLD y el aprendizaje cíclico, además de XGBoost federado en sus
variantes horizontal, \textit{bagging} y vertical~\parencite{nvflareAlgorithms,nvflareFedopt}.

En cuanto al aprendizaje asíncrono, la receta \texttt{EdgeFedBuffRecipe}, orientada a
dispositivos de borde, admite entrenamiento síncrono y asíncrono: en modo síncrono el
servidor espera todas las actualizaciones antes de generar un nuevo modelo, mientras
que en modo asíncrono lo genera en cuanto recibe la primera, lo que reduce el tiempo
total cuando los dispositivos presentan latencias heterogéneas~\parencite{nvflareEdge,nvflareWelcome}.
La receta expone parámetros como el número máximo de versiones de modelo activas o el
tamaño de selección de dispositivos por ronda. No obstante, la propia documentación
advierte de que la asincronía general aún no está disponible para todas las recetas y
se concentra en el escenario de borde, por lo que su valoración se mantiene como
soporte documental parcial.

En materia de seguridad y privacidad, la documentación es notablemente extensa.
NVIDIA FLARE implementa la privacidad mediante un mecanismo general de filtros que
procesan los objetos compartibles antes o después de la comunicación, configurables en
la receta como \texttt{task\_data\_filters} o \texttt{task\_result\_filters}. Soporta
privacidad diferencial a nivel de muestra mediante la integración de DP-SGD con Opacus
—controlada por los parámetros \texttt{epsilon} y \texttt{delta}— y filtros de DP local
como \texttt{PercentilePrivacy}, que comparte solo las diferencias de pesos por encima
de un percentil, \texttt{SVTPrivacy}, basado en la técnica del vector disperso con ruido
de Laplace, y \texttt{StatisticsPrivacyFilter} para estadísticas
federadas~\parencite{nvflareDP}. A un nivel superior, la solución de computación
confidencial ejecuta la agregación dentro de un entorno de ejecución confiable
(\textit{Trusted Execution Environment}) sobre hardware como AMD SEV-SNP o Intel TDX,
protege el modelo y el código del cliente en máquinas virtuales confidenciales y
permite combinar privacidad diferencial con cifrado homomórfico o agregación
segura~\parencite{nvflareSecurity,nvflareConfidential}.

Por último, NVIDIA FLARE documenta esquemas descentralizados mediante los
\textit{workflows} controlados por clientes, construidos sobre las clases base
\texttt{ServerSideController} y \texttt{ClientSideController} para escenarios donde el
servidor no es plenamente confiable~\parencite{nvflareSwarm}. Entre ellos destacan el
aprendizaje cíclico, en el que cada cliente entrena con el modelo recibido del anterior
según un orden fijo o aleatorio~\parencite{nvflareCyclic}; el aprendizaje \textit{swarm},
que designa un cliente agregador distinto por ronda y mejora la resiliencia frente a
fallos del servidor~\parencite{nvflareSwarm}; y la evaluación cruzada entre sitios, que
permite a los clientes validar modelos de otros mediante comunicación entre pares sin
exponerlos al servidor. A esto se suman el aprendizaje vertical con XGBoost y las
estadísticas federadas, así como las utilidades para integrar el seguimiento de
experimentos con TensorBoard, MLflow o Weights \& Biases sobre cualquier
receta~\parencite{nvflareTracking,nvflareRecipes}. La documentación reconoce, además,
dos cautelas: la API de recetas se encuentra en estado de \textit{technical preview} y
no cubre todavía todos los algoritmos, y la configuración de la computación confidencial
exige infraestructura especializada.

\subsection{Limitaciones observadas}\label{subsec:flex-nvflare-lim}

La limitación más concreta de la evaluación fue el fallo de FedAdam en GPU por un
error de inicialización de CUDA durante la actualización FedOpt del servidor. Debe
interpretarse como una incidencia técnica del entorno, no como una carencia del
framework, puesto que la misma configuración completó las tres rondas en CPU; aun
así, impide presentar FedAdam como ejecución correcta en GPU. La segunda limitación
es la complejidad práctica ya señalada: a diferencia de Flower, alterar la lógica de
agregación obligó a intervenir en componentes del servidor y a validar
cuidadosamente los registros, lo que incrementa la probabilidad de errores de
configuración. Por último, el modo concurrente quedó condicionado por los recursos
del equipo, en línea con lo expuesto en el capítulo de resultados de rendimiento;
aunque la documentación describe un sistema de ejecución multi-job con gestión de
recursos, esa capacidad no pudo validarse de forma estable en el hardware empleado.

\subsection{Valoración de la flexibilidad de NVIDIA FLARE}\label{subsec:flex-nvflare-score}

La Tabla~\ref{tab:flex-nvflare-score} recoge la valoración sobre los criterios
F1--F10, con el mismo convenio aplicado a Flower: puntuación máxima cuando existe
evidencia experimental directa, puntuación parcial cuando la capacidad solo consta
en la documentación y puntuación nula cuando no se identifica soporte. NVIDIA FLARE
obtiene la máxima puntuación en personalización del modelo y los datos (F1), del
entrenamiento local (F2), de la estrategia de agregación (F3) y de los algoritmos
federados disponibles (F4), todos respaldados por las ejecuciones, así como en los
modos de ejecución y despliegue (F6), la integración con herramientas externas (F7)
y la reproducibilidad (F10). La diferencia respecto a Flower aparece en el
aprendizaje asíncrono (F5): mientras que en Flower no se identificó soporte, NVIDIA
FLARE lo documenta de forma explícita mediante FedBuff, lo que justifica una
puntuación parcial. La seguridad y privacidad (F8) y la modificación de la
arquitectura federada (F9) reciben también puntuación parcial por estar respaldadas
solo documentalmente.

\begin{table}[htbp]
  \centering
  \caption{Valoración de la flexibilidad de NVIDIA FLARE según los criterios
  F1--F10. El tipo de evidencia distingue las capacidades ejecutadas en los
  experimentos de las descritas en la documentación oficial.}
  \label{tab:flex-nvflare-score}
  \begin{tabular}{clcl}
    \toprule
    \textbf{Criterio} & \textbf{Descripción} & \textbf{Punt.} & \textbf{Evidencia} \\
    \midrule
    F1  & Personalización del modelo y el dataset      & 2/2 & Exp.\ + documental \\
    F2  & Personalización del entrenamiento local      & 2/2 & Experimental \\
    F3  & Estrategias de agregación personalizadas     & 2/2 & Experimental \\
    F4  & Algoritmos federados disponibles             & 2/2 & Exp.\ + documental \\
    F5  & Aprendizaje federado asíncrono               & 1/2 & Documental \\
    F6  & Modos de ejecución y despliegue              & 2/2 & Exp.\ + documental \\
    F7  & Integración con herramientas externas        & 2/2 & Exp.\ + documental \\
    F8  & Seguridad, privacidad y extensiones          & 1/2 & Documental \\
    F9  & Modificación de la arquitectura federada     & 1/2 & Documental \\
    F10 & Reproducibilidad y configuración             & 2/2 & Exp.\ + documental \\
    \midrule
    \multicolumn{2}{r}{\textbf{Total}} & \multicolumn{2}{l}{\textbf{17/20}} \\
    \bottomrule
  \end{tabular}
\end{table}

En conjunto, NVIDIA FLARE alcanza una valoración de flexibilidad de 17/20. La
evidencia experimental
demuestra que permite adaptar el entrenamiento local —incluido el término proximal
de FedProx— y, sobre todo, la lógica de agregación del servidor mediante FedOpt, con
FedAvg, FedProx y FedAdagrad ejecutados correctamente y FedAdam validado en CPU. A
ello se suma una de las superficies de capacidades documentadas más amplias del
estudio, especialmente sólida en seguridad y privacidad, que abarca aprendizaje
asíncrono, workflows cíclicos y \textit{swarm}, privacidad
diferencial, cifrado homomórfico y computación confidencial. El contrapunto es una
complejidad de configuración superior a la de Flower y el hecho de que esas
capacidades avanzadas solo se verificaron documentalmente; además, la ejecución de
FedAdam en GPU falló por un error de CUDA y el modo concurrente quedó limitado por el
hardware. Esta lectura se contrasta con la de los otros frameworks en la comparativa
final del capítulo.

% =====================================================================
%  SECCIÓN AÑADIDA: Flexibilidad de FedML
% =====================================================================

\section{Flexibilidad de FedML}\label{sec:flex-fedml}

FedML es una biblioteca de investigación y banco de pruebas para aprendizaje
federado~\parencite{he2020fedml}, documentada actualmente dentro de la plataforma
TensorOpera aunque la librería y los ejemplos siguen empleando el paquete
\texttt{fedml}. Su documentación organiza el framework en varios escenarios
—simulación, cross-silo y cross-device— y ofrece un catálogo amplio de algoritmos
de referencia. La evaluación de su flexibilidad combina la implementación propia
realizada para la comparativa, que extiende las abstracciones de FedML, con una
prueba reducida de sustitución de estrategias. Conviene una precisión de alcance:
por el coste computacional, la prueba reducida de FedML se ejecutó con cuatro
clientes totales, dos por ronda y dos rondas, una configuración algo menor que la
empleada con Flower y NVIDIA FLARE; las métricas de aprendizaje son, por tanto, una
comprobación funcional y no una comparación de rendimiento.

\subsection{Configuración y componentes personalizados}\label{subsec:flex-fedml-config}

La configuración de FedML se apoya en un fichero \texttt{fedml\_config.yaml}
estructurado en secciones (\texttt{common\_args}, \texttt{data\_args},
\texttt{model\_args}, \texttt{train\_args}, \texttt{device\_args},
\texttt{comm\_args} y \texttt{tracking\_args}), donde se declaran los parámetros
federados como \texttt{federated\_optimizer}, \texttt{client\_num\_in\_total},
\texttt{client\_num\_per\_round}, \texttt{comm\_round}, \texttt{epochs},
\texttt{batch\_size}, \texttt{client\_optimizer} y \texttt{learning\_rate}~\parencite{fedmlSimApi}.
Esta organización permitió variar clientes, rondas, épocas e hiperparámetros entre
ejecuciones sin reescribir la aplicación.

Más allá del fichero de configuración, la flexibilidad de FedML en este trabajo se
sustenta en evidencia experimental directa: la implementación de la comparativa
define un entrenador y un agregador propios que extienden las abstracciones
\mintinline{python}{ClientTrainer} y \mintinline{python}{ServerAggregator} del
framework. El entrenador \texttt{FedMLCifar10Trainer} encapsula el bucle de
entrenamiento local sobre la ResNet-18 adaptada —descenso de gradiente estocástico,
entropía cruzada y métricas propias—, mientras que el agregador
\texttt{FedMLCifar10Aggregator} gobierna la evaluación del modelo global; ambos se
conectan al flujo mediante \texttt{FedMLRunner}. La documentación respalda esta vía,
al indicar que el usuario puede definir modelo, datos, optimizadores, funciones de
pérdida y lógica de entrenamiento arbitrarios, así como dataloaders propios que
respeten la estructura esperada por el simulador~\parencite{fedmlTrainer,fedmlDataloader}.
Esta capacidad de sustituir componentes centrales mediante herencia sostiene una
valoración alta en personalización del entrenamiento local, de la agregación y de la
integración.

\subsection{Sustitución de la estrategia federada}\label{subsec:flex-fedml-strategy}

La prueba de flexibilidad consistió en ejecutar cuatro estrategias modificando
únicamente la configuración, manteniendo constantes el conjunto de datos, el
modelo y el resto de hiperparámetros. FedAvg y FedProx se seleccionan asignando
\texttt{federated\_optimizer} a \texttt{"FedAvg"} o \texttt{"FedProx"}, mientras que
FedAdam y FedAdagrad se obtienen mediante \texttt{federated\_optimizer:\ "FedOpt"}
combinado con un optimizador de servidor \texttt{Adam} o \texttt{Adagrad}, tal como
ilustra el Código~\ref{cod:fedml-config}. Durante las primeras pruebas, las
variantes FedOpt fallaban porque FedML esperaba el atributo \texttt{ci} en los
argumentos de configuración; tras añadir \texttt{ci:\ 0} a la configuración común,
ambas pudieron completar la ejecución. La selección del algoritmo desde la
configuración está respaldada por la documentación, que controla el optimizador
federado a través del mismo campo~\parencite{fedmlSimApi}.

\begin{yamlcode}[title={Fragmento de fedml\_config.yaml (FedAdam vía FedOpt)}, label={cod:fedml-config}]
train_args:
  federated_optimizer: "FedOpt"
  server_optimizer: "Adam"
  ci: 0
  client_num_in_total: 4
  client_num_per_round: 2
  comm_round: 2
  epochs: 1
  batch_size: 32
  learning_rate: 0.02
\end{yamlcode}

\subsection{Mini-experimento de flexibilidad}\label{subsec:flex-fedml-mini}

La Tabla~\ref{tab:flex-fedml-resultados} resume las cuatro ejecuciones, todas
finalizadas con estado \texttt{success}. FedAvg y FedProx concluyeron con métricas
coherentes para una prueba de dos rondas, con precisiones del 43,01\,\% y el
42,60\,\% respectivamente, lo que confirma el soporte experimental de una estrategia
alternativa a FedAvg con comportamiento correcto. FedAdam y FedAdagrad también
completaron la ejecución, lo que demuestra que fue posible activar FedOpt con
optimizadores de servidor Adam y Adagrad; sin embargo, sus métricas finales fueron
claramente deficientes, con una precisión próxima al 10\,\% —equivalente al azar en
un problema de diez clases— y pérdidas muy elevadas.

\begin{table}[htbp]
  \centering
  \caption{Resultados del mini-experimento de flexibilidad de FedML. Las cifras de
  aprendizaje son auxiliares (dos rondas, una época, una semilla) y constituyen una
  comprobación funcional, no una comparación de rendimiento entre estrategias.}
  \label{tab:flex-fedml-resultados}
  \begin{tabular}{llrrrrr}
    \toprule
    \textbf{Estrategia} & \textbf{Estado} & \textbf{Dur. (s)} & \textbf{Acc.} &
    \textbf{Loss} & \textbf{GPU (MiB)} & \textbf{Util. (\%)} \\
    \midrule
    FedAvg     & success & 118 & 0,4301 & 1,5248   & 416 & 75,29 \\
    FedProx    & success & 123 & 0,4260 & 1,5137   & 484 & 79,01 \\
    FedAdam    & success & 117 & 0,0998 & 809,7697 & 416 & 76,09 \\
    FedAdagrad & success & 116 & 0,0998 & 191,2948 & 416 & 76,85 \\
    \bottomrule
  \end{tabular}
\end{table}

La lectura correcta de estos datos es que FedML permite cambiar la estrategia
federada mediante configuración, pero que la mera disponibilidad de una estrategia
no garantiza su estabilidad ni un buen rendimiento. Las variantes FedOpt
ejecutaron de principio a fin, de modo que la flexibilidad queda demostrada, pero su
falta de aprendizaje en esta prueba indica que requieren un ajuste específico de
hiperparámetros —como la tasa de aprendizaje del servidor o un mayor número de
rondas— para resultar competitivas. Este patrón es coherente con lo observado en los
otros frameworks: en Flower las estrategias adaptativas divergían con los valores por
defecto y exigieron ajustar \texttt{eta} y \texttt{tau}, y en NVIDIA FLARE FedAdam no
llegó a completarse en GPU. La fragilidad de las variantes FedOpt con poca
configuración es, por tanto, un comportamiento transversal y no una carencia
particular de FedML.

\subsection{Capacidades avanzadas según la documentación}\label{subsec:flex-fedml-doc}

La documentación de FedML respalda una superficie de capacidades más amplia que la
probada experimentalmente, organizada en varias capas. En el plano de uso, ofrece una
interfaz de línea de comandos —integrada en la plataforma TensorOpera— con órdenes para
autenticación y gestión de dispositivos (\texttt{fedml login}, \texttt{fedml device
bind}), lanzamiento y gestión de trabajos y clústeres (\texttt{fedml launch},
\texttt{fedml cluster}, \texttt{fedml run}), y gestión y despliegue de modelos como
servicios de inferencia (\texttt{fedml model})~\parencite{fedmlExamples}. En el plano de
programación, expone una API de Python con una llamada de una línea
(\texttt{fedml.run\_simulation()}) y una API más explícita de varias líneas que
inicializa dispositivo, datos y modelo y ejecuta la federación mediante
\texttt{SimulatorMPI} o las clases \texttt{Client}/\texttt{Server} envueltas por
\texttt{FedMLRunner}~\parencite{fedmlSimApi}.

En cuanto al catálogo algorítmico, el simulador Parrot enumera implementaciones de
referencia como FedOpt, FedNova, FedGKT, aprendizaje federado descentralizado, vertical
y jerárquico, FedNAS y \textit{Split Learning}, seleccionables cambiando el campo
\texttt{federated\_optimizer}~\parencite{fedmlSimOverview,fedmlSimApi}, y el módulo de
seguridad declara compatibilidad con optimizadores como FedAvg, FedSGD, FedOpt, FedProx,
FedGKT y FedNova~\parencite{fedmlSecurity}. A este catálogo se añaden la analítica
federada, que calcula estadísticas distribuidas como promedios o percentiles mediante un
ejecutor específico sin compartir datos crudos, y un proyecto independiente orientado al
ajuste federado de grandes modelos de lenguaje, que reutiliza las mismas API y la
plataforma de operaciones~\parencite{fedmlExamples}. En cuanto al aprendizaje asíncrono,
FedML Octopus documenta un escenario jerárquico heterogéneo en el que cada silo puede
ejecutar entrenamiento distribuido con varias GPU y orquestarse mediante optimización
federada asíncrona o síncrona~\parencite{fedmlCrossSilo}; se trata, no obstante, de un
soporte menos específico que el de NVIDIA FLARE, ya que no se documenta una receta directa
equivalente a FedBuff, aprendizaje cíclico o \textit{swarm}.

En privacidad y seguridad, el perfil de FedML difiere del de NVIDIA FLARE. Su
fortaleza documental está en la simulación de ataques y defensas mediante
\texttt{FedMLAttacker} y \texttt{FedMLDefender}, con ataques de tipo bizantino,
\textit{label flipping} y reemplazo de modelo, y defensas como recorte de norma, Krum,
mediana geométrica o FoolsGold, activables desde el fichero de configuración con
\texttt{enable\_attack} y \texttt{enable\_defense}~\parencite{fedmlSecurity}, además de un
ejemplo de agregación segura basada en computación multiparte~\parencite{fedmlExamples}.
En cambio, en las páginas revisadas no figura una integración de privacidad diferencial o
cifrado homomórfico comparable a la de NVIDIA FLARE, por lo que su valoración en este
criterio se sostiene en el \textit{benchmarking} de robustez más que en la privacidad
criptográfica. Por último, para integración y trazabilidad, FedML documenta el registro
de métricas, artefactos y modelos mediante \texttt{fedml.log()}, la captura automática de
métricas de hardware y GPU, los backends de comunicación MPI, gRPC, PyTorch RPC y MQTT+S3,
y la integración con la plataforma de operaciones de TensorOpera para gestionar agentes,
experimentos y despliegues~\parencite{fedmlTracking,fedmlExamples}.

\subsection{Limitaciones observadas}\label{subsec:flex-fedml-lim}

La limitación más visible es que las variantes FedOpt finalizaron sin aprender bajo
la configuración empleada, lo que confirma que su soporte es real pero condicionado a
un ajuste de hiperparámetros que queda fuera del alcance de esta prueba de
flexibilidad. A ello se añade el detalle de configuración del atributo \texttt{ci},
que tuvo que fijarse manualmente para que FedOpt funcionara, y el hecho de que la
prueba reducida de FedML emplea una configuración menor que la de los otros dos
frameworks, por lo que sus métricas auxiliares no son directamente comparables. En
cuanto a la ejecución, el modo concurrente de FedML mediante el backend MPI fue el
único que llegó a completarse entre los tres frameworks, según se documenta en el
capítulo de resultados de rendimiento, si bien se mostró sensible al consumo de
memoria en las configuraciones más grandes.

\subsection{Valoración de la flexibilidad de FedML}\label{subsec:flex-fedml-score}

La Tabla~\ref{tab:flex-fedml-score} recoge la valoración de FedML sobre los criterios
F1--F10, con el convenio ya aplicado a los otros frameworks. FedML obtiene la máxima
puntuación en personalización del modelo y los datos (F1), del entrenamiento local
(F2) y de la agregación (F3), respaldadas por la implementación de entrenador y
agregador propios; en los algoritmos federados disponibles (F4), por haber ejecutado
cuatro estrategias mediante configuración; en los modos de ejecución (F6), por ser el
único framework cuyo modo concurrente se completó; y en integración (F7) y
reproducibilidad (F10). El aprendizaje asíncrono (F5), la seguridad y privacidad (F8)
y la modificación de la arquitectura federada (F9) reciben puntuación parcial por
estar respaldados únicamente por la documentación.

\begin{table}[htbp]
  \centering
  \caption{Valoración de la flexibilidad de FedML según los criterios F1--F10. El
  tipo de evidencia distingue las capacidades ejecutadas en los experimentos de las
  descritas en la documentación oficial.}
  \label{tab:flex-fedml-score}
  \begin{tabular}{clcl}
    \toprule
    \textbf{Criterio} & \textbf{Descripción} & \textbf{Punt.} & \textbf{Evidencia} \\
    \midrule
    F1  & Personalización del modelo y el dataset      & 2/2 & Exp.\ + documental \\
    F2  & Personalización del entrenamiento local      & 2/2 & Exp.\ + documental \\
    F3  & Estrategias de agregación personalizadas     & 2/2 & Experimental \\
    F4  & Algoritmos federados disponibles             & 2/2 & Exp.\ + documental \\
    F5  & Aprendizaje federado asíncrono               & 1/2 & Documental \\
    F6  & Modos de ejecución y despliegue              & 2/2 & Exp.\ + documental \\
    F7  & Integración con herramientas externas        & 2/2 & Exp.\ + documental \\
    F8  & Seguridad, privacidad y extensiones          & 1/2 & Documental \\
    F9  & Modificación de la arquitectura federada     & 1/2 & Documental \\
    F10 & Reproducibilidad y configuración             & 2/2 & Exp.\ + documental \\
    \midrule
    \multicolumn{2}{r}{\textbf{Total}} & \multicolumn{2}{l}{\textbf{17/20}} \\
    \bottomrule
  \end{tabular}
\end{table}

En conjunto, FedML alcanza una valoración de flexibilidad de 17/20. Su mayor fortaleza
experimental reside en la facilidad para sustituir componentes centrales del sistema
—entrenador y agregador— mediante herencia, y en haber ejecutado el abanico más amplio
de estrategias del estudio a través de la configuración. Su perfil documental es
asimismo notable, con un catálogo extenso de algoritmos y escenarios y un soporte de
seguridad orientado al benchmarking de ataques y defensas. Las cautelas son dos: las
variantes FedOpt completaron la ejecución pero no aprendieron sin un ajuste
adicional, y algunas capacidades avanzadas —asincronía, privacidad criptográfica y
esquemas alternativos— solo se verificaron documentalmente.

% =====================================================================
%  SECCIÓN AÑADIDA: Comparativa de flexibilidad entre frameworks
% =====================================================================

\section{Comparativa de flexibilidad}\label{sec:flex-comparativa}

Una vez evaluados los tres frameworks por separado, esta sección los contrasta sobre
el mismo conjunto de criterios. La Tabla~\ref{tab:flex-comparativa} reúne las
puntuaciones F1--F10 de Flower, NVIDIA FLARE y FedML. Los tres obtienen la misma
valoración, 17/20, lo que confirma que las tres herramientas son, en lo esencial,
comparablemente flexibles para el tipo de adaptaciones que exige este trabajo y que la
diferenciación entre ellas debe buscarse, más que en la cifra global, en el perfil
cualitativo de cada una.

\begin{table}[htbp]
  \centering
  \caption{Comparativa de la valoración de flexibilidad de los tres frameworks según
  los criterios F1--F10 definidos en la metodología.}
  \label{tab:flex-comparativa}
  \begin{tabular}{clccc}
    \toprule
    \textbf{Criterio} & \textbf{Descripción} & \textbf{Flower} & \textbf{NVIDIA FLARE} & \textbf{FedML} \\
    \midrule
    F1  & Modelo y dataset                  & 2/2 & 2/2 & 2/2 \\
    F2  & Entrenamiento local               & 2/2 & 2/2 & 2/2 \\
    F3  & Agregación personalizada          & 2/2 & 2/2 & 2/2 \\
    F4  & Algoritmos disponibles            & 2/2 & 2/2 & 2/2 \\
    F5  & Aprendizaje asíncrono             & 1/2 & 1/2 & 1/2 \\
    F6  & Modos de ejecución y despliegue   & 2/2 & 2/2 & 2/2 \\
    F7  & Integración externa               & 2/2 & 2/2 & 2/2 \\
    F8  & Seguridad y privacidad            & 1/2 & 1/2 & 1/2 \\
    F9  & Arquitectura federada             & 1/2 & 1/2 & 1/2 \\
    F10 & Reproducibilidad y configuración  & 2/2 & 2/2 & 2/2 \\
    \midrule
    \multicolumn{2}{r}{\textbf{Total}} & \textbf{17/20} & \textbf{17/20} & \textbf{17/20} \\
    \bottomrule
  \end{tabular}
\end{table}

La coincidencia en los criterios experimentales (F1--F4, F6, F7 y F10) refleja un
hecho de fondo: bajo un mismo escenario de modelo, datos y algoritmo base, los tres
frameworks permitieron integrar la ResNet-18 sobre CIFAR-10, editar el entrenamiento
local, sustituir la estrategia de agregación, registrar métricas y reproducir los
experimentos mediante ficheros y scripts. La igualdad también alcanza a los tres
criterios documentales: ninguno validó experimentalmente el aprendizaje asíncrono (F5),
la seguridad avanzada (F8) ni esquemas alternativos a la arquitectura clásica (F9), pero
los tres los documentan, por lo que reciben idéntica puntuación parcial. En F5, en
concreto, los tres declaran soporte —Flower mediante la \textit{Driver API} y
\textit{Flower Next}, NVIDIA FLARE mediante FedBuff en el borde y FedML en su escenario
jerárquico cross-silo—, si bien en todos los casos se trata de capacidades
experimentales, limitadas a entornos concretos o de bajo nivel.

Más allá de las cifras, el estudio revela un comportamiento transversal en las
estrategias adaptativas de la familia FedOpt. En los tres frameworks fue posible
configurar FedAdam y FedAdagrad, pero su ejecución resultó frágil con poca
configuración: en Flower divergían con los valores por defecto y necesitaron ajustar
la tasa del servidor y el factor de adaptabilidad; en NVIDIA FLARE, FedAdam no llegó a
completarse en GPU por un error de CUDA; y en FedML, FedAdam y FedAdagrad finalizaron
pero sin aprender. La conclusión es clara y común: la disponibilidad de una estrategia
no equivale a su estabilidad, que depende de un ajuste deliberado de hiperparámetros,
especialmente con pocas rondas.

Por último, cada framework muestra un perfil de flexibilidad propio que la puntuación,
por su granularidad, no captura del todo. Flower es el más directo para sustituir la
estrategia de agregación —basta con cambiar una clase en el servidor— y su documentación
revela una superficie amplia que abarca \textit{Mods}, privacidad diferencial, FedBN y,
sobre todo, un modelo de despliegue muy desarrollado hacia la nube (Google Cloud, Azure,
OpenShift e interconexión de clústeres); su contrapartida es que el soporte asíncrono
sigue siendo experimental y que el backend Ray arrastra una huella de memoria elevada.
NVIDIA FLARE ofrece la mayor profundidad en seguridad y privacidad —filtros de privacidad
diferencial, cifrado homomórfico y computación confidencial sobre entornos de ejecución
confiables— junto con un catálogo de recetas y entornos (\texttt{SimEnv}, \texttt{PocEnv},
\texttt{ProdEnv}) orientado a producción, a costa de una complejidad de configuración
mayor, pues alterar la agregación exige intervenir en componentes del servidor. FedML
destaca por la sustitución de entrenador y agregador mediante herencia, por haber
ejecutado el mayor número de estrategias, por ser el único cuyo modo concurrente se
completó, y por una integración estrecha con su plataforma de operaciones y la analítica
federada, con un soporte de seguridad orientado al \textit{benchmarking} de ataques y
defensas. La elección entre ellos, en términos de flexibilidad, depende por tanto de la
prioridad del proyecto: la simplicidad y el despliegue en la nube de Flower, la
profundidad en seguridad y la orientación a producción de NVIDIA FLARE, o la
extensibilidad por componentes y la experimentación algorítmica de FedML.
